\documentclass[11pt]{article}

\usepackage[margin=1.1in]{geometry}

\usepackage{amssymb,mathtools,natbib}

% Anchored formal environments for causalsmith papers.
% First mandatory argument = obj_id (machine anchor joining the paper to the
% Lean crosswalk; invisible in the rendered PDF). Optional argument = name.
% Each environment also defines \label{obj:<obj_id>} so prose can use
% \cref{obj:T-1}; cleveref derives the printed kind and number from the target.
\usepackage{amsmath}
\usepackage{mathrsfs}
\usepackage{amsthm}
\usepackage{booktabs}
% Long displays may break across pages; let TeX stretch a little harder before an
% overfull \hbox so wide formulas/tables are less likely to run off the page.
\allowdisplaybreaks
\setlength{\emergencystretch}{3em}
\newtheorem{theorem}{Theorem}
\newtheorem{lemma}{Lemma}
\newtheorem{assumption}{Assumption}
\newtheorem{definition}{Definition}
\newtheorem{proposition}{Proposition}
\newtheorem{remark}{Remark}
\newtheorem{citedresult}{Cited result}
% The amsthm counter is `algorithmbox` (not `algorithm`) so a later `\usepackage{algorithm}` cannot clash.
\newcounter{algorithmbox}
\renewcommand{\thealgorithmbox}{\arabic{algorithmbox}}
\NewDocumentEnvironment{theoremv}{m o s}{\begin{theorem}\label{obj:#1}{\normalfont[\texttt{\detokenize{#1}}]\IfValueT{#2}{\ (#2)}\IfBooleanT{#3}{\verificationfootnotemark}\ }}{\end{theorem}}
\NewDocumentEnvironment{lemmav}{m o s}{\begin{lemma}\label{obj:#1}{\normalfont[\texttt{\detokenize{#1}}]\IfValueT{#2}{\ (#2)}\IfBooleanT{#3}{\verificationfootnotemark}\ }}{\end{lemma}}
\NewDocumentEnvironment{assumptionv}{m o s}{\begin{assumption}\label{obj:#1}{\normalfont[\texttt{\detokenize{#1}}]\IfValueT{#2}{\ (#2)}\IfBooleanT{#3}{\verificationfootnotemark}\ }}{\end{assumption}}
\NewDocumentEnvironment{definitionv}{m o s}{\begin{definition}\label{obj:#1}{\normalfont[\texttt{\detokenize{#1}}]\IfValueT{#2}{\ (#2)}\IfBooleanT{#3}{\verificationfootnotemark}\ }}{\end{definition}}
% A `propositionv` is a proved result tiered as a Proposition (theorem-grade, machine-verified).
\NewDocumentEnvironment{propositionv}{m o s}{\begin{proposition}\label{obj:#1}{\normalfont[\texttt{\detokenize{#1}}]\IfValueT{#2}{\ (#2)}\IfBooleanT{#3}{\verificationfootnotemark}\ }}{\end{proposition}}
% A `remarkv` holds NON-load-bearing interpretive / open-question content (e.g. a causalsmith `oeq:`
% open question). It is neither proved nor assumed; no theorem may depend on it.
\NewDocumentEnvironment{remarkv}{m o s}{\begin{remark}\label{obj:#1}{\normalfont[\texttt{\detokenize{#1}}]\IfValueT{#2}{\ (#2)}\IfBooleanT{#3}{\verificationfootnotemark}\ }}{\end{remark}}
% An `algorithmv` is a DEFINITION-kind object presented as a boxed, numbered-step procedure (an
% estimator / selection rule built in stages). Same anchor contract as `definitionv`; its body is
% ordinary LaTeX whose steps are `\begin{enumerate}` items, so tex2html needs no extra machinery.
% Presented in the CONVENTIONAL algorithm style readers expect from `algorithm`+`algorithmic`:
% a rule above, an "Algorithm N (Title)" caption line, a rule under the caption, the numbered
% steps, and a closing rule. It is NOT a float — the anchor contract requires the object to stay
% where the outline places it, and floats drift. The obj-id tag is suppressed here (it is
% bookkeeping, and a caption line carrying it reads as debug output in a finished paper).
\NewDocumentEnvironment{algorithmv}{m o s}%
  {\par\addvspace{\medskipamount}\noindent\rule{\linewidth}{0.8pt}\par\nobreak\vspace{2pt}%
   \refstepcounter{algorithmbox}\label{obj:#1}%
   \noindent\textbf{Algorithm \thealgorithmbox}\IfValueT{#2}{ \textnormal{(#2)}}%
   \IfBooleanT{#3}{\verificationfootnotemark}\par\nobreak\vspace{2pt}%
   \noindent\rule{\linewidth}{0.4pt}\par\nobreak\vspace{2pt}\normalfont}%
  {\par\nobreak\vspace{2pt}\noindent\rule{\linewidth}{0.8pt}\par\addvspace{\medskipamount}}
% A `citedv` is an IMPORTED external result the paper relies on but does NOT prove
% (a causalsmith `gate_class:"cited"` node). It is machine-checked against the cited
% source at F2.5, never proved here; the fixed provenance note makes that explicit
% so the environment can never be read as a contribution of this paper. The \citep
% attribution is supplied by the surrounding prose (P2), keyed to references.bib.
\NewDocumentEnvironment{citedv}{m o s}{\begin{citedresult}\label{obj:#1}{\normalfont[\texttt{\detokenize{#1}}]\IfValueT{#2}{\ (#2)}\IfBooleanT{#3}{\verificationfootnotemark}\ \textit{(imported result; machine-checked against the cited source, not proved here.)}\ }}{\end{citedresult}}
% \leanref{obj_id}{display text}: an inline first-mention link to a from-note object's
% verified Lean code. In the PDF it renders the display text plainly (the Lean link is a
% web-only affordance); tex2html turns it into a drawer-opening span.
\newcommand{\leanref}[2]{#2}
% For a theorem/lemma whose Lean declaration accepts a source-matched published proposition as an
% input, the starred anchored environment puts the mark in its heading; the generated text remains
% outside the frozen mathematical environment. Keep the combined macro for older paper bundles.
\newcommand{\verificationfootnotemark}{\footnotemark}
\newcommand{\verificationfootnotetext}[1]{\footnotetext{#1}}
\newcommand{\verificationfootnote}[1]{\verificationfootnotemark\verificationfootnotetext{#1}}
% xcolor before hyperref (hyperref would otherwise pull in plain `color` for colorlinks, and a
% later xcolor load clashes). The page-1 identity stamp at the end of this file needs a grey.
\usepackage{xcolor}
\usepackage{graphicx}
\usepackage{eso-pic}
% hyperref follows natbib and the environment macros; cleveref must follow hyperref.
\usepackage[colorlinks=true,linkcolor=blue,citecolor=blue,urlcolor=blue]{hyperref}
\usepackage[capitalise,nameinlink,noabbrev]{cleveref}
% Pin the names explicitly: labels are placed inside the underlying amsthm environments, and these
% declarations make the PDF contract independent of cleveref's theorem-name inference. House style:
% reference names always print with a capital first letter ("Theorem 1", "Definition 2"), in any
% sentence position — hence `capitalise` above and capitalized \crefname forms below.
\crefname{theorem}{Theorem}{Theorems}
\Crefname{theorem}{Theorem}{Theorems}
\crefname{lemma}{Lemma}{Lemmas}
\Crefname{lemma}{Lemma}{Lemmas}
\crefname{assumption}{Assumption}{Assumptions}
\Crefname{assumption}{Assumption}{Assumptions}
\crefname{definition}{Definition}{Definitions}
\Crefname{definition}{Definition}{Definitions}
\crefname{proposition}{Proposition}{Propositions}
\Crefname{proposition}{Proposition}{Propositions}
\crefname{remark}{Remark}{Remarks}
\Crefname{remark}{Remark}{Remarks}
\crefname{citedresult}{Cited result}{Cited results}
\Crefname{citedresult}{Cited result}{Cited results}
\crefname{algorithmbox}{Algorithm}{Algorithms}
\Crefname{algorithmbox}{Algorithm}{Algorithms}
% ---------------------------------------------------------------------------
% Page-1 identity stamp (arXiv style) and the pinned title-block date.
%
% Split of responsibility: THIS file owns the FORMAT (placement, font, colour, punctuation),
% the generated `paper_stamp.tex` owns only the VALUES (number+version, area, revised date,
% DOI, link target). P4 refreshes this template on every emit, so a layout fix reaches every
% bundle from a recompile; and because `paper_stamp.tex` is not one of the P2 assembly
% digest's authored inputs, a DOI that arrives after publication needs only that one small
% file rewritten plus a recompile — never a redraft, never a reassembly.
%
% Defaults first, so a bundle WITHOUT a stamp file compiles exactly as it always did:
% `\cspaperdate` falls back to `\today` and no stamp is drawn.
\providecommand*{\cspaperdate}{\today}  % the \date{} target
\providecommand*{\csstampid}{}          % e.g. CSWP-2026-008v3 (empty ⇒ no stamp at all)
\providecommand*{\csstamparea}{}        % e.g. Stat
\providecommand*{\csstampdate}{}        % e.g. 27 Aug 2026
\providecommand*{\csstampdoi}{}         % e.g. 10.5281/zenodo.123 (empty ⇒ DOI part omitted)
\providecommand*{\csstamplink}{}        % href target (DOI url, else the paper's page; may be empty)
\IfFileExists{paper_stamp.tex}{% GENERATED by CausalSmith P4 — paper identity stamp values. Do not hand-edit.
%
% Field order on the stamp (owned by paper_macros.tex):
%   <number>v<version> · doi:<doi> · [<Area>] <date>   — the DOI is never the last token, so a
%   text extractor cannot run it into whatever follows. Without a DOI that part is omitted.
% Format lives in paper_macros.tex; only the values are here, and this file is NOT one of
% the P2 assembly digest's authored inputs. A DOI arriving after publication therefore needs
% this file rewritten plus a `latexmk` recompile — never a redraft or a reassembly.
\renewcommand*{\csstampid}{CSWP-2026-039v3}
\renewcommand*{\csstamparea}{Experimentation}
\renewcommand*{\csstampdate}{3 Oct 2026}
\renewcommand*{\csstampdoi}{}
\renewcommand*{\csstamplink}{https://causalsmith.org/papers/exp_thinnedgraph_additive_risk_frontier_v1}
\renewcommand*{\cspaperdate}{3 October 2026}
}{}
\makeatletter
% Field order is deliberate: the DOI is NEVER the last token on the line. A trailing identifier
% runs into whatever a text extractor emits next, which makes it unsafe to match mechanically
% (the Zenodo stamp matcher reads exactly this line); the date is a safe terminator.
%   <number>v<version> · doi:<doi> · [<Area>] <date>      — with a DOI
%   <number>v<version> · [<Area>] <date>                  — without
\newcommand{\cs@stampsep}{\,\textperiodcentered\,}
\newcommand{\cs@stamptext}{%
  \csstampid
  \ifx\csstampdoi\@empty\else\cs@stampsep doi:\csstampdoi\fi
  \ifx\csstamparea\@empty\else\cs@stampsep[\csstamparea]\fi
  \ifx\csstampdate\@empty\else\quad\csstampdate\fi
}
\ifx\csstampid\@empty\else
  \definecolor{csstampgrey}{gray}{0.45}
  % Background shipout material: it occupies no space, so the text block and the \thanks
  % footnote are byte-identical to an unstamped compile. The starred form applies to the
  % NEXT page shipped only — i.e. page 1. Placed in the left margin (the text block starts
  % at 1.1in), reading bottom-to-top, in a Times-like face at 8pt.
  \AddToShipoutPictureBG*{%
    \AtPageLowerLeft{%
      \put(\LenToUnit{0.42in},\LenToUnit{1.1in}){%
        \rotatebox{90}{%
          \fontfamily{ptm}\fontsize{8}{10}\selectfont
          \ifx\csstamplink\@empty
            \textcolor{csstampgrey}{\cs@stamptext}%
          \else
            \expandafter\href\expandafter{\csstamplink}{\textcolor{csstampgrey}{\cs@stamptext}}%
          \fi
        }%
      }%
    }%
  }
\fi
\makeatother


\title{Minimax precision for total treatment effects with partially observed interference networks}

\author{CausalSmith\thanks{Machine-generated by the CausalSmith research pipeline.}}

\date{\cspaperdate}

\begin{document}

\maketitle

\begin{abstract}
We characterize minimax squared-error precision for the population all-treated-versus-all-control contrast under heterogeneous additive interference and sampled network information. The experiment assigns treatment independently with probability one half, retains each true directed off-diagonal arrow independently with known probability \(q\), the edge-retention probability, and observes the retained graph, every assignment, and every noiseless realized outcome. Fixed schedules have population size \(n\), common in- and out-degree bound \(d\), and absolute coefficient mass at most one for each outcome, including its baseline and own-treatment coefficients. Uniformly over integers \(n\ge4\) and \(1\le d\le n-1\), and retention probabilities \(0<q\le1\), minimax squared risk is of order
\[
\min\left\{1,\frac{d^2}{n[1-(1-q)^d]}\right\}
\asymp
\min\left\{1,\frac{d^2}{n}+\frac{d}{nq}\right\},
\]
with universal comparison constants; at zero retention it is of constant order. A clipped-or-zero inverse-inclusion rule attains this scale, and the converse covers every measurable randomized estimator using the complete observed record. For degree sequences \(d_n\) and retention sequences \(q_n\), uniform consistency is possible exactly when \(d_n^2/n\to0\) and \(d_n/(nq_n)\to0\), interpreting the second ratio as infinite at zero retention. Retention of order \(1/d\) or greater attains the supplied-graph precision order. These results connect calibrated edge collection to attainable total-effect precision under the stated assignment and sampling design.
\end{abstract}

\section{Introduction}

This paper characterizes the optimal squared-error precision for estimating a population total treatment effect from one randomized experiment with a sampled interference graph. The target compares assigning treatment to everyone with assigning treatment to no one, averaging the resulting outcome differences over a fixed population. We study heterogeneous additive responses under bounded absolute coefficient mass and bounded directed in- and out-degrees. Treatment is assigned independently with probability one half, and each true off-diagonal interference arrow is independently retained with a known common probability. Under this experiment, we establish matching minimax risk bounds, construct an observable attaining rule, and characterize exactly when uniformly consistent estimation is possible.

The precision question arises when network information is collected alongside an experiment. Assignments and realized outcomes are available for the full population, while recorded arrows reveal a sampled portion of the interference structure. The retained graph contributes information about which treatments enter each outcome; the outcomes themselves can also inform those associations. An optimal-precision analysis therefore evaluates the joint information in graph labels, assignments, and outcomes. Our minimax criterion allows every measurable randomized estimator to use that complete record, as specified in \cref{obj:synth_3,obj:def:risk}.

The schedule restrictions make the causal target and information comparison precise. Write \(n\), the population size, and \(d\), the common off-diagonal in-degree and out-degree bound, with \(n\ge4\) and \(1\le d\le n-1\). Each outcome is the sum of a fixed baseline, an own-treatment effect, and source-specific additive spillovers. For each recipient, the sum of the absolute baseline, own-treatment, and spillover coefficients is at most one. Coefficients may be heterogeneous and have either sign. The graph and coefficients remain fixed during treatment randomization and edge collection. This is the schedule class in \cref{obj:def:model}. Its correspondence with the published interaction-order-one SNIPE model adjoins every own-treatment coordinate and translates the off-diagonal degree bound into an augmented degree bound of \(d+1\), as recorded in \cref{obj:lem:published-model}.

Let \(q\), the known edge-retention probability, range over \([0,1]\). The principal result in \cref{obj:thm:observed-record-precision-frontier} establishes, uniformly over the stated parameter domain, that the minimax squared risk \(R_{n,d}(q)\) satisfies
\[
R_{n,d}(q)\asymp
\begin{cases}
\displaystyle
\min\left\{1,\frac{d^2}{n[1-(1-q)^d]}\right\},&q>0,\\[6pt]
1,&q=0,
\end{cases}
\]
with universal comparison constants. For positive retention, the equivalent expression
\[
R_{n,d}(q)\asymp
\min\left\{1,\frac{d^2}{n}+\frac{d}{nq}\right\}
\]
separates a degree contribution from an edge-collection contribution. The cap at one describes saturation of worst-case squared risk under the bounded target range.

The upper comparison is constructive. The inverse-inclusion score in \cref{obj:def:audit-score} combines each realized outcome with its own-treatment sign and the treatment signs of retained sources, weighting recorded arrows by \(q^{-1}\). For every fixed admissible schedule at positive retention, this score is unbiased over assignment and audit randomness. Its variance decomposes into assignment variation with the graph supplied and an additional audit term, as shown in \cref{obj:lem:score-audit}. The attaining rule clips the score to the target interval when its risk envelope is below one and uses zero otherwise. \Cref{obj:thm:observable-upper,obj:thm:observed-record-precision-frontier} establish its uniform guarantee and minimax order.

The lower comparison evaluates the same complete observation experiment. It uses opposite-sign priors over fixed schedules, with baseline and source-allocation draws occurring before the experiment. The resulting likelihood preserves retained-edge labels, all treatment coordinates, repeated recipient outcomes, and the treatment-count constraint on hidden source allocations. Translation bounds, all-order Walsh energy control, and hidden-allocation contraction yield the testing comparison in \cref{obj:thm:block-testing-scale}. Together with the two-prior risk reduction in \cref{obj:lem:full-record-two-prior-risk} and the supplied-graph lower comparison in \cref{obj:lem:supplied-block-record-lower}, these arguments establish the converse against arbitrary measurable randomized uses of the record throughout the admissible degree and retention ranges.

The characterization gives two collection implications. First, retention of order \(1/d\) or greater attains the supplied-graph risk order \(\min\{1,d^2/n\}\); below this transition, the collection contribution determines the uncapped scale. Second, along admissible population sequences with degree bounds \(d_n\) and retention probabilities \(q_n\), uniform consistency is possible exactly when
\[
d_n=o(\sqrt n),
\qquad
q_n>0\ \text{eventually},
\qquad
\frac{nq_n}{d_n}\longrightarrow\infty.
\]
Thus consistency can coexist with vanishing retention, and retention of order \(1/d\) or greater suffices to match supplied-graph precision. These conclusions follow from \cref{obj:thm:observed-record-precision-frontier,obj:thm:frontier-answer}.

The closest estimator comparison is known-neighborhood SNIPE, which studies the same population contrast under Bernoulli assignment and bounded-order polynomial responses \citep[Sections 3--5; Theorems 1--2]{CortezRodriguezEichhornYu2023SNIPE}. Its published variance bound specializes to order \((d+1)^3/n\) under the additive, half-Bernoulli, unit-envelope conventions used here; its published minimax construction has its own treatment-probability conditions. The sharper supplied-graph degree benchmark of \citet[Theorem \texttt{thm:degree-frontier}; non-archival research note]{CausalSmithKnownGraph2026} is restated for the present conventions in \cref{obj:thm:supported-boundaries}. Design-based minimax analysis with a supplied network further studies joint optimization over assignments and estimators under a specified response envelope \citep[Corollary 5.4]{KandirosHarshawSavje2026DesignBasedMinimax}. Our new contribution fixes the assignment law and characterizes sampled-network precision, including a complete-record converse and the joint degree--retention conditions for consistency.

Known network-error probabilities also support correction of exposure-specific means \citep[Section 3; Theorems 4--7]{LiSussmanKolaczyk2021NetworkNoise}, while supplied neighborhood supersets support unbiased total-effect estimation under additive and polynomial responses \citep[Theorems 1, 4, and 8; Appendix B.1]{ShankarEtAl2024UNITE}. These approaches clarify the role of network information in causal estimation. The present result characterizes worst-case total-effect precision under independent retention of true arrows and heterogeneous additive coefficient restrictions. Under bounded outcomes and first- and second-order influence restrictions, the pseudoinverse estimator of \citet[Assumptions 1--2 and Theorems 1--2]{JiangDengWangWang2024ApproximateNetworks} has asymptotic variance \(O\{d_{\mathcal G}^2/[n\pi(1-\pi)]\}\), with the same estimator-specific order for constant outcomes on a regular surrogate graph. Our result instead optimizes over all estimators under the calibrated retention experiment.

Proofs of the results appear in the appendices.

The paper proceeds through related work, the population model and observation design, minimax precision and attainable estimation, and the edge-collection comparison with the supplied-graph benchmark. A limitations and future-work discussion follows. The appendices develop the score calculations, complete-record block experiments, analytic and testing bounds, and minimax arguments, followed by the verification note and proofs of the main results.


\section{Related work}

The closest comparisons concern the population contrast between assigning treatment to everyone and assigning treatment to no one. For this target, precision depends jointly on the response class, the assignment mechanism, and the network information supplied to the estimator. The present analysis combines additive responses with bounded coefficient mass and bounded in- and out-degrees, independent treatment assignment with probability one half, and independent edge retention with a known probability. Over the schedule class in \cref{obj:def:model}, the risk in \cref{obj:def:risk} allows every measurable randomized estimator to use the retained graph, the full assignment vector, and all realized outcomes. The characterization in \cref{obj:thm:observed-record-precision-frontier} therefore describes optimal squared-error precision for this complete observation experiment.

Known-neighborhood SNIPE provides the closest estimator comparison. It targets the same population contrast under independent Bernoulli assignment and polynomial neighborhood responses of bounded interaction order. Its published analysis establishes unbiasedness, a variance upper bound, and a minimax lower bound \citep[Sections 3--5; Theorems 1--2]{CortezRodriguezEichhornYu2023SNIPE}. The present additive model corresponds to interaction order one. Because the published neighborhoods include each recipient itself, the common off-diagonal degree bound \leanref{sym:d}{\ensuremath{d}} in \cref{obj:def:model} translates into augmented in- and out-degree bounds of \(d+1\); the coefficient-mass normalization is preserved by this embedding, as recorded in \cref{obj:lem:published-model}. For population size \leanref{sym:n}{\ensuremath{n}}, the published variance bound specializes to order \((d+1)^3/n\) at treatment probability one half. Its minimax lower bound uses Gaussian coefficient priors and, at interaction order one, a common treatment probability below \(0.16\). These statements furnish distinct comparisons for the bounded additive class and assignment mechanism studied here. Under full edge retention, \cref{obj:thm:observable-upper} supplies a quadratic degree upper bound. The sharper supplied-graph benchmark is stated in \cref{obj:thm:supported-boundaries}.
% [SNIPE source](https://arxiv.org/pdf/2208.05553)

A second comparison concerns optimization of the experiment itself. With the interference network supplied, design-based minimax analysis can optimize jointly over assignment mechanisms and estimators under an arbitrary neighborhood response model and its specified potential-outcome envelope. For random regular networks whose degree grows more slowly than the square root of population size, the global average treatment effect has minimax risk bounded below by degree divided by population size and above by squared degree divided by population size, with high probability over the network \citep[Corollary 5.4]{KandirosHarshawSavje2026DesignBasedMinimax}. That comparison concerns a fixed supplied network and design optimization under its own response envelope. Here, the assignment law is fixed, the response envelope bounds absolute additive coefficient mass, and the worst case ranges over graphs and coefficients. The supplied-graph risk in \cref{obj:def:supplied-risk} isolates the information comparison while retaining the paper's assignment mechanism and schedule class.
% [Design-based minimax source](https://arxiv.org/html/2608.04909v1)

Network measurement error also motivates correction of exposure classifications. Under a four-level exposure model and Bernoulli treatment assignment, known false-positive and false-negative edge probabilities support method-of-moments correction of exposure-specific mean outcomes. The resulting guarantees include asymptotic bias and variance control and, under additional conditions, asymptotic normality \citep[Section 3; Theorems 4--7]{LiSussmanKolaczyk2021NetworkNoise}. This work credits the role of known network-error probabilities in correcting causal estimators. The present observation model instead supplies a randomly retained subgraph of the true directed graph, and the response class permits heterogeneous additive effects from individual sources. Its precision guarantee concerns the all-treated-versus-all-control contrast and worst-case squared loss over all estimators.
% [Network-noise source](https://arxiv.org/pdf/2105.04518)

Supplied supergraphs offer another form of partial network information. UNITE estimates the global average treatment effect under Bernoulli assignment using a supplied neighborhood superset. Its additive and bounded-order polynomial analyses establish unbiasedness and consistency under the stated containment and response conditions, together with variance bounds whose constants depend on neighborhood sizes, assignment probabilities, and the response envelope \citep[Theorems 1, 4, and 8; Appendix B.1]{ShankarEtAl2024UNITE}. A containing supergraph supplies every relevant source as a candidate, whereas a retained subgraph supplies sampled true arrows together with their inclusion probability. Baseline information provides a further route: under heterogeneous additive responses and equal marginal treatment probabilities along interfering edges, an observed baseline supports an unbiased graph-free estimator of the total treatment effect \citep[Corollary 3]{YuAiroldiBorgsChayes2022}. These results clarify how neighborhood containment and baseline observations change the estimator's information set. In the present experiment, baselines are fixed schedule coefficients, and estimation uses the retained graph, assignments, and realized outcomes.
% [UNITE source](https://proceedings.mlr.press/v238/shankar24a/shankar24a.pdf), [Observed-baseline source](https://pmc.ncbi.nlm.nih.gov/articles/PMC9636977/)

Approximate-network analysis studies a related bias--variance trade-off for a specified surrogate-network estimator. Under bounded outcomes and restrictions on first- and second-order treatment influences, its asymptotic variance is \(O\{d_{\mathcal G}^2/[n\pi(1-\pi)]\}\), where \(d_{\mathcal G}\) is surrogate-network degree and \(\pi\) is treatment probability; constant outcomes on a regular surrogate graph give the same estimator-specific order. A separate omitted-influence condition controls relative bias \citep[Assumptions 1--3; Theorems 1--2; Lemma 2]{JiangDengWangWang2024ApproximateNetworks}. Those guarantees describe the performance of a chosen estimator as the surrogate network captures different amounts of influence. The present characterization optimizes squared risk over the full estimator class under the explicit edge-retention experiment. Its lower bound applies to unrestricted use of the observed record, while its upper bound supplies an attainable procedure under the same schedule restrictions. This common experiment makes the contribution interpretable as a precision characterization indexed by population size, degree, and edge retention.
% [Approximate-network source](https://arxiv.org/html/2408.04441v1)


\section{Population, causal target, and observation design}

Throughout, indices range over the fixed population unless another range is specified. Absolute-value signs denote scalar magnitude or set cardinality according to their argument. Generic positive constants may vary across statements; the notation \(A\asymp B\) denotes upper and lower bounds by positive constant multiples, with the stated uniformity.

Fix a \leanref{S-1}{finite-population setup} with population \leanref{sym:V}{\ensuremath{V=\{1,\ldots,n\}}}, population size \(n\ge4\), and degree index \(1\le d\le n-1\). The recipient index is \leanref{sym:i}{\ensuremath{i}}, and the source index is \leanref{sym:j}{\ensuremath{j}}. A simple directed off-diagonal interference graph \leanref{sym:G}{\ensuremath{G}} contains arrows \((j,i)\) from source \(j\) to recipient \(i\). A fixed graph-and-coefficient schedule \leanref{sym:theta}{\ensuremath{\theta=(G,a,t,b)}} specifies the baseline vector \leanref{sym:a}{\ensuremath{a}}, own-treatment vector \leanref{sym:t}{\ensuremath{t}}, and spillover array \leanref{sym:b}{\ensuremath{b}}. Thus \(b_{ij}\) records the effect of source \(j\)'s treatment on recipient \(i\).

The population-and-outcome conventions begin with the potential-outcome map \leanref{sym:Y}{\ensuremath{Y}}. Its argument \leanref{sym:z}{\ensuremath{z\in\{0,1\}^V}} is a binary assignment vector, and its additive form separates baseline, own-treatment, and source-specific spillover contributions.

\begin{definitionv}{synth_1}[Additive potential outcomes]
For a population \(V=\{1,\ldots,n\}\) and a fixed schedule \(\theta=(G,a,t,b)\), where \(G\subseteq\{(j,i)\in V^2:j\ne i\}\) is a simple directed off-diagonal graph, define the potential-outcome map \(Y:\{0,1\}^V\to\mathbb R^V\) by
\[
Y(z)=(Y_i(z))_{i\in V},
\qquad
Y_i(z)=a_i+t_i z_i+\sum_{j:(j,i)\in G}b_{ij}z_j,
\qquad z\in\{0,1\}^V.
\]
Here \(a_i\) is the fixed baseline coefficient, \(t_i\) is the fixed own-treatment coefficient, and \(b_{ij}\) is the fixed spillover coefficient on the arrow from source \(j\) to recipient \(i\). For the realized assignment \(Z\), the realized outcome vector is \(Y(Z)\).
\end{definitionv}

For the outcome map in \cref{obj:synth_1}, let \(\boldsymbol 1\) and \(\boldsymbol 0\) denote the assignments treating every unit and treating no unit, respectively. Define the population all-treated-versus-all-control contrast \leanref{sym:tau}{\ensuremath{\tau(\theta)}} by
\[
\tau(\theta)
=
\frac1n\sum_{i\in V}\bigl(Y_i(\boldsymbol 1)-Y_i(\boldsymbol 0)\bigr)
=
\frac1n\sum_{i\in V}
\left(t_i+\sum_{j:(j,i)\in G}b_{ij}\right).
\]
This contrast averages the combined own-treatment and spillover effects of moving the entire population between the two assignments. The schedule, and hence this causal target, remains fixed throughout the experiment.

We impose three restrictions on schedules: bounds on the number of sources affecting each recipient, the number of recipients affected by each source, and the total absolute coefficient mass for each outcome.

\begin{assumptionv}{ass:in-degree}[In-degree bound]
The interference graph \(G\) on the finite population \(V\) satisfies the in-degree bound \(d\):
\[
\max_{i\in V}\bigl|\{j:(j,i)\in G\}\bigr|\le d.
\]
\end{assumptionv}

The recipient in-degree condition bounds the number of other units whose treatments can enter a given outcome. Including the compulsory own-treatment coordinate gives the augmented in-degree bound \(d+1\), corresponding to the neighborhood-size restriction in the published additive model \citep{CortezRodriguezEichhornYu2023SNIPE}. The source-side restriction controls how widely a single treatment coordinate can affect outcomes.

\begin{assumptionv}{ass:out-degree}[Out-degree bound]
The interference graph \(G\) on the finite population \(V\) satisfies the out-degree bound \(d\):
\[
\max_{j\in V}\bigl|\{i:(j,i)\in G\}\bigr|\le d.
\]
\end{assumptionv}

On the source side, the out-degree condition bounds the number of other recipients reached by one treatment coordinate. Its augmented counterpart is the published out-degree bound \(d+1\) \citep{CortezRodriguezEichhornYu2023SNIPE}. Together, \cref{obj:ass:in-degree,obj:ass:out-degree} control both the breadth of each unit's exposure and the reach of each source's treatment. The coefficient-mass restriction above controls outcome magnitudes.

\begin{assumptionv}{ass:coefficient-mass}[Unit coefficient-mass bound]
The baseline coefficients \(a_i\), own-treatment coefficients \(t_i\), and spillover coefficients \(b_{ij}\) on the interference graph \(G\) satisfy
\[
\max_{i\in V}
\left(
|a_i|+|t_i|+\sum_{j:(j,i)\in G}|b_{ij}|
\right)\le 1,
\]
where \(V\) is the finite population.
\end{assumptionv}

The coefficient-mass condition is the published interaction-order-one SNIPE restriction with \(Y_{\max}\), the maximum absolute coefficient mass, bounded by one; it includes both the baseline and the diagonal singleton coefficient \citep{CortezRodriguezEichhornYu2023SNIPE}. It supplies a common outcome scale while permitting heterogeneous coefficients of either sign. In particular, \cref{obj:synth_1,obj:ass:coefficient-mass} imply \(|Y_i(z)|\le1\) for every unit and assignment, and \(|\tau(\theta)|\le1\). Collecting these restrictions defines the schedule class over which precision is evaluated.

\begin{definitionv}{def:model}[Schedule class $\mathcal M_{n,d}$]
\[
\mathcal M_{n,d}
=
\left\{\theta=(G,a,t,b):
\max_{i\in V}|\{j:(j,i)\in G\}|\le d,\quad
\max_{j\in V}|\{i:(j,i)\in G\}|\le d,\quad
\max_{i\in V}\left(|a_i|+|t_i|+\sum_{j:(j,i)\in G}|b_{ij}|\right)\le1
\right\}.
\]
Here \(V\) is the fixed finite population, \(n\) is its size, and \(\theta=(G,a,t,b)\) is a fixed graph-and-coefficient schedule, with the following conventions:
\begin{itemize}
\item \textbf{(Parameters.)} \(n\ge4\) and \(1\le d\le n-1\).
\item \textbf{(Graph.)} \(G\) is directed, simple, and off-diagonal; an arrow \((j,i)\) points from source \(j\) to recipient \(i\).
\item \textbf{(Bounds.)} The degree and coefficient-mass requirements are those of \cref{obj:ass:in-degree,obj:ass:out-degree,obj:ass:coefficient-mass}.
\end{itemize}
The coordinates \(a_i\), \(t_i\), and \(b_{ij}\) are, respectively, baseline, own-treatment, and source-to-recipient spillover coefficients. This is the interaction-order-one coefficient-mass model of \citet{CortezRodriguezEichhornYu2023SNIPE}, with published interaction order \(\beta=1\), restricted to compulsory diagonal coordinates after graph augmentation. The corresponding published total in-degree and out-degree bound is \(d+1\), while \(d\) counts off-diagonal arrows. The own-treatment coordinate is always available and unthinned. Off-diagonal arrows are observed through the audit channel specified in the design assumptions.
\end{definitionv}

The schedule class \leanref{sym:Mclass}{\ensuremath{\mathcal M_{n,d}}} in \cref{obj:def:model} uses separate notation for own-treatment effects and off-diagonal spillovers. Adding every diagonal arrow represents the same outcomes in the published neighborhood convention, with each degree increasing by one. The compulsory diagonal coordinates remain available even when their coefficients are zero. The coefficient correspondence and degree translation are given in \cref{obj:lem:published-model}.

The \leanref{S-2}{experimental design} combines treatment randomization with an edge audit on this fixed schedule. Write \leanref{sym:Z}{\ensuremath{Z}} for the random treatment-assignment vector. Its coordinates follow the independent half-Bernoulli design.

\begin{assumptionv}{ass:assignment}[Independent treatment assignment]
The treatment-assignment vector \(Z\), indexed by the finite population \(V\), has law
\[
\operatorname{Law}(Z)
=
\bigotimes_{j\in V}\operatorname{Bernoulli}(1/2).
\]
\end{assumptionv}

Independent unit Bernoulli randomization is the treatment-design condition used by SNIPE, here with treatment probability one half \citep{CortezRodriguezEichhornYu2023SNIPE}. Equal assignment probabilities give each treatment coordinate the same experimental variation. For subsequent score formulas, define the centered treatment signs \leanref{sym:X}{\ensuremath{X_k=2Z_k-1}} for \(k\in V\); \(X_i\) and \(X_j\) denote the own-treatment and source-treatment signs under the same convention.

Graph observation is governed by the audit-mark array \leanref{sym:W}{\ensuremath{W}}. The coordinate \(W_{ij}\) is the mark for the ordered pair \((j,i)\), and \leanref{sym:q}{\ensuremath{q\in[0,1]}} is the known edge-retention probability.

\begin{assumptionv}{ass:audit}[Independent edge audit]
The audit-mark array \(W\), indexed by ordered pairs of distinct units in the finite population \(V\), has law
\[
\operatorname{Law}(W)
=
\bigotimes_{(j,i)\in V^2:j\ne i}\operatorname{Bernoulli}(q),
\]
where \(q\) is the known edge-retention probability.
\end{assumptionv}

Independent homogeneous edge error supplies the measurement model; the audit adopts its directed, false-negative-only variant \citep{LiSussmanKolaczyk2021NetworkNoise}. Each true off-diagonal arrow is retained with probability \(q\), independently across ordered pairs, and every recorded arrow is a true arrow. Marks on pairs outside the true graph leave the recorded graph unchanged. The independence condition below makes this measurement mechanism exogenous to treatment randomization.

\begin{assumptionv}{ass:design-independence}[Assignment and audit independence]
The audit-mark array \(W\) and the treatment-assignment vector \(Z\) are independent:
\[
W\mathbin{\perp\!\!\!\perp}Z.
\]
\end{assumptionv}

Independence of network measurement error and treatment assignment makes the audit mechanism exogenous to the randomized treatments, as in the network-error design of \citep{LiSussmanKolaczyk2021NetworkNoise}. Conditional on a fixed schedule, \cref{obj:ass:assignment,obj:ass:audit,obj:ass:design-independence} therefore specify the joint design randomness. The assignment-and-record conventions below assemble that randomness into the retained graph \leanref{sym:H}{\ensuremath{H}}, the complete observed record \leanref{sym:O}{\ensuremath{O}}, and the law used to evaluate estimators.

\begin{definitionv}{synth_3}[Fixed-schedule experiment law]
Fix \(V=\{1,\ldots,n\}\), a schedule \(\theta=(G,a,t,b)\), and a known edge-retention probability \(q\in[0,1]\). Draw the treatment coordinates independently as
\[
Z_j\sim\operatorname{Bernoulli}(1/2),\qquad j\in V.
\]
Independently of \(Z\), draw mutually independent audit marks
\[
W_{ij}\sim\operatorname{Bernoulli}(q),
\qquad (j,i)\in V^2,\ j\ne i.
\]
Define the retained graph and complete observed record by
\[
H=\{(j,i)\in G:W_{ij}=1\},
\qquad
O=(H,Z,Y(Z)),
\]
where
\[
Y_i(Z)=a_i+t_iZ_i+\sum_{j:(j,i)\in G}b_{ij}Z_j,
\qquad i\in V.
\]
The record contains every retained-edge label, every treatment coordinate, and every noiseless realized outcome. Draw \(\xi\sim\operatorname{Uniform}[0,1]\) independently of \((Z,W)\) to represent optional estimator randomization. Define
\[
P_{\theta,q}=\operatorname{Law}(O,\xi)
\]
with \(\theta\) fixed. This law is taken on
\[
2^{\{(j,i)\in V^2:j\ne i\}}
\times\{0,1\}^V\times\mathbb R^V\times[0,1],
\]
with discrete sigma-algebras on the finite graph and assignment coordinates and Borel sigma-algebras on the outcome and seed coordinates.
\end{definitionv}

The experiment law \leanref{sym:P}{\ensuremath{P_{\theta,q}}} in \cref{obj:synth_3} holds the graph and coefficients fixed while averaging over assignment, audit, and optional estimator randomization. The independent uniform seed \leanref{sym:xi}{\ensuremath{\xi}} permits randomized rules, with deterministic rules represented by estimators that ignore it. Every treatment coordinate and realized outcome is observed, including those associated with arrows omitted by the audit. Accordingly, estimation can use the full joint information in the record.

We evaluate that estimation problem by expected squared error for the population contrast and take the worst case over the bounded schedule class.

\begin{definitionv}{def:risk}[Minimax squared risk $R_{n,d}(q)$]
\[
R_{n,d}(q)
=
\inf_T\sup_{\theta\in\mathcal M_{n,d}}\mathcal L(T;\theta,q),
\qquad
\mathcal L(T;\theta,q)
=
\mathbb E_{P_{\theta,q}}
\left[(T(O,\xi)-\tau(\theta))^2\right].
\]
The infimum ranges over measurable randomized estimators \(T\). Here \(O\) is the complete observed graph, assignment, and realized-outcome record, \(\tau(\theta)\) is the population all-treated-versus-all-control contrast, and \(P_{\theta,q}\) is the experiment law for the fixed schedule \(\theta\) and known edge-retention probability \(q\). The estimator-randomization seed \(\xi\) is uniform on \([0,1]\) and independent of the experiment.
\end{definitionv}

For a measurable estimator \leanref{sym:T}{\ensuremath{T}}, the loss \leanref{sym:L}{\ensuremath{\mathcal L(T;\theta,q)}} in \cref{obj:def:risk} averages its squared error under the fixed-schedule experiment. The minimax squared risk \leanref{sym:R}{\ensuremath{R_{n,d}(q)}} optimizes this worst-case loss over all measurable uses of the complete record and independent seed. Thus population size, off-diagonal degree, and known edge retention index a common precision criterion for the declared causal target and schedule class.


\section{Minimax precision and attainable estimation}

The precision criterion in \cref{obj:def:risk} concerns estimation of the population total treatment effect from the retained graph, treatment assignments, and realized outcomes. An attainable bound therefore begins with a score computed from those observations. Centering the binary assignments provides the treatment contrasts used in the score.

\begin{definitionv}{synth_2}[Centered treatment signs]
For a binary treatment-assignment vector \(Z\in\{0,1\}^V\), define its centered treatment-sign vector \(X\in\{-1,1\}^V\) by
\[
X_k=2Z_k-1,\qquad k\in V.
\]
Thus \(X_i=2Z_i-1\) is the own-treatment sign for recipient \(i\), and \(X_j=2Z_j-1\) is the treatment sign for source \(j\).
\end{definitionv}

For positive retention, the observable inverse-inclusion score \leanref{sym:Tq}{\ensuremath{T_q}} in \cref{obj:def:audit-score} combines each outcome with its own-treatment sign and the treatment signs of recorded sources. Weighting each recorded arrow by the reciprocal of its retention probability compensates for edge thinning under the independent audit design.

\begin{definitionv}{def:audit-score}[Inverse-inclusion score $T_q$]
\[
T_q
=
\frac2n\sum_{i\in V}Y_i(Z)
\left(X_i+q^{-1}\sum_{j:(j,i)\in H}X_j\right),
\qquad q>0.
\]
Here \(V\) is the population, \(n\) is its size, \(Z\) is the treatment-assignment vector, \(Y_i(Z)\) is unit \(i\)'s realized outcome, and \(H\) is the recorded retained-edge graph. The quantities \(X_i\) and \(X_j\) are the centered own-treatment and source-treatment signs, respectively, and \(q\) is the known edge-retention probability.
\end{definitionv}

The supplied-graph score \leanref{sym:TG}{\ensuremath{T_G}} in \cref{obj:def:oracle-score} provides the corresponding comparison when the true interference graph is available. It replaces the weighted retained-source sum with the sum over all true sources.

\begin{definitionv}{def:oracle-score}[Supplied-graph score $T_G$]
\[
T_G
=
\frac2n\sum_{i\in V}Y_i(Z)
\left(X_i+\sum_{j:(j,i)\in G}X_j\right).
\]
Here \(V\) is the population, \(n\) is its size, \(G\) is the supplied true interference graph, \(Z\) is the treatment-assignment vector, and \(Y_i(Z)\) is unit \(i\)'s realized outcome. The quantities \(X_i\) and \(X_j\) are the centered own-treatment and source-treatment signs, respectively.
\end{definitionv}

The audit decomposition in \cref{obj:lem:score-audit} separates assignment variation in the supplied-graph score from the additional variation generated by inverse-inclusion weighting. Conditional on the assignment, averaging the observable score over the audit recovers the supplied-graph score. Independent audit marks then isolate the extra variance attributable to edge collection. The resulting upper-risk envelope \leanref{sym:upper}{\ensuremath{u(n,d,q)}} in \cref{obj:def:upper-envelope} records these two components.

\begin{definitionv}{def:upper-envelope}[Upper-risk envelope $u(n,d,q)$]
\[
u(n,d,q)
=
\begin{cases}
5(d+1)^2/n+4d(1-q)/(nq),&q>0,\\
+\infty,&q=0.
\end{cases}
\]
Here \(n\) is the population size, \(d\) is the common off-diagonal in-degree and out-degree bound, and \(q\) is the known edge-retention probability.
\end{definitionv}

The first term reflects assignment variation under the common in-degree and out-degree bound. The second measures the variance cost of retaining each true arrow with probability \(q\); it decreases to zero as retention approaches one. These components give a prospective error bound in terms of population size, degree, and the collection probability.

To obtain a rule for every admissible retention probability, use the envelope to choose between a clipped score and zero. The coefficient-mass restriction in \cref{obj:ass:coefficient-mass} bounds the absolute population contrast by one, so clipping to that interval preserves the squared-error guarantee. The observable upper rule \leanref{sym:Tup}{\ensuremath{T^{\mathrm{up}}}} in \cref{obj:def:upper-rule} makes this choice explicit.

\begin{definitionv}{def:upper-rule}[Observable upper rule $T^{\mathrm{up}}$]
\[
T^{\mathrm{up}}
=
\begin{cases}
\max\{-1,\min\{1,T_q\}\},&q>0\text{ and }u(n,d,q)<1,\\
0,&\text{otherwise}.
\end{cases}
\]
Here \(T_q\) is the observable inverse-inclusion score, \(u(n,d,q)\) is the upper-risk envelope, \(n\) is the population size, \(d\) is the common off-diagonal in-degree and out-degree bound, and \(q\) is the known edge-retention probability.
\end{definitionv}

The score branch applies when its envelope improves upon the unit bound supplied by the zero rule. This construction yields the following uniform guarantee over the schedule class.

\begin{definitionv}{def:frontier-scale}[Risk scale $F(n,d,q)$]
Define the risk scale by
\[
F(n,d,q)=
\begin{cases}
\min\left\{1,\dfrac{d^2}{n[1-(1-q)^d]}\right\},&q>0,\\
1,&q=0.
\end{cases}
\]
Its domain is specified by
\begin{itemize}
\item \textbf{(Population size.)} \(n\ge4\).
\item \textbf{(Degree bound.)} \(1\le d\le n-1\).
\item \textbf{(Retention probability.)} \(q\in[0,1]\).
\end{itemize}
On this domain, \(F(n,d,q)\) is positive and finite.
\end{definitionv}

The guarantee combines variance control with the bounded range of the causal target. When assignment and audit variation are sufficiently small, the rule uses the inverse-inclusion score; when the envelope reaches one, the zero branch supplies a bounded worst-case squared loss. Thus the same observable construction covers both improving precision and saturation.

Matching lower bounds identify how this guarantee compares with the best measurable use of the complete record. The explicit risk scale \leanref{sym:F}{\ensuremath{F(n,d,q)}} in \cref{obj:def:frontier-scale} and the named attaining rule \leanref{sym:Tstar}{\ensuremath{T^\star_{n,d,q}}} in \cref{obj:def:frontier-rule} state the scale and estimator used in that characterization.

\begin{definitionv}{def:frontier-rule}[Attaining rule $T^\star_{n,d,q}$]
For every admissible \(n,d,q\), define the attaining rule on the observed record \(O\) by
\[
T^\star_{n,d,q}(O)=T^{\mathrm{up}}(O)=
\begin{cases}
\max\{-1,\min\{1,T_q(O)\}\},&q>0,\ u(n,d,q)<1,\\
0,&\text{otherwise},
\end{cases}
\]
where \(T_q\) is the observable inverse-inclusion score and \(u(n,d,q)\) is the upper-risk envelope.
\end{definitionv}

\begin{theoremv}{thm:observable-upper}[Observable upper risk bound]
Let \(V\) be a finite population with size \(n=|V|\). Suppose that:
\begin{itemize}
\item \textbf{(Population size.)} \(n\ge4\).
\item \textbf{(Degree bound.)} The integer degree bound \(d\) satisfies \(1\le d\le n-1\).
\item \textbf{(Retention probability.)} The known edge-retention probability \(q\) belongs to \([0,1]\).
\end{itemize}
Under the independent assignment--audit product design, the minimax risk \(R_{n,d}(q)\) and expected squared loss \(\mathcal L\) of \cref{obj:def:risk}, the observable rule \(T^{\mathrm{up}}\) of \cref{obj:def:upper-rule}, and the upper-risk envelope \(u(n,d,q)\) of \cref{obj:def:upper-envelope} satisfy
\[
R_{n,d}(q)
\le
\sup_{\theta\in\mathcal M_{n,d}}
\mathcal L(T^{\mathrm{up}};\theta,q)
\le
\min\{1,u(n,d,q)\},
\]
where \(\mathcal M_{n,d}\) is the fixed schedule class underlying the risk in \cref{obj:def:risk}. The rule \(T^{\mathrm{up}}\) is real-valued and measurable on the complete observed-record space.
\end{theoremv}

\begin{theoremv}{thm:observed-record-precision-frontier}[Observed-record precision frontier]
Let \(F\) be the risk-scale function in \cref{obj:def:frontier-scale}, and let the attaining rule be
\[
T^\star_{n,d,q}=T^{\mathrm{up}},
\]
where \(T^{\mathrm{up}}\) is the total observable upper rule. There exist universal real constants \(c,C_0>0\) such that, simultaneously for all parameters satisfying
\begin{itemize}
\item \textbf{(Population.)} The population size \(n\) is an integer with \(n\ge4\).
\item \textbf{(Degree.)} The degree bound \(d\) is an integer with \(1\le d\le n-1\).
\item \textbf{(Retention.)} The known retention probability \(q\) belongs to \([0,1]\).
\end{itemize}
the canonical independent assignment--audit design satisfies
\[
cF(n,d,q)
\le R_{n,d}(q)
\le
\sup_{\theta\in\mathcal M_{n,d}}
\mathcal L(T^\star_{n,d,q};\theta,q)
\le C_0F(n,d,q).
\]
Here \(\mathcal M_{n,d}\) is the fixed schedule class, \(\mathcal L(T;\theta,q)\) is expected squared estimation loss, and \(R_{n,d}(q)\) is its infimum over all measurable real-valued randomized estimators \(T(O,\xi)\), where \(O\) is the complete observed record and \(\xi\) is the independent randomization seed.

For every pair of degree and retention sequences \(d_n,q_n\) satisfying
\[
1\le d_n\le n-1,\qquad q_n\in[0,1]
\qquad\text{for every integer }n\ge4,
\]
there exists a sequence of measurable real-valued randomized estimators \(T_n(O,\xi)\) such that
\[
\sup_{\theta\in\mathcal M_{n,d_n}}
\mathcal L(T_n;\theta,q_n)\longrightarrow0
\qquad\text{as }n\longrightarrow\infty
\]
if and only if
\[
\frac{d_n^2}{n}\longrightarrow0
\qquad\text{and}\qquad
\begin{cases}
+\infty,&q_n=0,\\[2pt]
\dfrac{d_n}{nq_n},&q_n>0
\end{cases}
\longrightarrow0.
\]
The second convergence is in the extended nonnegative real numbers.

For every integer \(n\ge4\) and \(1\le d\le n-1\), the endpoint scales are
\[
F(n,d,0)=1,
\qquad
F(n,d,1)=\min\left\{1,\frac{d^2}{n}\right\}.
\]
Moreover, there exist universal real constants \(c_1,C_1>0\) such that, simultaneously for every integer \(n\ge4\), every integer \(1\le d\le n-1\), and every \(0<q\le1\),
\[
c_1\min\left\{1,\frac{d^2}{n}+\frac{d}{nq}\right\}
\le F(n,d,q)
\le
C_1\min\left\{1,\frac{d^2}{n}+\frac{d}{nq}\right\}.
\]
\end{theoremv}

The characterization separates two sources of imprecision: the degree component \(d^2/n\) and the edge-collection component \(d/(nq)\). The universal lower and upper constants \leanref{sym:c}{\ensuremath{c}} and \leanref{sym:C0}{\ensuremath{C_0}} apply simultaneously across the stated parameter domain. Consequently, the observable rule attains the minimax scale even when competing estimators use all assignment and outcome coordinates and independent randomization.

The factor \leanref{sym:p}{\ensuremath{1-(1-q)^d}} is the probability that at least one of \(d\) independently audited arrows is retained. Its role in the scale describes the collection transition: when \(dq\) is small, edge collection determines the leading precision cost; once \(dq\) is of constant order or larger, the degree component determines the scale up to universal factors. The transition occurs around retention of order \(1/d\). The minimum with one expresses saturation of worst-case squared risk: whenever the uncapped scale is at least one, minimax risk is bounded above and below by positive universal constants.

For admissible degree sequences \leanref{sym:dseq}{\ensuremath{d_n}} and retention sequences \leanref{sym:qseq}{\ensuremath{q_n}}, the consistency statement gives an exact joint requirement. Degree must grow more slowly than the square root of population size, and retention must be positive eventually with \(nq_n/d_n\) diverging. Under these conditions, the named attaining rules have uniformly vanishing expected squared loss. Necessity follows from the lower bound over all measurable randomized estimators, so these conditions characterize the estimation problem under the declared design and schedule class.

The converse uses blocks of recipients that share response values while their source allocations are only partly revealed by the audit. Before treatment and audit draws, each block receives a fixed but unknown baseline and a random allocation of sources. The observed outcomes can therefore teach the estimator about hidden associations, so the comparison retains repeated responses and conditions the likelihood on the observed global treated-source count. Translation control, Walsh energy, and hidden-allocation contraction show that opposite-sign block schedules remain difficult to distinguish even with that outcome-driven learning; see \cref{obj:lem:block-law,obj:lem:baseline-translation-affinity,obj:lem:walsh-channel-energy,obj:lem:hidden-allocation-contraction,obj:thm:block-testing-scale}. The two-prior reduction in \cref{obj:lem:full-record-two-prior-risk} then turns this complete-record testing bound into the minimax lower bound. Endpoint and explicit-constant deductions are collected in \cref{obj:thm:supported-boundaries,obj:thm:frontier-answer}.


\section{Edge collection and the supplied-graph benchmark}

The characterization in \cref{obj:thm:observed-record-precision-frontier} gives calibrated edge retention a prospective interpretation: the known retention probability \(q\) is a collection choice whose precision consequences depend on population size and the common degree bound. Under the independent audit design in \cref{obj:ass:audit,obj:ass:design-independence}, increasing retention reduces the contribution of edge collection to squared risk. The comparison between \(d/(nq)\) and \(d^2/n\) places the transition at retention of order \(1/d\). At or above this order, the minimax risk has the same order as the full-retention scale \(\min\{1,d^2/n\}\). Below the transition, the collection term determines the uncapped risk scale.

Population size determines how this threshold translates into a requirement for consistent estimation. For the admissible sequences in \cref{obj:thm:observed-record-precision-frontier}, uniform consistency is equivalent to
\[
d_n=o(\sqrt n),
\qquad
q_n>0\ \text{eventually},
\qquad
\frac{nq_n}{d_n}\longrightarrow\infty.
\]
Thus a collection plan can use a vanishing retention probability provided it exceeds \(d_n/n\) by a diverging factor and the degree condition is satisfied. Retention of order \(1/d_n\) or greater suffices to match the full-retention precision order. When \(d_n^2/n\to0\), consistent collection plans can have retention below this transition. For a fixed degree bound, for example, consistency requires \(nq_n\to\infty\), while retention bounded away from zero attains the full-retention order.

The upper-risk envelope in \cref{obj:def:upper-envelope} supplies a complementary planning calculation. Its assignment component is \(5(d+1)^2/n\), and its audit component is \(4d(1-q)/(nq)\). The latter vanishes at full retention. Accordingly, the envelope makes the benefit of additional collection explicit, while the matching minimax characterization determines the attainable order over all measurable uses of the observed record.

To isolate the degree component, consider an experiment that supplies the true interference graph in addition to the original observations. The supplied-graph minimax squared risk \leanref{sym:RG}{\ensuremath{R^G_{n,d}(q)}} in \cref{obj:def:supplied-risk} retains the schedule class, causal target, assignment law, and estimator-randomization convention of the original decision problem.

\begin{definitionv}{def:supplied-risk}[Supplied-graph risk $R^G_{n,d}(q)$]
For fixed admissible \(n,d,q\), define the supplied-graph minimax squared risk by
\[
R^G_{n,d}(q)=
\inf_{T^G}\sup_{\theta\in\mathcal M_{n,d}}
\mathbb E_{P_{\theta,q}}
\left[
\bigl(T^G(O,G(\theta),\xi)-\tau(\theta)\bigr)^2
\right].
\]
Here \(\mathcal M_{n,d}\) is the schedule class, \(P_{\theta,q}\) is the fixed-schedule experiment law, \(O\) is the original observed record, \(\tau(\theta)\) is the population all-treated-versus-all-control contrast, and \(\xi\) is an independent uniform seed on \([0,1]\). The additional observation \(G(\theta)\) is the true off-diagonal graph component of the fixed schedule \(\theta\).

The infimum ranges over all Borel-measurable functions
\[
T^G:
\operatorname{range}(O)\times
2^{\{(j,i)\in V^2:j\ne i\}}\times[0,1]
\longrightarrow\mathbb R,
\]
where \(V\) is the fixed finite population. The finite graph coordinate carries the discrete sigma-algebra. The original record carries the Borel sigma-algebra inherited from its finite graph and assignment coordinates and its real outcome coordinates. The estimator may use the known \(n,d,q\). Expectations of the nonnegative loss take values in \([0,+\infty]\), and the estimator class includes every such measurable rule.
\end{definitionv}

Supplying the graph makes every true source-to-recipient arrow available to the estimator; the outcome coefficients and population contrast remain quantities to be learned from the randomized outcomes. The quadratic benchmark of \citet[Theorem \texttt{thm:degree-frontier}; non-archival research note]{CausalSmithKnownGraph2026} is translated to the present additive conventions in \cref{obj:thm:supported-boundaries}. After adjoining the compulsory diagonal coordinate as in \cref{obj:lem:published-model}, the supplied-graph score \(T_G\) is exactly additive, half-Bernoulli SNIPE: its own-treatment sign is the diagonal summand and its true-source signs are the off-diagonal summands. The present work adds inverse-inclusion correction for sampled arrows and analyzes that rule jointly with the unrestricted complete-record converse. Under the same schedule class in \cref{obj:def:model} and assignment law in \cref{obj:ass:assignment}, the supplied-graph lower bound and score calculation in \cref{obj:lem:supplied-block-record-lower,obj:lem:score-audit} give
\[
R^G_{n,d}(q)\asymp \min\left\{1,\frac{d^2}{n}\right\}
\]
uniformly over admissible \(n,d,q\), with universal comparison constants. The upper bound uses the supplied-graph score in \cref{obj:def:oracle-score}, clipped to the target range, together with the zero rule when appropriate. The lower bound identifies the same degree dependence even with the true graph supplied. This comparison therefore attributes the degree component to estimation under the fixed assignment law and schedule restrictions, while the additional retention dependence measures the precision cost of sampled graph information.

For a platform workflow based on sampled interaction records, this interpretation ties collection to explicit sampling and influence conditions. The experiment in \cref{obj:synth_3} requires independent arrow retention at a known common probability, independent half-Bernoulli treatment assignment, and observation of every assignment and realized outcome. Its influence conditions are the additive response model, both degree bounds, and the coefficient-mass restriction in \cref{obj:def:model}. The approximate-network analysis of \citet[Assumptions 1--3; Theorems 1--2; Lemma 2]{JiangDengWangWang2024ApproximateNetworks} instead gives variance \(O\{d_{\mathcal G}^2/[n\pi(1-\pi)]\}\) for its pseudoinverse estimator under bounded outcomes and first- and second-order influence restrictions, with an omitted-influence condition controlling relative bias. Here, the independent retention experiment supports inverse-inclusion correction, a complete-record converse, and exact joint consistency conditions, making population size, degree, and collection probability inputs to an unrestricted minimax planning result.


\section{Limitations and future work}

The uniform squared-risk comparison in \cref{obj:thm:observed-record-precision-frontier} characterizes attainable precision for a single experiment under the schedule restrictions in \cref{obj:def:model} and the observation design in \cref{obj:ass:assignment,obj:ass:audit,obj:ass:design-independence}. Its comparison constants establish the minimax order across the admissible population sizes, degree bounds, and retention probabilities. Exact minimax constants remain a separate question, as does uncertainty quantification with specified confidence-interval coverage. A squared-risk guarantee controls expected estimation error; constructing informative intervals requires an additional coverage argument that accounts for assignment variation and sampled graph information. The behavior of such intervals near regimes in which the risk remains bounded away from zero warrants particular attention.

Known retention is central to the inverse-inclusion score in \cref{obj:def:audit-score}. When retention probabilities are estimated from calibration data, their estimation error enters the weights and can affect both bias and variance, especially at low retention. An extension would need to specify the calibration experiment, its relationship to the treatment experiment, and conditions ensuring sufficiently accurate estimation of inclusion probabilities. Heterogeneous retention probabilities present a related problem: calibration and precision guarantees would have to account for which arrows are sampled less frequently and how their spillover coefficients contribute to the target.

The audit design also distinguishes missing arrows from falsely recorded arrows. Under \cref{obj:ass:audit}, the recorded graph is formed by retaining true arrows independently. Allowing false-positive edges changes the observation law and the behavior of the score. A useful extension would specify separate inclusion and false-positive mechanisms, identify which aspects of those mechanisms are known or estimable, and determine how their uncertainty affects total-effect estimation.

Connecting the calibrated audit experiment to a platform collection mechanism requires calibration at the level of directed interference arrows. For the present guarantees to apply, the collection procedure must induce the common known inclusion probability, independent arrow retention, and independence from treatment assignment required by \cref{obj:ass:audit,obj:ass:design-independence}. It must also record assignments and realized outcomes for the full study population, while the interference schedule satisfies \cref{obj:def:model}. Sampling interaction events can induce dependent or heterogeneous arrow inclusion when several events represent the same arrow. Translating such a workflow into the present experiment therefore requires an explicit mapping from event records to interference arrows and justification of the resulting inclusion law. Treatment-dependent activity or outcome-dependent recording requires a different observation model.

Alternative assignments and repeated experiments offer further directions. Unequal treatment probabilities, complete randomization, and cluster assignment alter both the score construction and the dependence relevant to precision bounds. Repeated experiments would require specifying whether the graph and response coefficients remain fixed, how assignments vary across rounds, and whether graph audits are shared or renewed. Their potential precision gains depend on these choices. Richer response models, including treatment interactions and nonlinear exposure responses, likewise require an expanded schedule class and restrictions controlling how joint treatment changes affect outcomes. The additive coefficient-mass restriction in \cref{obj:ass:coefficient-mass} provides the starting point for such extensions.

Graph recovery is a separate inferential objective from estimating the population total treatment effect. Recovering individual arrows requires a recovery criterion and identifying restrictions connecting arrow presence to observable responses, together with an observation design that supplies sufficient information about individual associations. Studying those requirements alongside total-effect precision would clarify when collection sufficient for a population contrast also supports learning the interference structure.


\appendix

\section*{Appendices}

\section{Observable estimation and supplied-graph comparison}

The observable upper guarantee rests on a separation of treatment-assignment variation from edge-audit variation. The audit-score moment identity makes that separation exact, while the supplied-graph comparison identifies a component of minimax risk that persists throughout the retention range. The correspondence with the published additive model then fixes the degree convention and observation channel needed to interpret these comparisons.

\subsection{Audit-score moments and observable estimation}

For a fixed schedule, the audit score in \cref{obj:def:audit-score} weights recorded spillover coordinates by their inverse inclusion probability. Its assignment-variance benchmark is the oracle score in \cref{obj:def:oracle-score}, which uses the true graph. The following moment statement connects these scores to the population causal target.

\begin{lemmav}{lem:score-audit}[Audit-score moments and oracle bound]
Let \(V\) be a finite population with \(n=|V|\), and let \(\theta=(G,a,t,b)\) be a fixed schedule. Write \(Y_i(z)\) for unit \(i\)'s additive potential outcome under binary assignment \(z\). Suppose:
\begin{itemize}
\item \textbf{(Population and degree.)} \(n\ge4\), and \(d\) is an integer satisfying \(1\le d\le n-1\).
\item \textbf{(Schedule class.)} \(\theta\in\mathcal M_{n,d}\) as in \cref{obj:def:model}.
\item \textbf{(Assignment.)} The treatment-assignment vector \(Z\) satisfies \cref{obj:ass:assignment}.
\item \textbf{(Audit.)} The off-diagonal audit-mark array \(W\) satisfies \cref{obj:ass:audit} with retention probability \(q>0\).
\item \textbf{(Design independence.)} \(Z\) and \(W\) satisfy \cref{obj:ass:design-independence}.
\end{itemize}
Let \(P_{\theta,q}\) be the induced experiment law for this fixed schedule, with the observed record retaining precisely the audited arrows of \(G\), the assignment, and the realized outcomes. Let \(T_q\) and \(T_G\) be the scores in \cref{obj:def:audit-score,obj:def:oracle-score}, respectively, and define the all-treated-versus-all-control population contrast by
\[
\tau(\theta)
=\frac1n\sum_{i\in V}
\bigl(Y_i(\mathbf1)-Y_i(\mathbf0)\bigr),
\]
where \(\mathbf1\) and \(\mathbf0\) are the all-treated and all-control assignments. Then
\[
\mathbb E_{P_{\theta,q}}[T_q]=\tau(\theta),
\]
\[
\operatorname{Var}_{P_{\theta,q}}(T_q)
=
\operatorname{Var}_{Z}(T_G)
+\frac{4(1-q)}{n^2q}
\sum_{i\in V}
\bigl|\{j\in V:(j,i)\in G\}\bigr|
\,\mathbb E_Z\!\left[Y_i(Z)^2\right],
\]
and
\[
\operatorname{Var}_{Z}(T_G)\le\frac{5(d+1)^2}{n},
\]
where expectations and variances indexed by \(Z\) use the independent half-Bernoulli assignment law.
\end{lemmav}

\begin{proof}[Proof of \cref{obj:lem:score-audit}]
\begin{enumerate}
\item Fix the schedule and design in the hypotheses. For \(i,j\in V\) and \(z\in\{0,1\}^{V}\), write
\[
N_i^-:=\{j\in V:(j,i)\in G\},
\qquad
X_j(z):=2z_j-1,
\qquad
X_j:=X_j(Z).
\]
The additive outcome map is
\[
Y_i(z)=a_i+t_i z_i+\sum_{j\in N_i^-}b_{ij}z_j.
\]
Index an audit mark by its recipient and source:
\[
W_{ij}:=W(j,i)\quad(j\ne i),
\qquad
H:=\{(j,i)\in G:W_{ij}=1\}.
\]
By \cref{obj:ass:assignment,obj:ass:audit,obj:ass:design-independence},
\[
0<q\le1,
\qquad
\operatorname{Law}(Z,W)
=
\left(\bigotimes_{j\in V}\operatorname{Bernoulli}(1/2)\right)
\otimes
\left(\bigotimes_{\substack{(j,i)\in V^2\\j\ne i}}
\operatorname{Bernoulli}(q)\right).
\]
We use \(\mathbb E_W\) for expectation over the audit marks with the assignment fixed, and \(\mathbb E_Z\) for expectation over the assignment. All the sums and expectations below are finite.

The evaluations of the scores in
\cref{obj:def:audit-score,obj:def:oracle-score} are
\[
\begin{aligned}
T_q(z,W)
&=\frac2n\sum_{i\in V}Y_i(z)
 \left(X_i(z)+q^{-1}\sum_{j\in N_i^-}W_{ij}X_j(z)\right),\\
T_G(z)
&=\frac2n\sum_{i\in V}Y_i(z)
 \left(X_i(z)+\sum_{j\in N_i^-}X_j(z)\right),
\end{aligned}
\qquad
T_q=T_q(Z,W),\quad T_G=T_G(Z).
\]
Since \(\mathbb E_W[W_{ij}]=q\), for every assignment \(z\),
\[
\mathbb E_W\!\left[\sum_{j\in N_i^-}W_{ij}X_j(z)\right]
=q\sum_{j\in N_i^-}X_j(z),
\qquad
\mathbb E_W[T_q(z,W)]=T_G(z).
\]

Under the half-Bernoulli assignment law, for \(j,k\in V\),
\[
\mathbb E_Z[X_j]=0,
\qquad
\mathbb E_Z[Z_kX_j]
=
\begin{cases}
\frac12,&k=j,\\
0,&k\ne j.
\end{cases}
\]
The first case follows from \(Z_j(2Z_j-1)=Z_j\); the second follows from independence and the zero mean of \(X_j\). Expanding the additive outcome therefore gives
\[
\mathbb E_Z[Y_i(Z)X_j]
=
\frac12\left(
t_i\mathbf1_{\{i=j\}}
+\sum_{k\in N_i^-}b_{ik}\mathbf1_{\{k=j\}}
\right).
\]
Because \(i\notin N_i^-\), summing this identity over the own coordinate and the incoming sources yields
\[
\mathbb E_Z\!\left[
Y_i(Z)\left(X_i+\sum_{j\in N_i^-}X_j\right)
\right]
=
\frac12\left(t_i+\sum_{j\in N_i^-}b_{ij}\right).
\]
Also, evaluating the outcome map at the two constant assignments gives
\[
Y_i(\mathbf1)-Y_i(\mathbf0)
=t_i+\sum_{j\in N_i^-}b_{ij}.
\]
Consequently,
\[
\mathbb E_Z[T_G]
=\frac1n\sum_{i\in V}\left(t_i+\sum_{j\in N_i^-}b_{ij}\right)
=\tau(\theta),
\qquad
\mathbb E_{P_{\theta,q}}[T_q]
=\mathbb E_Z\!\left[\mathbb E_W[T_q(Z,W)]\right]
=\tau(\theta).
\]
% lean: CausalSmith.Experimentation.ThinnedgraphAdditiveRiskFrontier.integral_auditScore_eq_tte

\item For a fixed assignment \(z\), subtracting the two score evaluations gives
\[
T_q(z,W)-T_G(z)
=
\frac{2}{nq}\sum_{i\in V}\sum_{j\in N_i^-}
Y_i(z)X_j(z)(W_{ij}-q).
\]
Each centered mark has mean zero and second moment
\[
\mathbb E_W[W_{ij}-q]=0,
\qquad
\mathbb E_W[(W_{ij}-q)^2]
=q(1-q)^2+(1-q)q^2
=q(1-q).
\]
For distinct off-diagonal pairs, independence gives
\[
\mathbb E_W[(W_{ij}-q)(W_{i'j'}-q)]=0
\qquad\bigl((j,i)\ne(j',i')\bigr).
\]
Thus variance is unchanged by subtracting the constant \(T_G(z)\), and the variance of the resulting weighted sum is
\[
\begin{aligned}
\operatorname{Var}_W(T_q(z,W))
&=\operatorname{Var}_W(T_q(z,W)-T_G(z))\\
&=\frac{4}{n^2q^2}\,q(1-q)
 \sum_{i\in V}\sum_{j\in N_i^-}Y_i(z)^2X_j(z)^2\\
&=\frac{4(1-q)}{n^2q}
 \sum_{i\in V}|N_i^-|\,Y_i(z)^2,
\end{aligned}
\]
where the last equality uses \(X_j(z)^2=1\). Combining this with the conditional mean already established gives the conditional second-moment identity
\[
\mathbb E_W[T_q(z,W)^2]
=
T_G(z)^2
+\frac{4(1-q)}{n^2q}
 \sum_{i\in V}|N_i^-|\,Y_i(z)^2.
\]
Empty incoming neighborhoods contribute zero. These identities also include \(q=1\), when every audit mark equals one almost surely.
% lean: CausalSmith.Experimentation.ThinnedgraphAdditiveRiskFrontier.integral_auditScore_sq_given_assignment

\item Integrating the preceding second-moment identity over the assignment gives
\[
\mathbb E_{P_{\theta,q}}[T_q^2]
=
\mathbb E_Z[T_G^2]
+\frac{4(1-q)}{n^2q}
 \sum_{i\in V}|N_i^-|\,\mathbb E_Z[Y_i(Z)^2].
\]
Both scores have mean \(\tau(\theta)\). Subtracting its square from the two second moments therefore yields
\[
\begin{aligned}
\operatorname{Var}_{P_{\theta,q}}(T_q)
&=\mathbb E_Z[T_G^2]-\tau(\theta)^2
 +\frac{4(1-q)}{n^2q}
  \sum_{i\in V}|N_i^-|\,\mathbb E_Z[Y_i(Z)^2]\\
&=\operatorname{Var}_Z(T_G)
 +\frac{4(1-q)}{n^2q}
  \sum_{i\in V}|N_i^-|\,\mathbb E_Z[Y_i(Z)^2].
\end{aligned}
\]
% lean: CausalSmith.Experimentation.ThinnedgraphAdditiveRiskFrontier.variance_auditScore_eq_oracle_add

\item To bound the oracle variance, augment each incoming neighborhood by its own coordinate. Define
\[
N_i:=\{i\}\cup N_i^-,
\qquad
\widetilde b_{ij}:=
\begin{cases}
t_i,&j=i,\\
b_{ij},&j\in N_i^-,
\end{cases}
\quad(i\in V,\ j\in N_i),
\qquad
\mu_i:=a_i+\frac12\sum_{j\in N_i}\widetilde b_{ij}.
\]
The two cases defining \(\widetilde b_{ij}\) are disjoint because the graph is off-diagonal. Substituting \(z_j=(1+X_j(z))/2\) into the outcome map gives
\[
Y_i(z)=\mu_i+\frac12\sum_{j\in N_i}\widetilde b_{ij}X_j(z),
\qquad
\tau(\theta)=\frac1n\sum_{i\in V}\sum_{j\in N_i}\widetilde b_{ij}.
\]
Separating the diagonal products in each row gives
\[
\begin{aligned}
\left(\sum_{j\in N_i}\widetilde b_{ij}X_j(z)\right)
\left(\sum_{k\in N_i}X_k(z)\right)
&=\sum_{j\in N_i}\widetilde b_{ij}\\
&\quad+\sum_{j\in N_i}\sum_{k\in N_i\setminus\{j\}}
\widetilde b_{ij}X_j(z)X_k(z),
\end{aligned}
\]
since \(X_j(z)^2=1\). Define, for every \(z\in\{0,1\}^{V}\),
\[
\begin{aligned}
T_{G,\mathrm{lin}}(z)
&:=\frac2n\sum_{i\in V}\mu_i\sum_{j\in N_i}X_j(z),\\
T_{G,\mathrm{quad}}(z)
&:=\frac1n\sum_{i\in V}\sum_{j\in N_i}
 \sum_{k\in N_i\setminus\{j\}}\widetilde b_{ij}X_j(z)X_k(z),
\end{aligned}
\qquad
\begin{aligned}
T_{G,\mathrm{lin}}&:=T_{G,\mathrm{lin}}(Z),\\
T_{G,\mathrm{quad}}&:=T_{G,\mathrm{quad}}(Z).
\end{aligned}
\]
Expanding the oracle score and subtracting its constant coefficient now gives
\[
T_G(z)-\tau(\theta)
=T_{G,\mathrm{lin}}(z)+T_{G,\mathrm{quad}}(z).
\]
% lean: CausalSmith.Experimentation.ThinnedgraphAdditiveRiskFrontier.oracleScore_centered_expansion

\item Since \(\mathbb E_Z[T_G]=\tau(\theta)\), the expansion gives
\[
\operatorname{Var}_Z(T_G)
=
\mathbb E_Z[T_{G,\mathrm{lin}}^2]
+\mathbb E_Z[2T_{G,\mathrm{lin}}T_{G,\mathrm{quad}}]
+\mathbb E_Z[T_{G,\mathrm{quad}}^2].
\]
Collecting the linear terms by their source coordinate gives
\[
T_{G,\mathrm{lin}}(z)
=
\frac2n\sum_{j\in V}
\left(\sum_{\substack{i\in V\\j\in N_i}}\mu_i\right)X_j(z).
\]

For any subsets \(E_1,E_2\subseteq V\), independence and centering imply
\[
\mathbb E_Z\!\left[
\left(\prod_{j\in E_1}(Z_j-\tfrac12)\right)
\left(\prod_{k\in E_2}(Z_k-\tfrac12)\right)
\right]
=
\begin{cases}
4^{-|E_1|},&E_1=E_2,\\
0,&E_1\ne E_2.
\end{cases}
\]
Indeed, when the sets agree, every factor is a square equal to \(1/4\). When they differ, a coordinate in their symmetric difference contributes an independent factor with mean zero. Using
\[
\prod_{j\in E_1}X_j
=2^{|E_1|}\prod_{j\in E_1}(Z_j-\tfrac12)
\]
and the corresponding identity for \(E_2\), we obtain
\[
\mathbb E_Z\!\left[
\left(\prod_{j\in E_1}X_j\right)
\left(\prod_{k\in E_2}X_k\right)
\right]
=\mathbf1_{\{E_1=E_2\}}.
\]
The empty product equals one, so these formulas include empty sets.

In particular,
\[
\mathbb E_Z[X_\ell X_jX_k]=0
\qquad(\ell,j,k\in V,\ j\ne k),
\]
because the singleton \(\{\ell\}\) differs from the two-element set \(\{j,k\}\). Expanding the collected linear term against the quadratic term expresses every summand of their cross moment as a constant times such a triple moment. Hence
\[
\mathbb E_Z[2T_{G,\mathrm{lin}}T_{G,\mathrm{quad}}]=0,
\qquad
\operatorname{Var}_Z(T_G)
=\mathbb E_Z[T_{G,\mathrm{lin}}^2]
+\mathbb E_Z[T_{G,\mathrm{quad}}^2].
\]
% lean: CausalSmith.Experimentation.ThinnedgraphAdditiveRiskFrontier.oracle_mixed_second_moment_zero

\item Set
\[
d_+:=d+1.
\]
The degree and coefficient-mass bounds in \cref{obj:def:model} imply
\[
|N_i|=1+|N_i^-|\le d_+,
\qquad
\left|\{i\in V:j\in N_i\}\right|
=1+\left|\{i\in V:(j,i)\in G\}\right|
\le d_+,
\]
and
\[
|a_i|+\sum_{j\in N_i}|\widetilde b_{ij}|
=
|a_i|+|t_i|+\sum_{j\in N_i^-}|b_{ij}|
\le1.
\]
In particular,
\[
\sum_{j\in N_i}|\widetilde b_{ij}|\le1,
\qquad
|\mu_i|
\le |a_i|+\frac12\left|\sum_{j\in N_i}\widetilde b_{ij}\right|
\le |a_i|+\sum_{j\in N_i}|\widetilde b_{ij}|
\le1.
\]

The sign orthogonality above, applied to singleton sets, gives
\[
\mathbb E_Z[T_{G,\mathrm{lin}}^2]
=
\frac4{n^2}\sum_{j\in V}
\left(\sum_{\substack{i\in V\\j\in N_i}}\mu_i\right)^2.
\]
For each \(j\), finite Cauchy--Schwarz and the coordinate multiplicity bound yield
\[
\left(\sum_{\substack{i\in V\\j\in N_i}}\mu_i\right)^2
\le
\left|\{i\in V:j\in N_i\}\right|
\sum_{\substack{i\in V\\j\in N_i}}\mu_i^2
\le
d_+\sum_{\substack{i\in V\\j\in N_i}}\mu_i^2.
\]
Summing and reversing the finite sums gives
\[
\begin{aligned}
\sum_{j\in V}
\left(\sum_{\substack{i\in V\\j\in N_i}}\mu_i\right)^2
&\le d_+\sum_{j\in V}\sum_{\substack{i\in V\\j\in N_i}}\mu_i^2\\
&=d_+\sum_{i\in V}|N_i|\mu_i^2\\
&\le d_+\sum_{i\in V}d_+
=nd_+^2.
\end{aligned}
\]
Since \(n\ge4\), scaling by \(4/n^2\) therefore gives
\[
\mathbb E_Z[T_{G,\mathrm{lin}}^2]\le\frac{4d_+^2}{n}.
\]
% lean: CausalSmith.Experimentation.ThinnedgraphAdditiveRiskFrontier.oracle_linear_scaled_second_moment_le

\item Define the collected ordered-pair coefficients for all \(j,k\in V\) by
\[
b^{\mathrm{pair}}_{jk}
:=
\sum_{\substack{i\in V\\j\in N_i,\ k\in N_i\setminus\{j\}}}
\widetilde b_{ij}.
\]
An empty sum equals zero; in particular,
\[
b^{\mathrm{pair}}_{jj}=0.
\]
Reversing the finite sums in the quadratic term gives
\[
nT_{G,\mathrm{quad}}
=\sum_{j\in V}\sum_{k\in V}b^{\mathrm{pair}}_{jk}X_jX_k.
\]
For \(j,k,j',k'\in V\) with \(j\ne k\) and \(j'\ne k'\), sign orthogonality gives
\[
\mathbb E_Z[X_jX_kX_{j'}X_{k'}]
=
\begin{cases}
1,&(j',k')=(j,k)\ \text{or}\ (j',k')=(k,j),\\
0,&\text{otherwise}.
\end{cases}
\]
Indeed, the two products correspond to the sets \(\{j,k\}\) and \(\{j',k'\}\), which agree precisely in those two orientations. Thus, for \(j\ne k\),
\[
\mathbb E_Z\!\left[
X_jX_k\sum_{j'\in V}\sum_{k'\in V}
b^{\mathrm{pair}}_{j'k'}X_{j'}X_{k'}
\right]
=b^{\mathrm{pair}}_{jk}+b^{\mathrm{pair}}_{kj}.
\]
Expanding the square and using \(b^{\mathrm{pair}}_{jj}=0\) for its diagonal terms consequently yields
\[
\mathbb E_Z[T_{G,\mathrm{quad}}^2]
=
\frac1{n^2}\sum_{j\in V}\sum_{k\in V}
b^{\mathrm{pair}}_{jk}
\left(b^{\mathrm{pair}}_{jk}+b^{\mathrm{pair}}_{kj}\right).
\]
% lean: CausalSmith.Experimentation.ThinnedgraphAdditiveRiskFrontier.oracle_quadratic_second_moment

\item The incidence condition defining \(b^{\mathrm{pair}}_{jk}\) is unchanged by swapping \(j\) and \(k\). Hence
\[
b^{\mathrm{pair}}_{jk}+b^{\mathrm{pair}}_{kj}
=
\sum_{\substack{i\in V\\j\in N_i,\ k\in N_i\setminus\{j\}}}
(\widetilde b_{ij}+\widetilde b_{ik}).
\]
The set of recipients in this sum is contained in
\(\{i\in V:j\in N_i\}\), whose size is at most \(d_+\). Cauchy--Schwarz therefore gives, for every \(j,k\in V\),
\[
\left(b^{\mathrm{pair}}_{jk}+b^{\mathrm{pair}}_{kj}\right)^2
\le
d_+\sum_{\substack{i\in V\\j\in N_i,\ k\in N_i\setminus\{j\}}}
(\widetilde b_{ij}+\widetilde b_{ik})^2.
\]
For \(j=k\), both sides are zero.

For \(j\in N_i\) and \(k\in N_i\setminus\{j\}\), the coefficient-mass bound gives
\[
|\widetilde b_{ij}|+|\widetilde b_{ik}|
\le\sum_{\ell\in N_i}|\widetilde b_{i\ell}|
\le1.
\]
The triangle inequality then gives the rowwise pair estimate
\[
(\widetilde b_{ij}+\widetilde b_{ik})^2
\le
\left(|\widetilde b_{ij}|+|\widetilde b_{ik}|\right)^2
\le
|\widetilde b_{ij}|+|\widetilde b_{ik}|.
\]
Summing this estimate and enlarging each inner sum by its nonnegative diagonal term yields
\[
\begin{aligned}
\sum_{j\in N_i}\sum_{k\in N_i\setminus\{j\}}
(\widetilde b_{ij}+\widetilde b_{ik})^2
&\le
\sum_{j\in N_i}\sum_{k\in N_i\setminus\{j\}}
\left(|\widetilde b_{ij}|+|\widetilde b_{ik}|\right)\\
&\le
\sum_{j\in N_i}\sum_{k\in N_i}
\left(|\widetilde b_{ij}|+|\widetilde b_{ik}|\right)\\
&=2|N_i|\sum_{j\in N_i}|\widetilde b_{ij}|\\
&\le2d_+.
\end{aligned}
\]
When \(N_i=\{i\}\), the sums over distinct coordinates are empty and the same bound applies.

Summing the pairwise Cauchy--Schwarz inequalities and reindexing by recipient now gives
\[
\begin{aligned}
\sum_{j\in V}\sum_{k\in V}
\left(b^{\mathrm{pair}}_{jk}+b^{\mathrm{pair}}_{kj}\right)^2
&\le
d_+\sum_{i\in V}\sum_{j\in N_i}
\sum_{k\in N_i\setminus\{j\}}
(\widetilde b_{ij}+\widetilde b_{ik})^2\\
&\le d_+\sum_{i\in V}2d_+
=2nd_+^2.
\end{aligned}
\]
Finally, expanding the square and exchanging \(j\) and \(k\) in one of the resulting sums gives
\[
\sum_{j\in V}\sum_{k\in V}
\left(b^{\mathrm{pair}}_{jk}+b^{\mathrm{pair}}_{kj}\right)^2
=
2\sum_{j\in V}\sum_{k\in V}
b^{\mathrm{pair}}_{jk}
\left(b^{\mathrm{pair}}_{jk}+b^{\mathrm{pair}}_{kj}\right).
\]
It follows that
\[
\sum_{j\in V}\sum_{k\in V}
b^{\mathrm{pair}}_{jk}
\left(b^{\mathrm{pair}}_{jk}+b^{\mathrm{pair}}_{kj}\right)
\le nd_+^2.
\]
% lean: CausalSmith.Experimentation.ThinnedgraphAdditiveRiskFrontier.oracle_quadratic_energy_le

\item Combining the exact quadratic second moment with this energy bound gives
\[
\mathbb E_Z[T_{G,\mathrm{quad}}^2]
\le\frac{nd_+^2}{n^2}
=\frac{d_+^2}{n}.
\]
Together with the vanishing cross moment and the linear bound, this proves
\[
\operatorname{Var}_Z(T_G)
\le\frac{4d_+^2}{n}+\frac{d_+^2}{n}
=\frac{5(d+1)^2}{n}.
\]
The unbiasedness and variance identity established above complete the proof.
% lean: CausalSmith.Experimentation.ThinnedgraphAdditiveRiskFrontier.oracle_variance_le
\end{enumerate}
% lean: CausalSmith.Experimentation.ThinnedgraphAdditiveRiskFrontier.score_audit
\end{proof}

The unbiasedness statement concerns expectation over both parts of the experimental design for each fixed schedule. In the variance identity, the oracle term measures assignment variation, and the additional term measures the variation introduced by independently retaining true arrows. Each recipient's contribution to this additional term is weighted by its true in-degree and its assignment-averaged squared outcome. Thus the identity distinguishes the variance due to recording edges from the variance present when the graph is supplied.

The coefficient-mass restriction in \cref{obj:ass:coefficient-mass} controls both components of the upper-risk calculation. It bounds every potential outcome in absolute value by one and places the population contrast in \([-1,1]\). Together with the in-degree bound, it bounds the audit contribution by \(4d(1-q)/(nq)\). The common in-degree and out-degree restrictions enter the oracle bound in \cref{obj:lem:score-audit}, which is uniform over the schedule class. These two contributions are the terms in the upper envelope of \cref{obj:def:upper-envelope}.

The observable guarantee in \cref{obj:thm:observable-upper} uses the total rule of \cref{obj:def:upper-rule}. At positive retention and an upper envelope below one, this rule projects the audit score onto \([-1,1]\); projection preserves its squared-error upper bound because the causal target lies in that interval. For the remaining parameter values, the zero branch has squared loss at most one. Measurability follows from the score's finite sums and products of observed graph, assignment, and outcome coordinates, together with the continuous clipping map. The branch choice depends on the known design parameters. Consequently, the rule belongs to the estimator class in \cref{obj:def:risk}, including at zero retention.

\subsection{The supplied-graph lower comparison}

The oracle variance bound gives an attainable precision benchmark when the true graph is supplied. A complementary minimax statement establishes the degree-dependent lower scale for both observation experiments. The complete-record testing arguments in \cref{obj:lem:full-record-two-prior-risk} provide the risk-reduction framework for this comparison.

\begin{lemmav}{lem:supplied-block-record-lower}[Complete-record minimax lower bounds]
Let \(n\) be the population size, \(d\) the common off-diagonal in-degree and out-degree bound, and \(q\) the known edge-retention probability. Suppose that:
\begin{itemize}
\item \textbf{(Parameter range.)} \(n,d\) are integers with \(n\ge4\) and \(1\le d\le n-1\), and \(q\in[0,1]\).
\item \textbf{(Assignment law.)} The joint law of the treatment assignment and audit marks satisfies \cref{obj:ass:assignment}.
\item \textbf{(Audit law.)} The same joint law satisfies \cref{obj:ass:audit} at retention probability \(q\).
\item \textbf{(Design independence.)} The same joint law satisfies \cref{obj:ass:design-independence}.
\end{itemize}
Then the complete-record minimax squared risk \(R_{n,d}(q)\) of \cref{obj:def:risk} and the supplied-graph minimax squared risk \(R^G_{n,d}(q)\) of \cref{obj:def:supplied-risk} satisfy
\[
R_{n,d}(q)\ge
\frac{1}{160000\pi^2}
\min\left\{1,\frac{(d+1)^2}{n}\right\},
\qquad
R^G_{n,d}(q)\ge
\frac{1}{160000\pi^2}
\min\left\{1,\frac{(d+1)^2}{n}\right\}.
\]
These bounds apply to every measurable randomized estimator in the respective definitions, with the true off-diagonal graph supplied in the second experiment.
\end{lemmav}

\begin{proof}[Proof of \cref{obj:lem:supplied-block-record-lower}]
\begin{enumerate}
\item The design assumptions identify the joint experimental law. Relabel the population as
\[
V=\{1,\ldots,n\},
\qquad
D=\operatorname{Law}(Z,W).
\]
The assignment marginal is a probability measure, so \(D\) is a probability measure. Independence and the two specified marginals give
\[
D
=
\operatorname{Law}(Z)\otimes\operatorname{Law}(W)
=
\left(\bigotimes_{j\in V}\operatorname{Bernoulli}(1/2)\right)
\otimes
\left(\bigotimes_{\substack{i,j\in V\\i\ne j}}
\operatorname{Bernoulli}(q)\right),
\]
by \cref{obj:ass:assignment,obj:ass:audit,obj:ass:design-independence}.
We work with this product design, first for the complete-record risk and then for the supplied-graph risk.
% lean: CausalSmith.Experimentation.ThinnedgraphAdditiveRiskFrontier.design_eq_thinnedDesign

\item For the complete-record case, define the block size, block count, covered fraction, and amplitude by
\[
k_0=d+1,
\qquad
r_0=\left\lfloor\frac{n}{k_0}\right\rfloor,
\qquad
\rho=\frac{r_0k_0}{n},
\qquad
h=\frac{1}{100\pi}
\sqrt{\min\left\{1,\frac{k_0^2}{n}\right\}}.
\]
Here \(2\le k_0\le n\), so \(r_0\ge1\) and every denominator is positive. The quantity under the square root is positive and at most one. Since \(\pi\ge2\),
\[
0<h\le\frac{1}{100\pi}\le\frac14,
\qquad
h^2=
\frac{1}{10000\pi^2}
\min\left\{1,\frac{k_0^2}{n}\right\}.
\]
To check coverage, define the integer remainder by
\[
\operatorname{rem}_0=n-r_0k_0.
\]
Euclidean division and \(r_0\ge1\) imply
\[
0\le\operatorname{rem}_0<k_0\le r_0k_0,
\qquad
n=r_0k_0+\operatorname{rem}_0\le2r_0k_0.
\]
Consequently,
\[
\frac12\le\rho\le1,
\qquad
\Delta:=\rho h>0.
\]
% lean: CausalSmith.Experimentation.ThinnedgraphAdditiveRiskFrontier.completeBlockLowerAmplitude_properties
% lean: CausalSmith.Experimentation.ThinnedgraphAdditiveRiskFrontier.completeBlock_coverage_ge_half

\item Construct two priors on fixed schedules. Define the labeled blocks and their common graph by
\[
\mathcal B_v=\{(v-1)k_0+1,\ldots,vk_0\},
\quad v=1,\ldots,r_0,
\qquad
G_{n,d}
=
\bigcup_{v=1}^{r_0}
\{(j,i):i,j\in\mathcal B_v,\ j\ne i\}.
\]
Each block has \(k_0\) members. Both off-diagonal degrees are \(k_0-1=d\) on the blocks and zero on the remaining units.

Use the baseline density of \cref{obj:lem:baseline-translation-affinity},
\[
f(w)=
\begin{cases}
4\cos^2(2\pi w),& |w|\le1/4,\\
0,& |w|>1/4.
\end{cases}
\]
It is nonnegative and integrable, and
\[
\int_{\mathbb R}f(w)\,dw
=
\int_{-1/4}^{1/4}4\cos^2(2\pi w)\,dw
=
\left[2w+\frac{\sin(4\pi w)}{2\pi}\right]_{-1/4}^{1/4}
=1.
\]
Define the baseline-vector law by
\[
U=(U_1,\ldots,U_{r_0}),
\qquad
\operatorname{Law}(U)=\bigotimes_{v=1}^{r_0}f(w)\,dw.
\]
This is a probability law, and \(|U_v|\le1/4\) for every \(v\) almost surely.

For \(\sigma\in\{-1,1\}\) and \(U\in\mathbb R^{r_0}\), let
\(\theta^{\mathrm{cb}}_{\sigma,h}(U)\) be the schedule with graph \(G_{n,d}\) and coefficients
\[
a_i=
\begin{cases}
U_v-\sigma h/2,&i\in\mathcal B_v,\\
0,&i>r_0k_0,
\end{cases}
\qquad
t_i=
\begin{cases}
\sigma h/k_0,&i\le r_0k_0,\\
0,&i>r_0k_0,
\end{cases}
\qquad
b_{ij}=\frac{\sigma h}{k_0}
\quad\bigl((j,i)\in G_{n,d}\bigr).
\]
For a supported baseline draw and \(i\in\mathcal B_v\),
\[
|a_i|+|t_i|+\sum_{j:(j,i)\in G_{n,d}}|b_{ij}|
=
|U_v-\sigma h/2|+\frac{h}{k_0}
+d\frac{h}{k_0}
\le\frac14+\frac h2+h
\le\frac58\le1.
\]
The corresponding expression is zero for a padding unit. Thus these schedules belong almost surely to the class in \cref{obj:def:model}.

For a binary assignment \(z\in\{0,1\}^{V}\), define
\[
X_j(z)=2z_j-1,
\qquad j\in V.
\]
The additive outcome formula gives
\[
Y_i(z)=
\begin{cases}
\displaystyle
U_v-\frac{\sigma h}{2}
+\frac{\sigma h}{k_0}\sum_{j\in\mathcal B_v}z_j
=
U_v+\frac{\sigma h}{2k_0}
\sum_{j\in\mathcal B_v}X_j(z),
&i\in\mathcal B_v,\\[6pt]
0,&i>r_0k_0.
\end{cases}
\]
In particular,
\[
Y_i(\mathbf1)-Y_i(\mathbf0)=
\begin{cases}
\sigma h,&i\le r_0k_0,\\
0,&i>r_0k_0,
\end{cases}
\qquad
\tau\bigl(\theta^{\mathrm{cb}}_{\sigma,h}(U)\bigr)
=
\sigma\frac{r_0k_0h}{n}
=\sigma\Delta.
\]
Define the priors by
\[
\Pi^{\mathrm{cb}}_{\sigma,h}
=
\operatorname{Law}\bigl(\theta^{\mathrm{cb}}_{\sigma,h}(U)\bigr).
\]
Under each prior the schedule is drawn before, and independently of, the experiment. The recording map is jointly measurable: its graph and assignment coordinates are finite, and its outcome coordinates are affine functions of \(U\) for each design draw.
% lean: CausalSmith.Experimentation.ThinnedgraphAdditiveRiskFrontier.completeBlock_support
% lean: CausalSmith.Experimentation.ThinnedgraphAdditiveRiskFrontier.completeBlock_record_measurable

\item Define the complete-record mixture laws by
\[
\mathsf Q_+
=
\int P_{\theta,q}\,\Pi^{\mathrm{cb}}_{1,h}(d\theta),
\qquad
\mathsf Q_-
=
\int P_{\theta,q}\,\Pi^{\mathrm{cb}}_{-1,h}(d\theta).
\]
The support, constant-target, and measurability requirements of
\cref{obj:lem:full-record-two-prior-risk} have just been verified. It therefore supplies
\[
R_{n,d}(q)
\ge
\frac{\Delta^2}{2}
\left(1-\operatorname{TV}(\mathsf Q_+,\mathsf Q_-)\right).
\]
% lean: CausalSmith.Experimentation.ThinnedgraphAdditiveRiskFrontier.full_record_two_prior_risk

\item To bound this testing distance, retain the entire design draw \((z,w)\) and the vector \(y\in\mathbb R^{r_0}\) of distinct block responses. Define the shifts and response densities by
\[
\delta_{\sigma,v}(z)
=
\frac{\sigma h}{2k_0}
\sum_{j\in\mathcal B_v}X_j(z),
\qquad
\sigma\in\{-1,1\},\quad v=1,\ldots,r_0,
\]
\[
\lambda^{\mathrm{cb}}_{\sigma,h}(y\mid z,w)
=
\prod_{v=1}^{r_0}f\bigl(y_v-\delta_{\sigma,v}(z)\bigr).
\]
For every design draw these densities are nonnegative and normalized:
\[
\int_{\mathbb R^{r_0}}
\lambda^{\mathrm{cb}}_{\sigma,h}(y\mid z,w)\,dy
=
\prod_{v=1}^{r_0}
\int_{\mathbb R}f\bigl(y_v-\delta_{\sigma,v}(z)\bigr)\,dy_v
=1.
\]
Indeed, translating \(U_v\) by \(\delta_{\sigma,v}(z)\) gives the density
\(f(y_v-\delta_{\sigma,v}(z))\); taking the product gives the displayed vector density.

Define the joint laws by
\[
\mathsf J_\pm(dz,dw,dy)
=
D(dz,dw)\,
\lambda^{\mathrm{cb}}_{\pm1,h}(y\mid z,w)\,dy.
\]
These are precisely the laws of the full design draw together with
\((U_v+\delta_{\pm1,v}(Z))_{v=1}^{r_0}\).

For an audit array \(w\), define the retained graph by
\[
H(w)_{ij}
=
\mathbf1\{(j,i)\in G_{n,d}\}\,w_{ij},
\qquad i\ne j.
\]
The common recording map is
\[
\mathcal C_{\mathrm{rec}}((z,w),y)
=
\left(
H(w),\,z,\,
\left(
\begin{cases}
y_v,&i\in\mathcal B_v,\\
0,&i>r_0k_0
\end{cases}
\right)_{i\in V}
\right).
\]
The outcome identity in the preceding construction proves
\[
O
=
\mathcal C_{\mathrm{rec}}
\left((Z,W),(U_v+\delta_{\sigma,v}(Z))_{v=1}^{r_0}\right),
\qquad
\mathsf Q_\pm=(\mathcal C_{\mathrm{rec}})_\#\mathsf J_\pm.
\]
Here the subscript \(\#\) denotes the pushforward of a measure by the displayed map. The map is measurable, and its pushforward preserves the labeled retained graph and every treatment coordinate. Taking preimages of measurable events yields
\[
\operatorname{TV}(\mathsf Q_+,\mathsf Q_-)
\le
\operatorname{TV}(\mathsf J_+,\mathsf J_-).
\]
% lean: CausalSmith.Experimentation.ThinnedgraphAdditiveRiskFrontier.completeBlock_mixture_density_representation
% lean: CausalSmith.Experimentation.ThinnedgraphAdditiveRiskFrontier.completeBlock_mixture_tv_population_le

\item Define the conditional squared Hellinger energy by
\[
\mathcal E_{\mathrm{cb}}(z,w)
=
\int_{\mathbb R^{r_0}}
\left(
\sqrt{\lambda^{\mathrm{cb}}_{1,h}(y\mid z,w)}
-
\sqrt{\lambda^{\mathrm{cb}}_{-1,h}(y\mid z,w)}
\right)^2dy.
\]
The product-translation bound in
\cref{obj:lem:baseline-translation-affinity} gives
\[
\mathcal E_{\mathrm{cb}}(z,w)
\le
\sum_{v=1}^{r_0}
4\pi^2
\bigl(\delta_{1,v}(z)-\delta_{-1,v}(z)\bigr)^2.
\]

Write \(X_j=X_j(Z)\) for the random treatment signs. Under the assignment marginal of \(D\),
\[
\mathbb E_D[X_j]=0,
\qquad
\mathbb E_D[X_j^2]=1,
\qquad
\mathbb E_D[X_iX_j]=0\quad(i\ne j).
\]
Expanding the square therefore gives, for every block,
\[
\mathbb E_D
\left[\left(\sum_{j\in\mathcal B_v}X_j\right)^2\right]
=
\sum_{j\in\mathcal B_v}\mathbb E_D[X_j^2]
+
\sum_{\substack{i,j\in\mathcal B_v\\i\ne j}}
\mathbb E_D[X_iX_j]
=k_0.
\]
Since
\[
\delta_{1,v}(Z)-\delta_{-1,v}(Z)
=
\frac{h}{k_0}\sum_{j\in\mathcal B_v}X_j,
\]
it follows that
\[
\mathbb E_D
\left[
\bigl(\delta_{1,v}(Z)-\delta_{-1,v}(Z)\bigr)^2
\right]
=
\frac{h^2}{k_0}.
\]
Define the averaged energy by
\[
\mathcal E_{\mathrm{joint}}
=
\int\mathcal E_{\mathrm{cb}}(z,w)\,D(dz,dw).
\]
Integrating the conditional bound and using \(r_0k_0\le n\) gives
\[
\mathcal E_{\mathrm{joint}}
\le
4\pi^2h^2\frac{r_0}{k_0}
\le
4\pi^2h^2\frac{n}{k_0^2}.
\]

The common-marginal identity in
\cref{obj:lem:baseline-translation-affinity} identifies
\(\mathcal E_{\mathrm{joint}}\) as the unhalved squared Hellinger distance between \(\mathsf J_+\) and \(\mathsf J_-\).
For completeness, the sharper total-variation comparison used here follows directly from Cauchy--Schwarz. Define the common reference measure by
\[
\mu_{\mathrm{cb}}(dz,dw,dy)=D(dz,dw)\,dy.
\]
Normalization and the definition of the energy imply
\[
\int
\left(
\sqrt{\lambda^{\mathrm{cb}}_{1,h}}
+
\sqrt{\lambda^{\mathrm{cb}}_{-1,h}}
\right)^2d\mu_{\mathrm{cb}}
=
4-\mathcal E_{\mathrm{joint}}.
\]
Consequently,
\[
\begin{aligned}
\operatorname{TV}(\mathsf J_+,\mathsf J_-)
&=
\frac12\int
\left|\lambda^{\mathrm{cb}}_{1,h}
-\lambda^{\mathrm{cb}}_{-1,h}\right|d\mu_{\mathrm{cb}}\\
&\le
\frac12
\sqrt{\mathcal E_{\mathrm{joint}}}
\sqrt{4-\mathcal E_{\mathrm{joint}}}\\
&=
\sqrt{\mathcal E_{\mathrm{joint}}
\left(1-\mathcal E_{\mathrm{joint}}/4\right)}
\le
\sqrt{\mathcal E_{\mathrm{joint}}}.
\end{aligned}
\]
The densities in these integrals are evaluated at \((y\mid z,w)\). Together with the recording-map bound, this proves
\[
\operatorname{TV}(\mathsf Q_+,\mathsf Q_-)
\le
\sqrt{4\pi^2h^2\frac{n}{k_0^2}}.
\]
% lean: CausalSmith.Experimentation.ThinnedgraphAdditiveRiskFrontier.completeBlock_shift_separation_moment
% lean: CausalSmith.Experimentation.ThinnedgraphAdditiveRiskFrontier.completeBlock_response_energy_population_le
% lean: CausalSmith.Experimentation.ThinnedgraphAdditiveRiskFrontier.completeBlock_densityJoint_tv_le
% lean: CausalSmith.Experimentation.ThinnedgraphAdditiveRiskFrontier.completeBlock_mixture_tv_population_le

\item The amplitude calibration now gives
\[
4\pi^2h^2\frac{n}{k_0^2}
=
\min\left\{1,\frac{k_0^2}{n}\right\}
\frac{n}{2500k_0^2}
\le\frac1{2500},
\]
and hence
\[
\operatorname{TV}(\mathsf Q_+,\mathsf Q_-)\le\frac1{50}.
\]
Coverage and positivity give the separation comparisons
\[
\frac h2\le\rho h=\Delta,
\qquad
\frac{h^2}{4}\le\Delta^2.
\]
The exact squared-amplitude identity also gives
\[
\frac{1}{160000\pi^2}
\min\left\{1,\frac{(d+1)^2}{n}\right\}
=\frac{h^2}{16}.
\]
Substituting these comparisons into the complete-record two-prior bound yields
\[
\begin{aligned}
R_{n,d}(q)
&\ge
\frac{\Delta^2}{2}
\left(1-\operatorname{TV}(\mathsf Q_+,\mathsf Q_-)\right)\\
&\ge
\frac{h^2}{8}\left(1-\frac1{50}\right)
\ge\frac{h^2}{16}\\
&=
\frac{1}{160000\pi^2}
\min\left\{1,\frac{(d+1)^2}{n}\right\}.
\end{aligned}
\]
% lean: CausalSmith.Experimentation.ThinnedgraphAdditiveRiskFrontier.completeBlock_testing_distance
% lean: CausalSmith.Experimentation.ThinnedgraphAdditiveRiskFrontier.completeBlock_risk_lower_assembly

\item For the supplied-graph case, use the same amplitude and priors. Their support and target properties are
\[
\Pi^{\mathrm{cb}}_{\sigma,h}(\mathcal M_{n,d})=1,
\qquad
\tau\bigl(\theta^{\mathrm{cb}}_{\sigma,h}(U)\bigr)=\sigma\Delta,
\qquad
G\bigl(\theta^{\mathrm{cb}}_{\sigma,h}(U)\bigr)=G_{n,d}.
\]
Fix any measurable randomized rule \(T^G\) admitted by
\cref{obj:def:supplied-risk}. Define its rule with the graph fixed, and its supplied-graph worst-case risk, by
\[
T_{\mathrm{fix}}(o,\xi)=T^G(o,G_{n,d},\xi),
\qquad \xi\in[0,1],
\]
\[
\mathcal W^G(T^G)
=
\sup_{\theta\in\mathcal M_{n,d}}
\mathbb E_{P_{\theta,q}\otimes\nu_\xi}
\left[
\bigl(T^G(O,G(\theta),\xi)-\tau(\theta)\bigr)^2
\right],
\qquad
\nu_\xi=\operatorname{Uniform}[0,1].
\]
Inserting the constant finite graph is measurable. For every schedule in either prior,
\[
\bigl(T_{\mathrm{fix}}(O,\xi)-\tau(\theta)\bigr)^2
=
\bigl(T^G(O,G(\theta),\xi)-\tau(\theta)\bigr)^2.
\]
Integrating over the supported schedules therefore gives
\[
\int
\bigl(T_{\mathrm{fix}}(o,\xi)-\Delta\bigr)^2
\,d(\mathsf Q_+\otimes\nu_\xi)(o,\xi)
\le\mathcal W^G(T^G),
\]
\[
\int
\bigl(T_{\mathrm{fix}}(o,\xi)+\Delta\bigr)^2
\,d(\mathsf Q_-\otimes\nu_\xi)(o,\xi)
\le\mathcal W^G(T^G).
\]
These are nonnegative integrals, so the identities and inequalities also hold when the risk is infinite.
% lean: CausalSmith.Experimentation.ThinnedgraphAdditiveRiskFrontier.hardcodeSuppliedEstimator_sqLoss
% lean: CausalSmith.Experimentation.ThinnedgraphAdditiveRiskFrontier.mixture_seeded_risk_le_supplied_worst

\item Adding the common independent seed preserves total variation:
\[
\operatorname{TV}
(\mathsf Q_+\otimes\nu_\xi,\mathsf Q_-\otimes\nu_\xi)
=
\operatorname{TV}(\mathsf Q_+,\mathsf Q_-).
\]
Define the measurable testing event by
\[
\mathcal E_{\mathrm{err}}
=
\{(o,\xi):T_{\mathrm{fix}}(o,\xi)<0\}.
\]
The definition of total variation gives
\[
(\mathsf Q_+\otimes\nu_\xi)(\mathcal E_{\mathrm{err}})
+
(\mathsf Q_-\otimes\nu_\xi)(\mathcal E_{\mathrm{err}}^c)
\ge
1-\operatorname{TV}(\mathsf Q_+,\mathsf Q_-).
\]
Because \(\Delta>0\), the pointwise loss comparisons are
\[
\bigl(T_{\mathrm{fix}}(o,\xi)-\Delta\bigr)^2
\ge\Delta^2\mathbf1_{\mathcal E_{\mathrm{err}}}(o,\xi),
\]
\[
\bigl(T_{\mathrm{fix}}(o,\xi)+\Delta\bigr)^2
\ge\Delta^2\mathbf1_{\mathcal E_{\mathrm{err}}^c}(o,\xi).
\]
Integrating and using the two worst-case-risk comparisons proves
\[
2\mathcal W^G(T^G)
\ge
\Delta^2\left(1-\operatorname{TV}(\mathsf Q_+,\mathsf Q_-)\right).
\]
Division by the positive finite constant \(2\) is valid in the extended nonnegative reals. Infimizing over \(T^G\) thus yields
\[
R^G_{n,d}(q)
\ge
\frac{\Delta^2}{2}
\left(1-\operatorname{TV}(\mathsf Q_+,\mathsf Q_-)\right).
\]
The same testing-distance and separation bounds give
\[
\begin{aligned}
R^G_{n,d}(q)
&\ge
\frac{h^2}{8}\left(1-\frac1{50}\right)
\ge\frac{h^2}{16}\\
&=
\frac{1}{160000\pi^2}
\min\left\{1,\frac{(d+1)^2}{n}\right\}.
\end{aligned}
\]
The construction has at least one full block throughout the stated parameter range, and the density comparison integrates over the full design law. These arguments apply for every \(q\in[0,1]\), including both endpoints, and for every measurable randomized estimator in the two risk definitions.
% lean: CausalSmith.Experimentation.ThinnedgraphAdditiveRiskFrontier.two_prior_squared_testing
% lean: CausalSmith.Experimentation.ThinnedgraphAdditiveRiskFrontier.completeBlock_supplied_two_prior
% lean: CausalSmith.Experimentation.ThinnedgraphAdditiveRiskFrontier.completeBlock_supplied_risk_lower_assembly
\end{enumerate}
\end{proof}

Both lower bounds are uniform in the retention probability and cover arbitrary measurable use of the available record and estimator randomization. In particular, the supplied-graph bound establishes the degree-dependent scale even when every true arrow is available to the estimator. Combined with the oracle bound in \cref{obj:lem:score-audit} and clipping to the target interval, it characterizes supplied-graph minimax precision up to universal constants. The audit-score variance identity then quantifies the additional variance of the observable score relative to that benchmark.

\subsection{Published-model correspondence}

A comparison with the published SNIPE model requires matching its graph and coefficient conventions. In the interaction-order-one specification of \citet[Sections 3.1--3.2, Assumptions 1--2, equations (3.2)--(3.3)]{CortezRodriguezEichhornYu2023SNIPE}, outcomes are polynomials with a constant term and singleton-coordinate terms indexed by each recipient's in-neighborhood. Those neighborhoods may contain the recipient itself, and the coefficient-mass envelope includes the constant coefficient. The following correspondence therefore adjoins each unit's own coordinate to the off-diagonal graph used here.

\begin{propositionv}{lem:published-model}[Published-model correspondence and risk transfer]*
Let \(V\) be the fixed finite population. Use the published SNIPE model conventions of \citep[Sections 3.1--3.2, Assumptions 1--2, equations (3.2)--(3.3)]{CortezRodriguezEichhornYu2023SNIPE}. For every additive schedule \(\theta=(G,a,t,b)\), define its augmented graph \(\widetilde G\) and published polynomial coefficients by
\[
\widetilde G=G\cup\{(i,i):i\in V\},\qquad
c^{\mathrm{pub}}_{i,\varnothing}=a_i,\qquad
c^{\mathrm{pub}}_{i,\{i\}}=t_i,\qquad
c^{\mathrm{pub}}_{i,\{j\}}=b_{ij}\quad ((j,i)\in G),
\]
with all remaining polynomial coefficients zero.

For every nonnegative integer \(d\), if \(\theta\in\mathcal M_{n,d}\), the schedule class in \cref{obj:def:model}, its augmentation has every diagonal arrow and belongs to the published specialization with interaction order \(\beta=1\), coefficient-mass envelope
\[
Y_{\max}
=\max_{i\in V}
\left(
|c^{\mathrm{pub}}_{i,\varnothing}|
+\sum_{j:(j,i)\in\widetilde G}
|c^{\mathrm{pub}}_{i,\{j\}}|
\right)
\le 1,
\]
and total in-degree and out-degree at most \(d+1\). Conversely, every published schedule satisfying these order, envelope, and degree restrictions and containing every diagonal arrow yields a member of \(\mathcal M_{n,d}\) by deleting the diagonal arrows and reading its empty-set coefficient as \(a_i\), its own singleton coefficient as \(t_i\), and its singleton coefficient for source \(j\) as \(b_{ij}\). Augmenting this stripped schedule recovers the original graph and every polynomial coefficient indexed by a subset of a recipient's in-neighborhood. For every additive schedule, stripping its augmentation recovers its graph, baseline vector, own-treatment vector, and spillover coefficients on all graph arrows.

For every additive schedule, every unit \(i\), and every binary assignment \(z\), augmentation preserves \(Y_i(z)\), the additive potential outcome, and preserves \(\tau(\theta)\), the population all-treated-versus-all-control contrast. For every assignment and audit-mark array, adjoining the known diagonal arrows to \(H\), the retained graph, gives the published record through the deterministic transformation
\[
O=(H,Z,Y(Z))
\longmapsto
\bigl(H\cup\{(i,i):i\in V\},\,Z,\,Y(Z)\bigr),
\]
where \(O\) is the complete observed record, \(Z\) is the assignment vector, and \(Y(Z)\) is the realized-outcome vector.

For every joint assignment-and-audit law satisfying \cref{obj:ass:assignment,obj:ass:audit,obj:ass:design-independence}, with \(q\) the edge-retention probability, and every nonnegative integer \(d\), consider any class of published schedules containing the augmentation of every member of \(\mathcal M_{n,d}\). Under the same assignment law and off-diagonal audit channel, \(R_{n,d}(q)\), the minimax squared risk defined in \cref{obj:def:risk}, is at most the minimax squared risk over that published class.
\end{propositionv}
% CAUSALSMITH-CITED-SCOPE-BEGIN lem:published-model
\verificationfootnotetext{\textbf{Formalization scope.} This result uses the cited conclusion \citep{CortezRodriguezEichhornYu2023SNIPE} (Sections 3.1--3.2, Assumptions 1--2 and equations (3.2)--(3.3), arXiv version 4.). The cited source proof is not formalized here: Lean verifies its use and all remaining steps. If that cited conclusion is false, the portions of this result that depend on it are not certified.}
% CAUSALSMITH-CITED-SCOPE-END lem:published-model

The augmented graph \(\widetilde G\) in \cref{obj:lem:published-model} has one additional incoming and outgoing arrow per unit. Accordingly, the published total-degree index is \(d+1\) when the off-diagonal degree index is \(d\). The own-treatment coefficient occupies the recipient's singleton coordinate, and the published coefficient-mass envelope agrees with the row restriction in \cref{obj:ass:coefficient-mass}. The cited model conventions permit zero coefficients on graph arrows, so the correspondence retains the declared graph together with its coefficients.

The risk-transfer clause uses both schedule inclusion and a common observation channel. Under the same assignment law and off-diagonal audit mechanism, adjoining the known diagonal arrows transforms the observed record while preserving realized outcomes and the causal target. The resulting risk ordering transfers a lower bound from the additive class to any published class containing all of its augmentations. Comparisons with a supplied true graph use the separate experiment in \cref{obj:def:supplied-risk}; the published-model correspondence specifies the channel under which its own class-inclusion comparison applies.

\subsection{Supported boundary cases}

The observable upper guarantee and supplied-graph lower comparison fix the full-retention scale. The complete-record testing bounds in \cref{obj:thm:block-testing-scale,obj:lem:full-record-two-prior-risk} supply the complementary retention-dependent lower scale. The following statement consolidates these ingredients at zero retention, full retention, and degree one.

\begin{theoremv}{thm:supported-boundaries}[Minimax bounds at supported boundaries]
There exist universal constants \(c,C_0>0\) such that the following holds for every population size \(n\), degree bound \(d\), known retention probability \(q\), and joint law of the treatment-assignment vector \(Z\) and the off-diagonal audit-mark array \(W\) satisfying:
\begin{itemize}
\item \textbf{(Population.)} \(n\) is an integer with \(n\ge4\).
\item \textbf{(Degree.)} \(d\) is an integer with \(1\le d\le n-1\).
\item \textbf{(Retention.)} \(q\in[0,1]\).
\item \textbf{(Assignment.)} The joint law satisfies \cref{obj:ass:assignment}.
\item \textbf{(Audit.)} The joint law satisfies \cref{obj:ass:audit} at retention probability \(q\).
\item \textbf{(Independence.)} The joint law satisfies \cref{obj:ass:design-independence}.
\end{itemize}
For each such law, the minimax squared risk \(R_{n,d}(q)\) of \cref{obj:def:risk}, with the infimum taken over all measurable randomized estimators of the complete observed record \(O\), satisfies
\[
R_{n,d}(q)\ge
c\max\left\{
\min\left\{1,\frac{(d+1)^2}{n}\right\},
\frac{1}{1+nq}
\right\}.
\]
The supported boundary cases satisfy
\[
\begin{aligned}
c&\le R_{n,d}(0)\le1,\\
c\min\left\{1,\frac{(d+1)^2}{n}\right\}
&\le R_{n,d}(1)
\le C_0\min\left\{1,\frac{(d+1)^2}{n}\right\},\\
\frac{c}{1+nq}
&\le R_{n,1}(q)\le\frac{C_0}{1+nq}.
\end{aligned}
\]
\end{theoremv}

The two terms in the lower bound express distinct sources of difficulty: the degree-dependent precision scale already present with a supplied graph, and the retention-dependent scale associated with partial observation. At zero retention, the complete-record minimax risk remains bounded below by a universal positive constant. At full retention, the audit contribution in \cref{obj:lem:score-audit} vanishes, and the upper and lower bounds agree up to universal constants at the degree-dependent scale. For degree one, the matching reciprocal scale \(1/(1+nq)\) describes the entire retention range, linking constant risk at sparse collection to order \(1/n\) risk at full retention.


\section{Complete-record block experiments}

The lower-bound analysis uses pairs of priors over fixed schedules whose induced experiments retain the complete observed record. The construction separates labeled sources from fixed groups of recipients, allowing the response likelihood to account jointly for retained arrows, treatment assignments, and repeated recipient outcomes. Its auxiliary parameters are the block count \leanref{sym:B}{\ensuremath{B}} and the source count \leanref{sym:m}{\ensuremath{m=Bd}}. The following definition fixes the recipient groups and the population capacity required by the construction.

\begin{definitionv}{synth_4}[Fixed recipient blocks]
For integers \(n\ge4\) and \(B,d\ge1\) satisfying \(2Bd\le n\), set \(m=Bd\) and \(V=\{1,\ldots,n\}\). The source units are \(\{1,\ldots,m\}\). Define the fixed recipient blocks by
\[
C_\ell=\{m+(\ell-1)d+1,\ldots,m+\ell d\},
\qquad \ell\in\{1,\ldots,B\}.
\]
Each \(C_\ell\) contains \(d\) units, and these pairwise disjoint blocks partition \(\{m+1,\ldots,2m\}\). The padding units are \(\{i\in V:i>2m\}\). Thus the sources, recipient blocks, and padding units partition \(V\). The recipient blocks remain fixed as the ordered source partition \(S=(S_1,\ldots,S_B)\) varies; \(S_\ell\) is the source group assigned to \(C_\ell\) in the block-prior construction.
\end{definitionv}

The recipient block \leanref{sym:Cblocks}{\ensuremath{C_\ell}}, indexed by \leanref{sym:ell}{\ensuremath{\ell}}, remains fixed while the ordered source partition \leanref{sym:S}{\ensuremath{S}} varies. This distinction makes the repeated-response groups observable from recipient labels while leaving source membership partially revealed by the audit. We next specify the partition domain and the baseline density used to randomize fixed schedules.



The baseline vector \leanref{sym:U}{\ensuremath{U}} has density determined by \leanref{sym:f}{\ensuremath{f}}, with real density argument \leanref{sym:w}{\ensuremath{w}}, as specified in \cref{obj:def:block-prior}. The sign \leanref{sym:sigma}{\ensuremath{\sigma}} selects the direction of the spillover effect, and the amplitude \leanref{sym:h}{\ensuremath{h}} controls its magnitude. Each baseline is shared by a recipient block and is drawn before the experimental assignment. Thus the baseline randomization defines a prior over schedules, with each realized schedule held fixed during its experiment.

To describe the observable information entering the likelihood, consider a retained graph compatible with the block layout. A retained outgoing arrow identifies its source's recipient block, and the assignment vector records the treatment of every labeled source. The following definition collects the resulting capacities, sign sums, and count-constrained likelihood notation.



The remaining capacities \leanref{sym:urow}{\ensuremath{u_\ell}} and revealed sign sums \leanref{sym:A}{\ensuremath{A_\ell}} in \cref{obj:lem:block-law} are functions of the labeled retained graph and assignments. The same record determines the total remaining capacity \leanref{sym:M}{\ensuremath{M}} and the global hidden treated-source count \leanref{sym:K}{\ensuremath{K}}. The block-response vector \leanref{sym:y}{\ensuremath{y}} records one outcome per fixed recipient group, preserving the repeated observations through the known groups. In the likelihood expression \leanref{sym:density}{\ensuremath{\lambda_{\sigma,h}(y\mid H,Z)}} defined in \cref{obj:lem:block-law}, the indeterminate \leanref{sym:x}{\ensuremath{x}} tracks treated counts, while \leanref{sym:k}{\ensuremath{k}} indexes a block's possible hidden treated count.

These definitions organize the information used by the complete-record experiment. The prior now specifies how the source partition and baselines jointly generate its fixed graph and coefficients.

\begin{definitionv}{def:block-prior}[Block schedule prior]
The block prior \(\Pi_{\sigma,h}\), indexed by the sign \(\sigma\in\{-1,1\}\) and amplitude \(h\in\mathbb R\), is the distribution over fixed schedules \(\theta=(G,a,t,b)\) induced by the following construction. With \(B\) blocks and \(m=Bd\) sources, draw an ordered partition \(S=(S_1,\ldots,S_B)\) uniformly over partitions of the sources satisfying \(|S_\ell|=d\). Independently of \(S\), draw independent block baselines \(U_1,\ldots,U_B\) with density
\[
f(w)=
\begin{cases}
4\cos^2(2\pi w),& |w|\le 1/4,\\
0,& |w|>1/4.
\end{cases}
\]
Using the recipient blocks \(C_\ell\) of the block construction, set
\[
G=\bigcup_{\ell=1}^{B}\{(j,i):j\in S_\ell,\ i\in C_\ell\},
\qquad
a_i=\sum_{\ell=1}^{B}\mathbf1\{i\in C_\ell\}
       \left(U_\ell-\frac{\sigma h}{2}\right),
\qquad
t_i=0,
\qquad
b_{ij}=\frac{\sigma h}{d}.
\]
Spillover coefficients are read on graph arrows. Source and padding rows have zero outcomes. The baselines are drawn before treatment randomization and are fixed components of each schedule.

The corresponding complete-record mixture law \(P_{\sigma,h}\) is obtained by drawing the partition and baselines, fixing the resulting schedule, and then independently drawing treatment assignments and audit marks from the canonical product design with retention parameter \(q\). For every measurable set \(A\) of complete observed records \(O=(H,Z,Y(Z))\),
\[
P_{\sigma,h}(A)
=
\int
\Pr\!\left\{(H,Z,Y(Z))\in A\mid\theta\right\}
\,\Pi_{\sigma,h}(d\theta),
\]
where the conditional probability uses that assignment-audit design.
\end{definitionv}

The schedule prior \leanref{sym:Pi}{\ensuremath{\Pi_{\sigma,h}}} and complete-record mixture law \leanref{sym:Psign}{\ensuremath{P_{\sigma,h}}} in \cref{obj:def:block-prior} distinguish the two sources of randomization. Partition and baseline draws select a schedule; treatment and audit draws then produce an observation conditional on that schedule. Within a recipient block, the baseline centering gives the realized response
\[
Y_i(Z)
=
U_\ell+\frac{\sigma h}{2d}\sum_{j\in S_\ell}(2Z_j-1),
\qquad i\in C_\ell.
\]
Consequently, the same block response appears at each of its \(d\) recipients. The testing pair uses the following explicit amplitude to relate this construction to the estimation scale.

\begin{definitionv}{def:testing-handle}[Testing amplitude $h_\circ(B,d,q)$]
The testing amplitude is
\[
h_\circ(B,d,q)=\frac{1}{100\pi}
\begin{cases}
\min\left\{1,\dfrac{d}{\sqrt{mp}}\right\},&p>0,\\
1,&p=0,
\end{cases}
\]
where the source count \(m\) and association-revelation probability \(p\) are
\[
m=Bd,
\qquad
p=1-(1-q)^d.
\]
The parameters satisfy
\begin{itemize}
\item \textbf{(Block count.)} \(B\) is an integer with \(B\ge1\).
\item \textbf{(Degree.)} \(d\) is an integer with \(d\ge1\).
\item \textbf{(Population capacity.)} \(2m\le n\).
\end{itemize}
The testing handle consists of the opposite-sign block priors and their complete original-record mixture laws:
\[
\bigl(\Pi_{+1,h_\circ(B,d,q)},\Pi_{-1,h_\circ(B,d,q)}\bigr),
\qquad
\bigl(P_{+1,h_\circ(B,d,q)},P_{-1,h_\circ(B,d,q)}\bigr).
\]
The record includes detailed retained arrows and source labels in the retained graph \(H\), every coordinate of the assignment vector \(Z\), all noiseless realized outcomes \(Y(Z)\), repeated-outcome recipient grouping, remaining block capacities \(u_\ell\), and the observed global hidden treated-source count \(K\). Its conditional response density \(\lambda_{\sigma,h}\) uses coefficient extraction constrained by \(K\), including the convention for zero total remaining capacity \(M=0\).

Support, causal target separation, the uniform nonlocal testing bound at \(h_\circ\), and the all-degree Walsh contraction are proved properties of this construction. Together they yield the risk scale \(F\), its equivalent positive-retention elbow expression, the clipped-or-zero rule \(T^\star_{n,d,q}\), and the two consistency criteria
\[
\frac{d_n^2}{n}\longrightarrow0,
\qquad
\frac{d_n}{nq_n}\longrightarrow0,
\]
for admissible population sequences, with the second ratio assigned \(+\infty\) at zero retention.
\end{definitionv}

The amplitude \(h_\circ(B,d,q)\) in \cref{obj:def:testing-handle} lies within the support range \(0\le h\le1/4\). At positive association-revelation probability it combines a bounded amplitude with the scale \(d/\sqrt{mp}\); at zero revelation it uses the specified constant amplitude. The support and exact likelihood are supplied by the next result, while the contraction and testing bounds appear in \cref{obj:lem:hidden-allocation-contraction,obj:thm:block-testing-scale,obj:thm:nonlocal-testing-answer}.

\begin{lemmav}{lem:block-law}[Block mixture and reconstruction]
Let \(n,B,d\) be integers, let \(\sigma\in\{-1,1\}\) be the prior sign, let \(h\) be its amplitude, and let \(q\) be the edge-retention probability. Suppose:
\begin{itemize}
\item \textbf{(Population and blocks.)} \(n\ge4\), \(B\ge1\), \(d\ge1\), and \(2Bd\le n\).
\item \textbf{(Parameter ranges.)} \(0\le h\le1/4\) and \(0\le q\le1\).
\item \textbf{(Experimental design.)} The joint law of the assignment vector \(Z\) and audit-mark array \(W\) satisfies \cref{obj:ass:assignment,obj:ass:audit,obj:ass:design-independence}.
\end{itemize}

Set \(m=Bd\), the number of sources. Partition the \(m\) labeled sources uniformly into ordered blocks \(S_1,\ldots,S_B\) of size \(d\), independently of block baselines \(U_1,\ldots,U_B\), which are independent with density
\[
f(u)=
\begin{cases}
4\cos^2(2\pi u),& |u|\le1/4,\\
0,& |u|>1/4.
\end{cases}
\]
For the fixed recipient blocks \(C_1,\ldots,C_B\), each of size \(d\), include every arrow from \(S_\ell\) to \(C_\ell\). Set
\[
a_i=U_\ell-\frac{\sigma h}{2}\quad(i\in C_\ell),
\qquad t_i=0,
\qquad b_{ij}=\frac{\sigma h}{d}
\]
on these arrows, and give source and padding units zero outcomes. Let \(\Pi_{\sigma,h}\) be the resulting schedule law and \(P_{\sigma,h}\) its complete-record mixture law. Then, for \(\Pi_{\sigma,h}\)-almost every schedule \(\theta\),
\[
\theta\in\mathcal M_{n,d},
\qquad
\tau(\theta)=\frac{\sigma mh}{n},
\]
where \(\mathcal M_{n,d}\) is the schedule class in \cref{obj:def:model} and \(\tau(\theta)\) is the population all-treated-versus-all-control contrast.

Under the actual retained-graph marginal, the indicators that each labeled source has at least one retained outgoing arrow are independent Bernoulli variables with common probability
\[
p=1-(1-q)^d.
\]

Let \(H\) be the detailed labeled retained graph. For each block, define its remaining source capacity and revealed source-sign sum by
\[
u_\ell=d-\bigl|\{j\in S_\ell:j\text{ has a retained outgoing arrow}\}\bigr|,
\qquad
A_\ell=\sum_{\substack{j\in S_\ell\\j\text{ has a retained outgoing arrow}}}(2Z_j-1).
\]
Define the total remaining capacity and hidden treated-source count by
\[
M=\sum_{\ell=1}^B u_\ell,
\qquad
K=\sum_{\substack{j\text{ a source}\\j\text{ has no retained outgoing arrow}}}Z_j.
\]
For almost every \((H,Z)\) under its actual joint marginal, the conditional density of the distinct block-response vector \(y=(y_1,\ldots,y_B)\) is
\[
\lambda_{\sigma,h}(y\mid H,Z)
=
\frac{
[x^K]\displaystyle\prod_{\ell=1}^B
\left\{
\sum_{k=0}^{u_\ell}
\binom{u_\ell}{k}
f\!\left(y_\ell-\frac{\sigma h(A_\ell+2k-u_\ell)}{2d}\right)x^k
\right\}
}{
\binom{M}{K}
},
\]
where \([x^K]\) extracts the coefficient of \(x^K\). The law \(P_{\sigma,h}\) equals the law obtained by drawing \((H,Z)\) from that actual joint marginal, drawing \(y\) from this conditional density, copying \(y_\ell\) to every recipient in \(C_\ell\), and assigning zero outcomes to source and padding units. Its \((H,Z)\) marginal equals the actual joint marginal, including detailed retained-edge subsets, full source labels, and every treatment coordinate. For almost every \((H,Z)\) under this marginal,
\[
M=0\ \Longrightarrow\ K=0\ \text{ and }\ \binom{M}{K}=1.
\]
\end{lemmav}

\begin{proof}[Proof of \cref{obj:lem:block-law}]
\begin{enumerate}
\item By \cref{obj:ass:assignment,obj:ass:audit,obj:ass:design-independence}, the joint design law is
\[
\operatorname{Law}(Z,W)
=
\left(\bigotimes_{j\in V}\operatorname{Bernoulli}(1/2)\right)
\otimes
\left(\bigotimes_{\substack{(j,i)\in V^2\\j\ne i}}
\operatorname{Bernoulli}(q)\right).
\]
Indeed, independence identifies the joint law with the product of its stated marginals. Each factor is a probability law throughout \(0\le q\le1\), so the joint law has total mass one. We use this product representation below. % lean: CausalSmith.Experimentation.ThinnedgraphAdditiveRiskFrontier.design_eq_thinnedDesign

\item The baseline support in \cref{obj:def:block-prior} gives
\[
|U_\ell|\le\frac14
\qquad(1\le\ell\le B)
\]
almost surely. The condition \(2Bd\le n\) makes the source and recipient sets disjoint. For every source partition, the incoming and outgoing degrees therefore satisfy
\[
|\{j:(j,i)\in G\}|
=
\begin{cases}
d,&i\in\bigcup_{\ell=1}^{B}C_\ell,\\
0,&\text{otherwise},
\end{cases}
\qquad
|\{i:(j,i)\in G\}|
=
\begin{cases}
d,&j\text{ is a source},\\
0,&\text{otherwise}.
\end{cases}
\]
For \(i\in C_\ell\), positivity of \(d\) and nonnegativity of \(h\) give
\[
|a_i|+|t_i|+\sum_{j:(j,i)\in G}|b_{ij}|
=
\left|U_\ell-\frac{\sigma h}{2}\right|
+d\frac{h}{d}
\le\frac14+\frac{3h}{2}
\le1.
\]
Every other row has coefficient mass zero. These are precisely the requirements of \cref{obj:def:model}.

For the target, cancel the baselines in the all-treated versus all-control difference and interchange the finite edge sums:
\[
\begin{aligned}
\tau(\theta)
&=\frac1n\sum_{i\in V}\sum_{j:(j,i)\in G}\frac{\sigma h}{d}\\
&=\frac1n\sum_{j\in V}\sum_{i:(j,i)\in G}\frac{\sigma h}{d}
=\frac1n\sum_{\substack{j\in V\\j\text{ a source}}}\sigma h
=\frac{\sigma mh}{n}.
\end{aligned}
\]
There are \(m=Bd\) sources, each with \(d\) outgoing arrows. Thus the support and target conclusions hold almost surely under the prior. % lean: CausalSmith.Experimentation.ThinnedgraphAdditiveRiskFrontier.block_support

\item Define the source-label set and its reveal indicators by
\[
\mathcal J:=\bigcup_{\ell=1}^{B}S_\ell,
\qquad
R_j:=\mathbf1\{\exists i:(j,i)\in H\}
\quad(j\in\mathcal J).
\]
The set \(\mathcal J\) is fixed across partitions. Conditional on any partition \(S\), a source's indicator is the maximum of the \(d\) audit marks on its outgoing arrows. Hence
\[
\Pr(R_j=0\mid S)=(1-q)^d,
\qquad
\Pr(R_j=1\mid S)=p=1-(1-q)^d.
\]
The marks used by different sources are disjoint coordinates of the product audit law. Consequently,
\[
\operatorname{Law}\bigl((R_j)_{j\in\mathcal J}\mid S\bigr)
=
\bigotimes_{j\in\mathcal J}\operatorname{Bernoulli}(p).
\]
This law is independent of the chosen partition; averaging over its uniform probability law gives the same product law under the actual retained-graph marginal. The calculation includes \(q=0\) and \(q=1\). % lean: CausalSmith.Experimentation.ThinnedgraphAdditiveRiskFrontier.block_reveal_law

\item We next express the original mixture through its full graph-assignment marginal. Define the finite partition space, its associated graph, and the compatible partitions by
\[
\begin{aligned}
\mathcal S
&:=\left\{(S'_1,\ldots,S'_B):
\mathcal J=\bigsqcup_{\ell=1}^{B}S'_\ell,\ 
|S'_\ell|=d\right\},\\
G(S')&:=\bigcup_{\ell=1}^{B}(S'_\ell\times C_\ell)
\quad(S'\in\mathcal S),\\
\mathcal S(H)&:=\{S'\in\mathcal S:H\subseteq G(S')\}.
\end{aligned}
\]
Grouping an ordering of the \(m=Bd\) source labels into consecutive groups of size \(d\) exhibits an element of \(\mathcal S\); thus its uniform law has total mass one.

For \(S'\in\mathcal S\) and \(z\in\{0,1\}^{V}\), define the response shifts and their product density by
\[
\delta_\ell(S',z)
:=-\frac{\sigma h}{2}
+\frac{\sigma h}{d}\sum_{j\in S'_\ell}z_j,
\qquad
\rho_{S'}(y\mid z)
:=\prod_{\ell=1}^{B}f\bigl(y_\ell-\delta_\ell(S',z)\bigr),
\quad y\in\mathbb R^B.
\]
Substitution into the recipient response gives
\[
Y_i(z)=U_\ell+\delta_\ell(S',z)
\qquad(i\in C_\ell).
\]
The independent baselines therefore give density \(\rho_{S'}(\,\cdot\mid z)\).

For every detailed graph \(H\), assignment \(z\in\{0,1\}^{V}\), and response vector \(y\in\mathbb R^B\), define the copying map
\[
\mathcal C_{H,z}(y):=(H,z,\widetilde Y(y)),
\qquad
\widetilde Y_i(y):=
\begin{cases}
y_\ell,&i\in C_\ell,\\
0,&i\notin\bigcup_{\ell=1}^{B}C_\ell.
\end{cases}
\]
The recipient blocks are disjoint, so this definition is unambiguous. Define also the corresponding record measure
\[
\mathsf K_{S'}(H,z)
:=
(\mathcal C_{H,z})_{\#}
\bigl(\rho_{S'}(y\mid z)\,dy\bigr),
\qquad
S'\in\mathcal S,
\]
where the subscript \(\#\) denotes pushforward.

Every \(G(S')\) has \(md\) arrows. The product design yields
\[
\Pr(H,Z=z\mid S=S')
=
\begin{cases}
2^{-n}q^{|H|}(1-q)^{md-|H|},
&S'\in\mathcal S(H),\\
0,&S'\notin\mathcal S(H).
\end{cases}
\]
For a compatible partition, each retained true arrow contributes \(q\), each unretained true arrow contributes \(1-q\), and audit marks on all other pairs integrate to one. For an incompatible partition, an observed arrow outside \(G(S')\) has probability zero. We use the convention \(0^0=1\), so the displayed likelihood also applies at both audit endpoints.

Define
\[
\mu:=\operatorname{Law}(H,Z)
\]
under the actual experiment, averaging over the uniform source partition. On every atom of positive \(\mu\)-mass, the preceding common likelihood and uniform prior give
\[
\Pr(S=S'\mid H,Z=z)
=\frac1{|\mathcal S(H)|}
\qquad(S'\in\mathcal S(H)).
\]
Thus the original conditional record law is the uniform average
\[
\frac1{|\mathcal S(H)|}
\sum_{S'\in\mathcal S(H)}\mathsf K_{S'}(H,z).
\]
When \(\mathcal S(H)\) is empty, assign this expression the zero measure; the corresponding \(\mu\)-atom has mass zero. This convention makes the mixture representation valid over the whole observation space. % lean: CausalSmith.Experimentation.ThinnedgraphAdditiveRiskFrontier.blockMixtureLawOf_uniform_completion_bind

\item We now replace compatible partitions by completions of the original hidden labels. For a graph with \(\mathcal S(H)\ne\varnothing\), define
\[
\begin{aligned}
\mathcal R_\ell(H)
&:=\{j\in\mathcal J:\exists i\in C_\ell,\ (j,i)\in H\},\\
\mathcal J_{\mathrm{hid}}(H)
&:=\{j\in\mathcal J:\nexists i,\ (j,i)\in H\},\\
u_\ell&:=d-|\mathcal R_\ell(H)|,
&
M&:=\sum_{\ell=1}^{B}u_\ell,\\
A_\ell&:=\sum_{j\in\mathcal R_\ell(H)}(2z_j-1),
&
K&:=|\{j\in\mathcal J_{\mathrm{hid}}(H):z_j=1\}|.
\end{aligned}
\]
These agree with the quantities in the statement. A compatible partition places every label in \(\mathcal R_\ell(H)\) in its \(\ell\)-th block. In particular, these revealed sets are disjoint and have cardinality at most \(d\). For any \(S'\in\mathcal S(H)\),
\[
|\mathcal J_{\mathrm{hid}}(H)\cap S'_\ell|
=d-|\mathcal R_\ell(H)|=u_\ell,
\qquad
|\mathcal J_{\mathrm{hid}}(H)|=M.
\]

Define the hidden-label completion space by
\[
\mathcal G(H):=
\left\{g:\mathcal J_{\mathrm{hid}}(H)\to\{1,\ldots,B\}:
|g^{-1}(\ell)|=u_\ell\ \text{for every }\ell\right\}.
\]
Restriction to hidden labels and extension by the revealed labels give inverse maps
\[
\mathcal S(H)\longleftrightarrow\mathcal G(H),
\qquad
S'_\ell=\mathcal R_\ell(H)\cup g^{-1}(\ell).
\]
Count completions by ordering all \(M\) hidden labels and cutting the ordering into blocks of lengths \(u_1,\ldots,u_B\). Each completion arises from exactly \(\prod_\ell u_\ell!\) such orderings, because labels may be reordered within each block. Therefore
\[
|\mathcal G(H)|\prod_{\ell=1}^{B}u_\ell!=M!,
\qquad
\frac1{|\mathcal G(H)|}
=\frac{\prod_{\ell=1}^{B}u_\ell!}{M!}.
\]
The uniform compatible-partition average is consequently the same as the uniform hidden-label completion average. This identity includes \(M=0\), with \(0!=1\). For graphs with \(\mathcal S(H)=\varnothing\), the completion record measure is defined to be zero, as in the preceding step. % lean: CausalSmith.Experimentation.ThinnedgraphAdditiveRiskFrontier.blockMixtureLawOf_hidden_completion_bind

\item Fix \(H,z\) with \(\mathcal S(H)\ne\varnothing\). For \(g\in\mathcal G(H)\), define its hidden treated-count vector by
\[
k_\ell(g):=|\{j\in g^{-1}(\ell):z_j=1\}|,
\qquad
k(g):=(k_1(g),\ldots,k_B(g)).
\]
Then
\[
0\le k_\ell(g)\le u_\ell,
\qquad
\sum_{\ell=1}^{B}k_\ell(g)=K.
\]
More generally, for any compatible partition,
\[
K
=\sum_{\ell=1}^{B}
|\{j\in\mathcal J_{\mathrm{hid}}(H)\cap S'_\ell:z_j=1\}|
\le\sum_{\ell=1}^{B}u_\ell=M,
\]
so \(\binom{M}{K}>0\).

For an integer vector \(k=(k_1,\ldots,k_B)\) with \(0\le k_\ell\le u_\ell\), define its completion fiber by
\[
\mathcal G_k(H,z):=\{g\in\mathcal G(H):k(g)=k\}.
\]
If \(\sum_\ell k_\ell\ne K\), this fiber is empty. If \(\sum_\ell k_\ell=K\), restrict a completion separately to the \(K\) treated and \(M-K\) control labels. The two restrictions have block sizes \(k_\ell\) and \(u_\ell-k_\ell\), respectively, and together determine the completion. Applying the preceding ordering count to these two label sets gives
\[
|\mathcal G_k(H,z)|
\prod_{\ell=1}^{B}k_\ell!(u_\ell-k_\ell)!
=K!(M-K)!.
\]
Combining this identity with the total completion count gives
\[
\frac{|\mathcal G_k(H,z)|}{|\mathcal G(H)|}
=
\frac{K!(M-K)!\prod_{\ell=1}^{B}u_\ell!}
{M!\prod_{\ell=1}^{B}k_\ell!(u_\ell-k_\ell)!}
=
\frac{\prod_{\ell=1}^{B}\binom{u_\ell}{k_\ell}}
{\binom{M}{K}}
\quad\text{when }\sum_{\ell=1}^{B}k_\ell=K.
\]
All factorials are positive, including those of zero, so these identities also cover \(K=0\) and \(K=M\).

For the partition associated with \(g\), split the source-sign sum into revealed and hidden labels:
\[
\sum_{j\in S'_\ell}(2z_j-1)
=A_\ell+2k_\ell(g)-u_\ell.
\]
Since \(|S'_\ell|=d\), its response shift reduces to
\[
\delta_\ell(S',z)
=\frac{\sigma h}{2d}
\bigl(A_\ell+2k_\ell(g)-u_\ell\bigr).
\]
Thus its product response density depends on the completion through \(k(g)\). Regrouping the uniform completion average by these fibers gives
\[
\lambda_{\sigma,h}(y\mid H,z)
=
\frac1{\binom{M}{K}}
\sum_{\substack{0\le k_\ell\le u_\ell\\
\sum_{\ell=1}^{B}k_\ell=K}}
\prod_{\ell=1}^{B}
\binom{u_\ell}{k_\ell}
f\!\left(
y_\ell-\frac{\sigma h(A_\ell+2k_\ell-u_\ell)}{2d}
\right).
\]
Each translated product density integrates to one, and the fiber weights sum to one because the fibers partition \(\mathcal G(H)\). Hence this is a probability density.

Let \(x\) be a polynomial indeterminate, and let \([x^K]\) denote extraction of its degree-\(K\) coefficient. Expanding the finite product selects exactly the vectors whose coordinates sum to \(K\), so
\[
\lambda_{\sigma,h}(y\mid H,z)
=
\frac{
[x^K]\displaystyle\prod_{\ell=1}^{B}
\left\{
\sum_{k=0}^{u_\ell}
\binom{u_\ell}{k}
f\!\left(y_\ell-\frac{\sigma h(A_\ell+2k-u_\ell)}{2d}\right)x^k
\right\}
}{
\binom{M}{K}
}.
\]
This proves the stated density by averaging the translated baseline densities over the original labeled completions. % lean: CausalSmith.Experimentation.ThinnedgraphAdditiveRiskFrontier.hiddenCompletion_density_average

\item Every graph generated by the experiment is a subset of the graph of its generating partition. Therefore
\[
\mu\{(H,z):\mathcal S(H)\ne\varnothing\}=1.
\]
On this full-mass set, the conditional completion record measure equals
\[
(\mathcal C_{H,z})_{\#}
\bigl(\lambda_{\sigma,h}(y\mid H,z)\,dy\bigr),
\]
because a finite average of pushforward density measures is the pushforward of their averaged density. The preceding mixture representation consequently gives, for every measurable set \(\mathcal E_{\mathrm{rec}}\) of complete records,
\[
P_{\sigma,h}(\mathcal E_{\mathrm{rec}})
=
\sum_{H,z}\mu\{(H,z)\}
\int_{\mathbb R^B}
\mathbf1\{\mathcal C_{H,z}(y)\in\mathcal E_{\mathrm{rec}}\}
\lambda_{\sigma,h}(y\mid H,z)\,dy.
\]
The sum ranges over all detailed labeled retained graphs and all assignments in \(\{0,1\}^{V}\); zero-mass atoms contribute zero. This is exactly the reconstruction asserted in the statement. % lean: CausalSmith.Experimentation.ThinnedgraphAdditiveRiskFrontier.hiddenCompletionKernel_eq_blockDensity

\item Projection of an original record onto \((H,Z)\) depends on the partition and design draw, and is constant as the baselines, sign, and amplitude vary. Integrating the independent baseline probability law therefore yields
\[
\operatorname{Law}_{P_{\sigma,h}}(H,Z)
=
\frac1{|\mathcal S|}
\sum_{S'\in\mathcal S}
\operatorname{Law}(H,Z\mid S=S')
=\mu.
\]
Thus the marginal is the actual joint law of the detailed graph and every treatment coordinate. % lean: CausalSmith.Experimentation.ThinnedgraphAdditiveRiskFrontier.block_mixture_graphAssignMarginal

\item Finally, on the full-mass compatible-graph set, the bound already obtained gives
\[
0\le K\le M.
\]
Hence
\[
M=0\quad\Longrightarrow\quad K=0,
\qquad
\binom{M}{K}=\binom00=1.
\]
In this case every \(u_\ell\) is zero and the completion space consists of the unique map on the empty hidden-label set. The density formula therefore has a degree-zero numerator and the asserted denominator. % lean: CausalSmith.Experimentation.ThinnedgraphAdditiveRiskFrontier.block_hidden_zero_endpoint
\end{enumerate}
\end{proof}

\Cref{obj:lem:block-law} establishes support in the schedule class and a constant causal target under each prior. The construction's nonzero degrees are \(d\), and each recipient's absolute coefficient mass is bounded by
\[
\frac14+\frac h2+h\le\frac58
\qquad(0\le h\le1/4).
\]
Each of the \(m\) recipients has all-treated-versus-all-control contrast \(\sigma h\), giving the target \(\sigma mh/n\). The opposite-sign priors therefore have target separation \(2mh/n\), including at the amplitude specified in \cref{obj:def:testing-handle}. These are the support and separation properties used by the risk reduction in \cref{obj:lem:full-record-two-prior-risk}.

The conditional likelihood expresses how the observed assignments constrain the remaining source allocation. For a block receiving \(k\) hidden treated sources, the hidden sign sum is \(2k-u_\ell\); together with the revealed sum \(A_\ell\), it determines the translation of the baseline density. The coefficient extraction admits precisely those blockwise counts whose sum equals the observed global count \(K\). Its binomial normalization accounts for allocating that fixed treated count across the remaining capacities. At zero remaining capacity, the likelihood uses the stated \(K=0\) convention.

Finally, the reconstruction clause in \cref{obj:lem:block-law} preserves the complete observation channel. Detailed audit subsets and retained-edge labels enter through the actual graph marginal, every treatment coordinate remains in the joint marginal, and each generated block response is copied to its full recipient group. Although the conditional response formula summarizes the graph through observable capacities and sign sums, the reconstructed law retains the detailed graph itself. The resulting testing experiment therefore accommodates estimation procedures that use retained labels and repeated noiseless responses to infer source associations.


\section{Translation bounds and hidden-allocation contraction}

The conditional response law in \cref{obj:lem:block-law} connects the complete observation experiment to translations of the block baseline density. The bounds below quantify the information carried by these translations and by the remaining source allocation. Throughout, the retained graph retains its labels and detailed edge subsets, and averaging uses its actual marginal law.

\subsection{Baseline translations and common-marginal comparisons}

The first ingredient measures how a change in the block response location affects its square-root density. We fix the density and derivative conventions before stating the translation bound.



The square-root density in \cref{obj:lem:baseline-translation-affinity} supplies the geometry for comparing translated block responses. To carry these comparisons through conditioning and mixing, we also fix the probability-distance conventions.



These conventions allow one bound to address a single translated response, a vector of independent translated responses, and conditional response laws sharing the same marginal for the conditioning variables. The following statement collects the corresponding comparisons.

\begin{lemmav}{lem:baseline-translation-affinity}[Baseline translation and mixture bounds]
Let \(f\) be the baseline density and \(\varphi\) its square root:
\[
f(w)=
\begin{cases}
4\cos^2(2\pi w),& |w|\le 1/4,\\
0,& |w|>1/4,
\end{cases}
\qquad
\varphi(w)=\sqrt{f(w)}.
\]
Then \(\varphi\) is globally \(4\pi\)-Lipschitz, and
\[
\int_{-1/4}^{1/4}\bigl(4\pi\sin(2\pi w)\bigr)^2\,dw
=4\pi^2.
\]
For every \(s,t\in\mathbb R\),
\[
\int_{\mathbb R}
\bigl(\sqrt{f(w-s)}-\sqrt{f(w-t)}\bigr)^2\,dw
\le 4\pi^2(s-t)^2.
\]
For every nonnegative integer \(B\) and every pair of vectors
\(s,t\in\mathbb R^B\),
\[
\int_{\mathbb R^B}
\left(
\sqrt{\prod_{\ell=1}^{B}f(y_\ell-s_\ell)}
-
\sqrt{\prod_{\ell=1}^{B}f(y_\ell-t_\ell)}
\right)^2\,dy
\le
\sum_{\ell=1}^{B}4\pi^2(s_\ell-t_\ell)^2.
\]
Here real and Euclidean integrals use Lebesgue measure, including the usual empty-product convention when \(B=0\).

For probability laws, let \(\operatorname{TV}\) denote total variation, defined as the supremum of the absolute difference of probabilities over measurable events. The unhalved squared Hellinger distance between laws with densities \(v_+\) and \(v_-\) relative to a common measure \(\mu\) is
\[
\int_\Omega
\bigl(\sqrt{v_+(x)}-\sqrt{v_-(x)}\bigr)^2\,d\mu(x).
\]
For every measurable space \(\Omega\), every sigma-finite measure \(\mu\), and every pair of nonnegative measurable densities \(v_+,v_-\) satisfying
\[
\int_\Omega v_+(x)\,d\mu(x)
=
\int_\Omega v_-(x)\,d\mu(x)
=1,
\]
one has
\[
\operatorname{TV}\bigl(v_+\,\mu,v_-\,\mu\bigr)
\le
\left[
\int_\Omega
\bigl(\sqrt{v_+(x)}-\sqrt{v_-(x)}\bigr)^2\,d\mu(x)
\right]^{1/2}.
\]

For the common-marginal version, let \(\Xi,\Omega\) be measurable spaces and suppose:
\begin{itemize}
\item \textbf{(Reference measures.)} \(\nu\) is a probability measure on \(\Xi\), and \(\mu\) is a sigma-finite measure on \(\Omega\).
\item \textbf{(Conditional densities.)} \(v_+,v_-:\Xi\times\Omega\to\mathbb R\) are jointly measurable and nonnegative.
\item \textbf{(Normalization.)} For every \(\xi\in\Xi\),
\[
\int_\Omega v_+(\xi,x)\,d\mu(x)
=
\int_\Omega v_-(\xi,x)\,d\mu(x)
=1.
\]
\end{itemize}
The joint laws are
\[
\nu(d\xi)\,v_\pm(\xi,x)\,\mu(dx),
\]
and their response mixtures are
\[
\int_\Xi v_\pm(\xi,x)\,\mu(dx)\,\nu(d\xi).
\]
The unhalved squared Hellinger distance between these joint laws equals
\[
\int_\Xi
\left[
\int_\Omega
\bigl(\sqrt{v_+(\xi,x)}-\sqrt{v_-(\xi,x)}\bigr)^2\,d\mu(x)
\right]d\nu(\xi),
\]
and the response mixtures satisfy
\[
\operatorname{TV}\left(
\int_\Xi v_+(\xi,x)\,\mu(dx)\,\nu(d\xi),
\int_\Xi v_-(\xi,x)\,\mu(dx)\,\nu(d\xi)
\right)
\le
\left[
\int_\Xi
\left[
\int_\Omega
\bigl(\sqrt{v_+(\xi,x)}-\sqrt{v_-(\xi,x)}\bigr)^2\,d\mu(x)
\right]d\nu(\xi)
\right]^{1/2}.
\]

Define the derivative representative, with value zero at the two support endpoints, by
\[
\varphi'(w)=
\begin{cases}
-4\pi\sin(2\pi w),& |w|<1/4,\\
0,& |w|\ge 1/4.
\end{cases}
\]
At every \(w\in\mathbb R\setminus\{-1/4,1/4\}\), \(\varphi\) is differentiable with derivative \(\varphi'(w)\), and
\[
\int_{\mathbb R}\varphi'(w)^2\,dw=4\pi^2.
\]
\end{lemmav}

\begin{proof}[Proof of \cref{obj:lem:baseline-translation-affinity}]
\begin{enumerate}
\item \emph{The global Lipschitz bound.}
Define the clamping map on \(\mathbb R\) by
\[
\operatorname{clamp}(w)
=\max\{-1/4,\min\{1/4,w\}\}.
\]
On the support interval, \(\cos(2\pi w)\ge0\), so taking the square root of \(f(w)\) gives \(2\cos(2\pi w)\). Below the interval, the clamping map equals \(-1/4\); above it, the clamping map equals \(1/4\). Since the cosine vanishes at both corresponding angles, these cases give the identity
\[
\varphi(w)=2\cos\bigl(2\pi\operatorname{clamp}(w)\bigr)
\qquad(w\in\mathbb R).
\]
Taking a minimum or maximum with a fixed real number preserves the Lipschitz constant \(1\). Consequently, for \(w_1,w_2\in\mathbb R\),
\[
\bigl|\operatorname{clamp}(w_1)-\operatorname{clamp}(w_2)\bigr|
\le |w_1-w_2|.
\]
Using \(|\cos \omega_1-\cos \omega_2|\le|\omega_1-\omega_2|\) for real angles \(\omega_1,\omega_2\), we obtain
\[
\begin{aligned}
|\varphi(w_1)-\varphi(w_2)|
&\le 2\bigl|2\pi\operatorname{clamp}(w_1)
             -2\pi\operatorname{clamp}(w_2)\bigr|\\
&=4\pi\bigl|\operatorname{clamp}(w_1)
             -\operatorname{clamp}(w_2)\bigr|\\
&\le4\pi|w_1-w_2|.
\end{aligned}
\]
% lean: CausalSmith.Experimentation.ThinnedgraphAdditiveRiskFrontier.baseline_sqrt_lipschitz

\item \emph{The sine energy.}
Changing variables to the real angle \(\omega=2\pi w\), and using the antiderivative
\(\omega/2-\sin\omega\cos\omega/2\) of \(\sin^2\omega\), gives
\[
\begin{aligned}
\int_{-1/4}^{1/4}(4\pi\sin(2\pi w))^2\,dw
&=8\pi\int_{-\pi/2}^{\pi/2}\sin^2\omega\,d\omega\\
&=8\pi
  \left[\frac{\omega}{2}
        -\frac{\sin\omega\cos\omega}{2}\right]_{-\pi/2}^{\pi/2}\\
&=4\pi^2.
\end{aligned}
\]
The closed-interval Lebesgue integral has the same value as this interval integral because its endpoints have measure zero.
% lean: CausalSmith.Experimentation.ThinnedgraphAdditiveRiskFrontier.baseline_derivative_energy

\item \emph{An integral representation.}
For the translation argument, define the closed-support function
\[
\psi(w)=
\begin{cases}
-4\pi\sin(2\pi w),& |w|\le1/4,\\
0,& |w|>1/4,
\end{cases}
\qquad w\in\mathbb R.
\]
Both \(\psi\) and \(\psi^2\) are measurable and integrable by compact support and continuity of the sine expression. The preceding calculation gives
\[
\int_{\mathbb R}\psi(w)^2\,dw
=\int_{-1/4}^{1/4}(4\pi\sin(2\pi w))^2\,dw
=4\pi^2.
\]

Fix \(\rho_0,\rho_1\in\mathbb R\) with \(\rho_0\le\rho_1\), and define
\[
\rho_L=\max\{-1/4,\rho_0\},
\qquad
\rho_U=\min\{1/4,\rho_1\}.
\]
Restricting the integral to the support yields
\[
\int_{\rho_0}^{\rho_1}\psi(w)\,dw
=\int_{[\rho_L,\rho_U]}-4\pi\sin(2\pi w)\,dw,
\]
where the set interval is empty when \(\rho_U<\rho_L\).
If \(\rho_1<-1/4\), this interval is empty and both square-root density values are zero. If \(\rho_0>1/4\), the same conclusion holds. In the remaining case,
\(\rho_1\ge-1/4\) and \(\rho_0\le1/4\), so \(\rho_L\le\rho_U\). Integrating the derivative of \(2\cos(2\pi w)\) gives
\[
\begin{aligned}
\int_{\rho_0}^{\rho_1}\psi(w)\,dw
&=2\cos(2\pi\rho_U)-2\cos(2\pi\rho_L)\\
&=\varphi(\rho_1)-\varphi(\rho_0),
\end{aligned}
\]
because \(\rho_U=\operatorname{clamp}(\rho_1)\) and
\(\rho_L=\operatorname{clamp}(\rho_0)\) in this case. Thus, for every ordered pair of real endpoints,
\[
\int_{\rho_0}^{\rho_1}\psi(w)\,dw
=\varphi(\rho_1)-\varphi(\rho_0).
\]
% lean: CausalSmith.Experimentation.ThinnedgraphAdditiveRiskFrontier.baselineSqrtDerivative_integral

\item \emph{The translation estimate.}
Fix \(w\in\mathbb R\) and \(\rho_0\le\rho_1\), and define
\[
\mathcal I_{\mathrm{int}}(w;\rho_0,\rho_1)
=\int_{\rho_0}^{\rho_1}\psi(w+\rho)\,d\rho.
\]
When \(\rho_0=\rho_1\), this integral is zero. When \(\rho_0<\rho_1\), expanding a nonnegative square gives
\[
\begin{aligned}
0
&\le
\int_{\rho_0}^{\rho_1}
\left((\rho_1-\rho_0)\psi(w+\rho)
      -\mathcal I_{\mathrm{int}}(w;\rho_0,\rho_1)\right)^2\,d\rho\\
&=(\rho_1-\rho_0)
\left[
(\rho_1-\rho_0)\int_{\rho_0}^{\rho_1}\psi(w+\rho)^2\,d\rho
-\mathcal I_{\mathrm{int}}(w;\rho_0,\rho_1)^2
\right].
\end{aligned}
\]
Dividing by the positive interval length, and including the zero-length case, proves
\[
\mathcal I_{\mathrm{int}}(w;\rho_0,\rho_1)^2
\le
(\rho_1-\rho_0)\int_{\rho_0}^{\rho_1}\psi(w+\rho)^2\,d\rho.
\]
The integral representation from the preceding step identifies the left side with
\((\varphi(w+\rho_1)-\varphi(w+\rho_0))^2\).
Translation invariance gives
\[
\int_{\mathbb R}\psi(w+\rho)^2\,dw
=\int_{\mathbb R}\psi(w)^2\,dw.
\]
In particular, the square integrand is integrable on
\(\mathbb R\times[\rho_0,\rho_1]\). Integrating the pointwise bound and applying Fubini therefore gives
\[
\begin{aligned}
\int_{\mathbb R}
(\varphi(w+\rho_1)-\varphi(w+\rho_0))^2\,dw
&\le
(\rho_1-\rho_0)
\int_{\mathbb R}\int_{\rho_0}^{\rho_1}
\psi(w+\rho)^2\,d\rho\,dw\\
&=(\rho_1-\rho_0)
\int_{\rho_0}^{\rho_1}\int_{\mathbb R}
\psi(w+\rho)^2\,dw\,d\rho\\
&=(\rho_1-\rho_0)^2\int_{\mathbb R}\psi(w)^2\,dw\\
&=4\pi^2(\rho_1-\rho_0)^2.
\end{aligned}
\]
For real shifts \(\rho_{\mathrm{shift},1},\rho_{\mathrm{shift},2}\) with \(\rho_{\mathrm{shift},2}\le\rho_{\mathrm{shift},1}\), take \(\rho_0=-\rho_{\mathrm{shift},1}\) and \(\rho_1=-\rho_{\mathrm{shift},2}\). Reversing the difference inside the square then gives
\[
\int_{\mathbb R}(\varphi(w-\rho_{\mathrm{shift},1})-\varphi(w-\rho_{\mathrm{shift},2}))^2\,dw
\le4\pi^2(\rho_{\mathrm{shift},1}-\rho_{\mathrm{shift},2})^2.
\]
For \(\rho_{\mathrm{shift},1}\le\rho_{\mathrm{shift},2}\), apply the same argument with the shifts exchanged; both squared differences are unchanged. This covers all \(\rho_{\mathrm{shift},1},\rho_{\mathrm{shift},2}\), including equality.
% lean: CausalSmith.Experimentation.ThinnedgraphAdditiveRiskFrontier.baseline_translation_energy

\item \emph{Products of translated densities.}
Fix \(B\ge0\) and vectors \(\boldsymbol s,\boldsymbol t\in\mathbb R^B\). For \(1\le\ell\le B\) and \(w\in\mathbb R\), define
\[
f_{\ell,\boldsymbol s}(w)=f(w-s_\ell),
\qquad
f_{\ell,\boldsymbol t}(w)=f(w-t_\ell).
\]
The baseline density is nonnegative and integrable. Its normalization follows from
\[
\begin{aligned}
\int_{\mathbb R}f(w)\,dw
&=\frac{2}{\pi}\int_{-\pi/2}^{\pi/2}\cos^2\omega\,d\omega\\
&=\frac{2}{\pi}
  \left[\frac{\omega}{2}
        +\frac{\sin\omega\cos\omega}{2}\right]_{-\pi/2}^{\pi/2}
=1.
\end{aligned}
\]
Translation preserves integrability and the integral, so
\[
\int_{\mathbb R}f_{\ell,\boldsymbol s}(w)\,dw
=\int_{\mathbb R}f_{\ell,\boldsymbol t}(w)\,dw=1.
\]
Define the coordinate affinities and energies by
\[
\alpha_{\mathrm{aff},\ell}
=\int_{\mathbb R}
\sqrt{f_{\ell,\boldsymbol s}(w)f_{\ell,\boldsymbol t}(w)}\,dw,
\qquad
\mathcal E_{\mathrm{Hel},\ell}
=\int_{\mathbb R}
\left(\sqrt{f_{\ell,\boldsymbol s}(w)}
      -\sqrt{f_{\ell,\boldsymbol t}(w)}\right)^2\,dw.
\]
The square roots are square integrable, so their product is integrable. Expanding the square and using normalization yields
\[
\mathcal E_{\mathrm{Hel},\ell}
=2(1-\alpha_{\mathrm{aff},\ell}).
\]
The affinity is nonnegative, and the energy is nonnegative; hence
\[
0\le\alpha_{\mathrm{aff},\ell}\le1.
\]

For \(y\in\mathbb R^B\), define the product densities
\[
v_{\boldsymbol s}(y)=\prod_{\ell=1}^B f_{\ell,\boldsymbol s}(y_\ell),
\qquad
v_{\boldsymbol t}(y)=\prod_{\ell=1}^B f_{\ell,\boldsymbol t}(y_\ell).
\]
Product integration shows that both densities are integrable and normalized:
\[
\int_{\mathbb R^B}v_{\boldsymbol s}(y)\,dy
=\prod_{\ell=1}^B\int_{\mathbb R}f_{\ell,\boldsymbol s}(w)\,dw=1,
\]
and likewise for \(v_{\boldsymbol t}\).
Nonnegativity allows the square root to factor, so another application of product integration gives
\[
\begin{aligned}
\int_{\mathbb R^B}\sqrt{v_{\boldsymbol s}(y)v_{\boldsymbol t}(y)}\,dy
&=\int_{\mathbb R^B}
\prod_{\ell=1}^B
\sqrt{f_{\ell,\boldsymbol s}(y_\ell)
      f_{\ell,\boldsymbol t}(y_\ell)}\,dy\\
&=\prod_{\ell=1}^B\alpha_{\mathrm{aff},\ell}.
\end{aligned}
\]
Expanding the product-density square therefore gives
\[
\int_{\mathbb R^B}
\left(\sqrt{v_{\boldsymbol s}(y)}
      -\sqrt{v_{\boldsymbol t}(y)}\right)^2\,dy
=2\left(1-\prod_{\ell=1}^B\alpha_{\mathrm{aff},\ell}\right).
\]

For completeness, the product-defect bound follows by induction over subsets
\(\mathcal J\subseteq\{1,\ldots,B\}\). Every such product lies in \([0,1]\), and the empty subset gives equality. For
\(\ell_0\in\{1,\ldots,B\}\setminus\mathcal J\), the induction step is
\[
\begin{aligned}
1-\alpha_{\mathrm{aff},\ell_0}
      \prod_{\ell\in\mathcal J}\alpha_{\mathrm{aff},\ell}
&=(1-\alpha_{\mathrm{aff},\ell_0})
  +\alpha_{\mathrm{aff},\ell_0}
     \left(1-\prod_{\ell\in\mathcal J}\alpha_{\mathrm{aff},\ell}\right)\\
&\le (1-\alpha_{\mathrm{aff},\ell_0})
     +1-\prod_{\ell\in\mathcal J}\alpha_{\mathrm{aff},\ell}\\
&\le (1-\alpha_{\mathrm{aff},\ell_0})
     +\sum_{\ell\in\mathcal J}(1-\alpha_{\mathrm{aff},\ell}).
\end{aligned}
\]
Consequently,
\[
\begin{aligned}
\int_{\mathbb R^B}
\left(\sqrt{v_{\boldsymbol s}(y)}
      -\sqrt{v_{\boldsymbol t}(y)}\right)^2\,dy
&\le\sum_{\ell=1}^B\mathcal E_{\mathrm{Hel},\ell}\\
&\le\sum_{\ell=1}^B4\pi^2(s_\ell-t_\ell)^2,
\end{aligned}
\]
where the last inequality is the scalar translation estimate. When \(B=0\), both product densities equal \(1\) on the one-point space \(\mathbb R^0\), and both sides of the bound are zero.
% lean: CausalSmith.Experimentation.ThinnedgraphAdditiveRiskFrontier.hellingerEnergy_pi_le_sum

\item \emph{Total variation for normalized densities.}
Take \(v_+,v_-:\Omega\to[0,\infty)\) satisfying the stated measurability and normalization conditions. For \(\omega_{\mathrm{obs}}\in\Omega\), define
\[
\delta_{\mathrm{dens}}(\omega_{\mathrm{obs}})=v_+(\omega_{\mathrm{obs}})-v_-(\omega_{\mathrm{obs}}),
\qquad
\alpha_{\mathrm{aff}}
=\int_\Omega\sqrt{v_+(\omega_{\mathrm{obs}})v_-(\omega_{\mathrm{obs}})}\,d\mu(\omega_{\mathrm{obs}}).
\]
Normalization makes both densities integrable and gives
\[
\int_\Omega\delta_{\mathrm{dens}}\,d\mu=0.
\]
For any measurable event \(E_{\mathrm{ev}}\subseteq\Omega\), splitting this integral over the event and its complement gives
\[
\int_{E_{\mathrm{ev}}^c}\delta_{\mathrm{dens}}\,d\mu
=-\int_{E_{\mathrm{ev}}}\delta_{\mathrm{dens}}\,d\mu.
\]
Bounding the absolute value of each integral by the integral of the absolute value yields
\[
2\left|\int_{E_{\mathrm{ev}}}\delta_{\mathrm{dens}}\,d\mu\right|
\le
\int_{E_{\mathrm{ev}}}|\delta_{\mathrm{dens}}|\,d\mu
+\int_{E_{\mathrm{ev}}^c}|\delta_{\mathrm{dens}}|\,d\mu
=\int_\Omega|\delta_{\mathrm{dens}}|\,d\mu.
\]
Since the difference of event probabilities is the corresponding set integral, taking the supremum over measurable events proves
\[
\operatorname{TV}(v_+\mu,v_-\mu)
\le\frac12\int_\Omega|v_+-v_-|\,d\mu.
\]

The square roots belong to \(L^2(\mu)\), and pointwise
\[
|v_+-v_-|
=|\sqrt{v_+}-\sqrt{v_-}|\,
 |\sqrt{v_+}+\sqrt{v_-}|.
\]
Cauchy--Schwarz thus gives
\[
\int_\Omega|v_+-v_-|\,d\mu
\le
\left[\int_\Omega(\sqrt{v_+}-\sqrt{v_-})^2\,d\mu\right]^{1/2}
\left[\int_\Omega(\sqrt{v_+}+\sqrt{v_-})^2\,d\mu\right]^{1/2}.
\]
Expanding the first square gives
\[
\int_\Omega(\sqrt{v_+}-\sqrt{v_-})^2\,d\mu
=2(1-\alpha_{\mathrm{aff}}).
\]
For the second square, the pointwise inequality
\[
(\sqrt{v_+}+\sqrt{v_-})^2\le2v_++2v_-
\]
follows from \((\sqrt{v_+}-\sqrt{v_-})^2\ge0\). Therefore,
\[
\int_\Omega(\sqrt{v_+}+\sqrt{v_-})^2\,d\mu\le4.
\]
Combining these bounds and then substituting the energy identity proves
\[
\operatorname{TV}(v_+\mu,v_-\mu)
\le\sqrt{2(1-\alpha_{\mathrm{aff}})}
=
\left[\int_\Omega(\sqrt{v_+}-\sqrt{v_-})^2\,d\mu\right]^{1/2}.
\]
% lean: CausalSmith.Experimentation.ThinnedgraphAdditiveRiskFrontier.density_tv_le_sqrt_hellingerEnergy

\item \emph{Common-marginal averaging.}
Under the common-marginal hypotheses, use marginal and response coordinates
\[
\zeta_{\mathrm{marg}}\in\Xi,
\qquad
\omega_{\mathrm{obs}}\in\Omega,
\]
and define the joint and response laws by
\[
\mathsf P_{\mathrm{joint},\pm}(d\zeta_{\mathrm{marg}},d\omega_{\mathrm{obs}})
=\nu(d\zeta_{\mathrm{marg}})v_\pm(\zeta_{\mathrm{marg}},\omega_{\mathrm{obs}})\mu(d\omega_{\mathrm{obs}}),
\qquad
\mathsf P_{\mathrm{mix},\pm}(d\omega_{\mathrm{obs}})
=\left[\int_\Xi v_\pm(\zeta_{\mathrm{marg}},\omega_{\mathrm{obs}})\,d\nu(\zeta_{\mathrm{marg}})\right]\mu(d\omega_{\mathrm{obs}}).
\]
Joint measurability and integration of nonnegative functions give, for every measurable
\(D_{\mathrm{ev}}\subseteq\Xi\times\Omega\),
\[
\mathsf P_{\mathrm{joint},\pm}(D_{\mathrm{ev}})
=\int_\Xi\int_\Omega
\mathbf1_{D_{\mathrm{ev}}}(\zeta_{\mathrm{marg}},\omega_{\mathrm{obs}})v_\pm(\zeta_{\mathrm{marg}},\omega_{\mathrm{obs}})\,d\mu(\omega_{\mathrm{obs}})\,d\nu(\zeta_{\mathrm{marg}})
=\int_{D_{\mathrm{ev}}}v_\pm(\zeta_{\mathrm{marg}},\omega_{\mathrm{obs}})\,d(\nu\otimes\mu)(\zeta_{\mathrm{marg}},\omega_{\mathrm{obs}}).
\]
Fiber normalization gives
\[
\mathsf P_{\mathrm{joint},\pm}(\Xi\times\Omega)
=\int_\Xi 1\,d\nu=1.
\]
Thus the joint laws are probability laws with integrable densities \(v_\pm\) relative to the sigma-finite product reference measure.

To identify their squared Hellinger distance, define measures on \(\Xi\times\Omega\) by
\[
\Lambda_{\mathrm{joint}}=\nu\otimes\mu,
\qquad
\Lambda_{\mathrm{sum}}
=\mathsf P_{\mathrm{joint},+}+\mathsf P_{\mathrm{joint},-},
\]
and define their Radon--Nikodym density representatives by
\[
r_{\mathrm{sum}}
=\frac{d\Lambda_{\mathrm{sum}}}{d\Lambda_{\mathrm{joint}}},
\qquad
r_\pm
=\frac{d\mathsf P_{\mathrm{joint},\pm}}{d\Lambda_{\mathrm{sum}}}.
\]
These are nonnegative real-valued functions on \(\Xi\times\Omega\), with any infinite derivative values assigned zero; the relevant derivatives are finite almost everywhere relative to their reference measures. The Radon--Nikodym chain rule gives
\[
v_\pm(\zeta_{\mathrm{marg}},\omega_{\mathrm{obs}})=r_\pm(\zeta_{\mathrm{marg}},\omega_{\mathrm{obs}})r_{\mathrm{sum}}(\zeta_{\mathrm{marg}},\omega_{\mathrm{obs}})
\quad\text{for }\Lambda_{\mathrm{joint}}\text{-almost every }(\zeta_{\mathrm{marg}},\omega_{\mathrm{obs}}).
\]
Hence, almost everywhere,
\[
(\sqrt{r_+}-\sqrt{r_-})^2r_{\mathrm{sum}}
=
\left(\sqrt{r_+r_{\mathrm{sum}}}
      -\sqrt{r_-r_{\mathrm{sum}}}\right)^2
=(\sqrt{v_+}-\sqrt{v_-})^2.
\]
Changing reference measure therefore identifies the squared Hellinger distance of the joint laws as
\[
\begin{aligned}
\int_{\Xi\times\Omega}
(\sqrt{r_+}-\sqrt{r_-})^2\,d\Lambda_{\mathrm{sum}}
&=\int_{\Xi\times\Omega}
(\sqrt{v_+}-\sqrt{v_-})^2\,d\Lambda_{\mathrm{joint}}\\
&=\int_\Xi
\left[\int_\Omega
(\sqrt{v_+(\zeta_{\mathrm{marg}},\omega_{\mathrm{obs}})}-\sqrt{v_-(\zeta_{\mathrm{marg}},\omega_{\mathrm{obs}})})^2\,d\mu(\omega_{\mathrm{obs}})\right]d\nu(\zeta_{\mathrm{marg}}).
\end{aligned}
\]
For the last equality, the square-root difference has an integrable square because the joint densities are integrable and their square roots belong to \(L^2(\Lambda_{\mathrm{joint}})\); Fubini applies.

Applying the preceding total-variation bound to these joint densities gives
\[
\operatorname{TV}(\mathsf P_{\mathrm{joint},+},\mathsf P_{\mathrm{joint},-})
\le
\left[
\int_\Xi\int_\Omega
(\sqrt{v_+(\zeta_{\mathrm{marg}},\omega_{\mathrm{obs}})}-\sqrt{v_-(\zeta_{\mathrm{marg}},\omega_{\mathrm{obs}})})^2\,d\mu(\omega_{\mathrm{obs}})\,d\nu(\zeta_{\mathrm{marg}})
\right]^{1/2}.
\]
Finally, for every measurable \(E_{\mathrm{ev}}\subseteq\Omega\), integration of the joint law gives
\[
\mathsf P_{\mathrm{mix},\pm}(E_{\mathrm{ev}})
=\mathsf P_{\mathrm{joint},\pm}(\Xi\times E_{\mathrm{ev}}).
\]
The cylinder event is measurable, so
\[
\begin{aligned}
\left|\mathsf P_{\mathrm{mix},+}(E_{\mathrm{ev}})
      -\mathsf P_{\mathrm{mix},-}(E_{\mathrm{ev}})\right|
&\le
\operatorname{TV}(\mathsf P_{\mathrm{joint},+},\mathsf P_{\mathrm{joint},-}).
\end{aligned}
\]
Taking the supremum over response events proves the asserted mixture bound.
% lean: CausalSmith.Experimentation.ThinnedgraphAdditiveRiskFrontier.common_marginal_hellinger_averaging

\item \emph{Pointwise differentiation away from the endpoints.}
Fix \(w\in\mathbb R\setminus\{-1/4,1/4\}\).
If \(w<-1/4\), then \(\varphi\) agrees with zero throughout a neighborhood of \(w\), and its derivative is \(0=\varphi'(w)\).
For \(w>-1/4\), split further at \(1/4\). If \(w<1/4\), the clamping identity gives
\[
\varphi(w_{\mathrm{loc}})=2\cos(2\pi w_{\mathrm{loc}})
\quad\text{for real }w_{\mathrm{loc}}\text{ in a neighborhood of }w.
\]
The chain rule yields
\[
\left.\frac{d}{dw_{\mathrm{loc}}}\bigl(2\cos(2\pi w_{\mathrm{loc}})\bigr)\right|_{w_{\mathrm{loc}}=w}
=-4\pi\sin(2\pi w)=\varphi'(w).
\]
If \(w>1/4\), then \(\varphi\) again agrees with zero throughout a neighborhood of \(w\), giving derivative \(0=\varphi'(w)\). The two equality cases are excluded by the choice of \(w\).
% lean: CausalSmith.Experimentation.ThinnedgraphAdditiveRiskFrontier.sqrt_cosSqDensity_hasDerivAt

\item \emph{The energy of the derivative representative.}
The endpoint-zero convention in the statement gives the pointwise identity
\[
\varphi'(w)^2
=\mathbf1_{(-1/4,1/4)}(w)
  \bigl(-4\pi\sin(2\pi w)\bigr)^2
\qquad(w\in\mathbb R).
\]
Integrating this indicator, replacing the open interval by the closed interval up to its two null endpoints, and removing the sign inside the square gives
\[
\begin{aligned}
\int_{\mathbb R}\varphi'(w)^2\,dw
&=\int_{(-1/4,1/4)}(-4\pi\sin(2\pi w))^2\,dw\\
&=\int_{[-1/4,1/4]}(4\pi\sin(2\pi w))^2\,dw\\
&=4\pi^2.
\end{aligned}
\]
% lean: CausalSmith.Experimentation.ThinnedgraphAdditiveRiskFrontier.baselineSqrtDerivAt_energy
\end{enumerate}
\end{proof}

The translation inequality in \cref{obj:lem:baseline-translation-affinity} gives a quadratic bound in the change of location. Its product version adds the squared changes across distinct block responses. The common-marginal identity then preserves the conditioning variables in the joint experiment: conditional squared Hellinger distances are averaged under their shared marginal. The response-mixture inequality supplies the corresponding comparison after those variables are integrated out. These are the distance comparisons used to translate conditional response bounds into testing bounds.

\subsection{The translated channel and all-order Walsh energy}

For a block with \(d\) source signs, the response translation depends on their sum. A uniform sign mixture provides a reference density against which every translated response can be expanded. We first define this family at a fixed admissible amplitude.



The normalized translated densities in \cref{obj:lem:walsh-channel-energy} admit an expansion in products of source signs. Symmetry makes the coefficient depend on the number of signs in the product. The same result records the coefficients and the two energy aggregates needed for the allocation bound.



In \cref{obj:lem:walsh-channel-energy}, the binomial weight counts the sign subsets of each order. The additional factor \(s\) in the order-weighted aggregate records their sizes. The result controls both aggregates while retaining every nonconstant order of the response expansion.

\begin{lemmav}{lem:walsh-channel-energy}[Walsh channel energy bound]
Let \(d\ge1\) be an integer and let \(0\le h\le1/4\). With \(f\) denoting the baseline density, define, for each sign vector \(\varepsilon\in\{-1,1\}^d\),
\[
F_\varepsilon(w)
=f\!\left(w-\frac{h}{2d}\sum_{j=1}^d\varepsilon_j\right),
\qquad
g_d(w)=2^{-d}\sum_{\varepsilon\in\{-1,1\}^d}F_\varepsilon(w).
\]
For every integer \(s\ge0\), define the Walsh coefficient function by
\[
a_s(w)=
\begin{cases}
2^{-d}\displaystyle\sum_{\varepsilon\in\{-1,1\}^d}
\frac{F_\varepsilon(w)}{g_d(w)}
\displaystyle\prod_{j=1}^{\min\{s,d\}}\varepsilon_j,
& g_d(w)\ne0,\\
0,&g_d(w)=0.
\end{cases}
\]
The empty product equals one. Define the integrated Walsh energies and their two aggregates by
\[
\gamma_s=\int_{\mathbb R}a_s(w)^2g_d(w)\,dw
\quad(1\le s\le d),
\qquad
\eta=\sum_{s=1}^d\binom ds\gamma_s,
\qquad
\eta_1=\sum_{s=1}^d s\binom ds\gamma_s.
\]
For every \(w\) with \(g_d(w)>0\), \(a_0(w)=1\), and for every \(\varepsilon\in\{-1,1\}^d\),
\[
\frac{F_\varepsilon(w)}{g_d(w)}
=
\sum_{E\subseteq\{1,\ldots,d\}}
a_{|E|}(w)\prod_{j\in E}\varepsilon_j.
\]
Moreover, for every integer \(s\ge1\),
\[
\int_{\mathbb R}a_s(w)g_d(w)\,dw=0,
\]
and the aggregate energies satisfy
\[
0\le\eta\le\eta_1\le\frac{12\pi^2h^2}{d}.
\]
\end{lemmav}

\begin{proof}[Proof of \cref{obj:lem:walsh-channel-energy}]
\begin{enumerate}
\item For every subset \(E\subseteq\{1,\ldots,d\}\), define its sign character by
\[
\psi_E(\varepsilon)=\prod_{j\in E}\varepsilon_j,
\qquad \varepsilon\in\{-1,1\}^d.
\]
Also define the initial coordinate sets
\[
E_s=\{1,\ldots,\min\{s,d\}\},
\qquad s\in\mathbb Z_{\ge0}.
\]
Fix \(w\in\mathbb R\) with \(g_d(w)>0\). The empty-product convention and the definition of \(g_d\) give
\[
a_0(w)
=2^{-d}\sum_{\varepsilon\in\{-1,1\}^d}
\frac{F_\varepsilon(w)}{g_d(w)}
=\frac{g_d(w)}{g_d(w)}
=1.
\]

Relabeling coordinates preserves the sign sum and hence \(F_\varepsilon(w)\). More explicitly, for any permutation
\[
\rho:\{1,\ldots,d\}\longrightarrow\{1,\ldots,d\},
\qquad
(\varepsilon\circ\rho)_j=\varepsilon_{\rho(j)},
\]
one has
\[
F_{\varepsilon\circ\rho}(w)=F_\varepsilon(w),
\qquad
\psi_{\rho(E)}(\varepsilon)=\psi_E(\varepsilon\circ\rho).
\]
For a given \(E\), choose such a permutation with
\(\rho(E_{|E|})=E\), which is possible because both sets have cardinality \(|E|\). Reindexing the sign sum therefore yields
\[
a_{|E|}(w)
=2^{-d}\sum_{\varepsilon\in\{-1,1\}^d}
\frac{F_\varepsilon(w)}{g_d(w)}\psi_E(\varepsilon).
\]

The finite character kernel satisfies
\[
\sum_{E\subseteq\{1,\ldots,d\}}
\psi_E(\delta)\psi_E(\varepsilon)
=
\prod_{j=1}^d(1+\delta_j\varepsilon_j)
=
\begin{cases}
2^d,&\delta=\varepsilon,\\
0,&\delta\ne\varepsilon,
\end{cases}
\qquad \varepsilon,\delta\in\{-1,1\}^d.
\]
The first equality expands the product over all subsets. If the sign vectors agree, every factor equals two; otherwise a coordinate where they differ supplies a zero factor. Substituting the coefficient formula and using this kernel gives
\[
\begin{aligned}
\sum_{E\subseteq\{1,\ldots,d\}}
a_{|E|}(w)\psi_E(\varepsilon)
&=
2^{-d}\sum_{\delta\in\{-1,1\}^d}
\frac{F_\delta(w)}{g_d(w)}
\sum_{E\subseteq\{1,\ldots,d\}}
\psi_E(\delta)\psi_E(\varepsilon)\\
&=\frac{F_\varepsilon(w)}{g_d(w)}.
\end{aligned}
\]
% lean: CausalSmith.Experimentation.ThinnedgraphAdditiveRiskFrontier.rowDensity_walsh_expansion

\item Each \(F_\varepsilon\) is a translate of the baseline probability density, so
\[
F_\varepsilon(w)\ge0,
\qquad
\int_{\mathbb R}F_\varepsilon(w)\,dw=1.
\]
Consequently \(g_d(w)\ge0\), and \(g_d(w)=0\) forces every \(F_\varepsilon(w)=0\). Thus cancellation of the denominator, including the zero-density convention, gives the identity
\[
a_s(w)g_d(w)
=
2^{-d}\sum_{\varepsilon\in\{-1,1\}^d}
F_\varepsilon(w)\psi_{E_s}(\varepsilon),
\qquad w\in\mathbb R,\quad s\in\mathbb Z_{\ge0}.
\]
For \(s\ge1\), the assumptions \(d\ge1\) and \(s\ge1\) ensure that \(E_s\) contains coordinate one, including when \(s>d\). Summing signs coordinate by coordinate gives
\[
\sum_{\varepsilon\in\{-1,1\}^d}\psi_{E_s}(\varepsilon)
=
\prod_{j=1}^d
\left[
\sum_{\varepsilon_j\in\{-1,1\}}
\begin{cases}
\varepsilon_j,&j\in E_s,\\
1,&j\notin E_s
\end{cases}
\right]
=0,
\]
because the factor for coordinate one is \((-1)+1=0\). Integrating the preceding cancellation identity, with the finite sum exchanged with the integral, now yields
\[
\int_{\mathbb R}a_s(w)g_d(w)\,dw
=
2^{-d}\sum_{\varepsilon\in\{-1,1\}^d}
\psi_{E_s}(\varepsilon)
\int_{\mathbb R}F_\varepsilon(w)\,dw
=0.
\]
% lean: CausalSmith.Experimentation.ThinnedgraphAdditiveRiskFrontier.walshCoeff_centered

\item Nonnegativity of \(g_d\) gives
\[
\gamma_s=\int_{\mathbb R}a_s(w)^2g_d(w)\,dw\ge0
\qquad(1\le s\le d).
\]
Since \(s\ge1\) on this range,
\[
0\le \binom ds\gamma_s
\le s\binom ds\gamma_s.
\]
Summing proves
\[
0\le\eta\le\eta_1.
\]
% lean: CausalSmith.Experimentation.ThinnedgraphAdditiveRiskFrontier.eta_nonneg_le_etaOne

\item First consider \(h=0\). Then \(F_\varepsilon(w)=f(w)\) for every sign vector. At a point where \(g_d(w)=0\), the coefficient convention gives \(a_s(w)=0\). At any other point, for every \(s\ge1\), the character cancellation from the preceding argument gives
\[
a_s(w)
=
2^{-d}\frac{f(w)}{g_d(w)}
\sum_{\varepsilon\in\{-1,1\}^d}\psi_{E_s}(\varepsilon)
=0.
\]
Hence every \(\gamma_s\), \(1\le s\le d\), vanishes, and
\[
\eta_1=0=\frac{12\pi^2h^2}{d}.
\]
% lean: CausalSmith.Experimentation.ThinnedgraphAdditiveRiskFrontier.etaOne_zero_amplitude

\item Suppose henceforth that \(h\ne0\). Define the coordinate reversal and its averaged density energy by
\[
(\varepsilon^{(j)})_i=
\begin{cases}
-\varepsilon_j,&i=j,\\
\varepsilon_i,&i\ne j,
\end{cases}
\qquad
\varepsilon\in\{-1,1\}^d,\quad i,j\in\{1,\ldots,d\},
\]
and
\[
\mathcal D_j(w)=
\begin{cases}
2^{-d}\displaystyle\sum_{\varepsilon\in\{-1,1\}^d}
\frac{(F_\varepsilon(w)-F_{\varepsilon^{(j)}}(w))^2}{g_d(w)},
&g_d(w)>0,\\
0,&g_d(w)=0,
\end{cases}
\qquad w\in\mathbb R,\quad 1\le j\le d.
\]

To derive the weighted identity, let
\[
v_{\mathrm{cube}}:\{-1,1\}^d\longrightarrow\mathbb R,
\qquad
\widehat v_{\mathrm{cube}}(E)
=
2^{-d}\sum_{\varepsilon\in\{-1,1\}^d}
v_{\mathrm{cube}}(\varepsilon)\psi_E(\varepsilon),
\quad E\subseteq\{1,\ldots,d\}.
\]
Expanding the squared coefficients and using the character kernel established above proves finite Parseval:
\[
\begin{aligned}
\sum_{E\subseteq\{1,\ldots,d\}}
\widehat v_{\mathrm{cube}}(E)^2
&=
2^{-2d}\sum_{\varepsilon,\delta\in\{-1,1\}^d}
v_{\mathrm{cube}}(\varepsilon)v_{\mathrm{cube}}(\delta)
\sum_{E\subseteq\{1,\ldots,d\}}
\psi_E(\varepsilon)\psi_E(\delta)\\
&=
2^{-d}\sum_{\varepsilon\in\{-1,1\}^d}
v_{\mathrm{cube}}(\varepsilon)^2.
\end{aligned}
\]
A coordinate reversal is an involution, and
\[
\psi_E(\varepsilon^{(j)})
=
\begin{cases}
-\psi_E(\varepsilon),&j\in E,\\
\psi_E(\varepsilon),&j\notin E.
\end{cases}
\]
Reindexing by this involution therefore gives
\[
2^{-d}\sum_{\varepsilon\in\{-1,1\}^d}
\bigl(v_{\mathrm{cube}}(\varepsilon)
-v_{\mathrm{cube}}(\varepsilon^{(j)})\bigr)\psi_E(\varepsilon)
=
\begin{cases}
2\widehat v_{\mathrm{cube}}(E),&j\in E,\\
0,&j\notin E.
\end{cases}
\]
Applying finite Parseval to the coordinate difference yields
\[
\sum_{\substack{E\subseteq\{1,\ldots,d\}\\j\in E}}
\widehat v_{\mathrm{cube}}(E)^2
=
\frac14\,2^{-d}\sum_{\varepsilon\in\{-1,1\}^d}
\bigl(v_{\mathrm{cube}}(\varepsilon)
-v_{\mathrm{cube}}(\varepsilon^{(j)})\bigr)^2.
\]
Summing over \(j\) counts each subset exactly \(|E|\) times, so
\[
\sum_{E\subseteq\{1,\ldots,d\}}
|E|\widehat v_{\mathrm{cube}}(E)^2
=
\frac14\sum_{j=1}^d2^{-d}
\sum_{\varepsilon\in\{-1,1\}^d}
\bigl(v_{\mathrm{cube}}(\varepsilon)
-v_{\mathrm{cube}}(\varepsilon^{(j)})\bigr)^2.
\]

For \(g_d(w)>0\), apply this identity with
\[
v_{\mathrm{cube}}(\varepsilon)=\frac{F_\varepsilon(w)}{g_d(w)}.
\]
Its subset coefficients are \(a_{|E|}(w)\). Grouping subsets by cardinality and multiplying by \(g_d(w)\) gives
\[
\sum_{s=0}^d s\binom ds a_s(w)^2g_d(w)
=
\frac14\sum_{j=1}^d\mathcal D_j(w).
\]
At \(g_d(w)=0\), both sides vanish by their conventions, so this identity holds for every \(w\).

The coefficient densities are integrable: the triangle inequality gives, at positive \(g_d(w)\),
\[
|a_s(w)|
\le
2^{-d}\sum_{\varepsilon\in\{-1,1\}^d}
\frac{F_\varepsilon(w)}{g_d(w)}
=1,
\]
and the zero-density convention preserves this bound. Thus
\[
0\le a_s(w)^2g_d(w)\le g_d(w),
\qquad
\int_{\mathbb R}g_d(w)\,dw=1.
\]
The flip energies are integrable by the bounded compact-support envelope established in the next step. Integrating the pointwise identity and discarding its zero \(s=0\) term therefore gives
\[
\eta_1=\frac14\sum_{j=1}^d\int_{\mathbb R}\mathcal D_j(w)\,dw.
\]
% lean: CausalSmith.Experimentation.ThinnedgraphAdditiveRiskFrontier.walsh_weighted_integral_parseval

\item By \cref{obj:lem:baseline-translation-affinity}, the square root of the baseline density is globally \(4\pi\)-Lipschitz. Hence, for \(w_{\mathrm a},w_{\mathrm b}\in\mathbb R\),
\[
\bigl(\sqrt{f(w_{\mathrm a})}-\sqrt{f(w_{\mathrm b})}\bigr)^2
\le16\pi^2(w_{\mathrm a}-w_{\mathrm b})^2.
\]
Also,
\[
\bigl(\sqrt{f(w_{\mathrm a})}+\sqrt{f(w_{\mathrm b})}\bigr)^2
\le2\bigl(f(w_{\mathrm a})+f(w_{\mathrm b})\bigr),
\]
as follows by expanding the nonnegative square of their difference. Factoring the difference of the densities and multiplying these bounds gives
\[
\begin{aligned}
\bigl(f(w_{\mathrm a})-f(w_{\mathrm b})\bigr)^2
&=
\bigl(\sqrt{f(w_{\mathrm a})}-\sqrt{f(w_{\mathrm b})}\bigr)^2
\bigl(\sqrt{f(w_{\mathrm a})}+\sqrt{f(w_{\mathrm b})}\bigr)^2\\
&\le
32\pi^2(w_{\mathrm a}-w_{\mathrm b})^2
\bigl(f(w_{\mathrm a})+f(w_{\mathrm b})\bigr).
\end{aligned}
\]

Define the translation amounts by
\[
\zeta_\varepsilon=\frac{h}{2d}\sum_{i=1}^d\varepsilon_i,
\qquad \varepsilon\in\{-1,1\}^d.
\]
Reversing coordinate \(j\) gives
\[
\sum_{i=1}^d(\varepsilon^{(j)})_i
=\sum_{i=1}^d\varepsilon_i-2\varepsilon_j,
\qquad
\zeta_{\varepsilon^{(j)}}=\zeta_\varepsilon-\frac hd\varepsilon_j.
\]
Consequently the two translated arguments have squared difference \(h^2/d^2\), and the preceding density bound becomes
\[
\bigl(F_\varepsilon(w)-F_{\varepsilon^{(j)}}(w)\bigr)^2
\le
\frac{32\pi^2h^2}{d^2}
\bigl(F_\varepsilon(w)+F_{\varepsilon^{(j)}}(w)\bigr).
\]
Since the reversal permutes the sign cube,
\[
2^{-d}\sum_{\varepsilon\in\{-1,1\}^d}
\bigl(F_\varepsilon(w)+F_{\varepsilon^{(j)}}(w)\bigr)
=2g_d(w).
\]
Summing the density bound and dividing by \(g_d(w)>0\) proves
\[
0\le\mathcal D_j(w)\le\frac{64\pi^2h^2}{d^2}.
\]
At \(g_d(w)=0\), the same inequality follows from \(\mathcal D_j(w)=0\).

For the support bound, define
\[
I_h=\left[-\frac14-\frac h2,\frac14+\frac h2\right].
\]
The triangle inequality gives
\[
|\zeta_\varepsilon|
\le\frac{h}{2d}\sum_{i=1}^d|\varepsilon_i|
=\frac h2.
\]
Thus, if \(w\notin I_h\),
\[
|w-\zeta_\varepsilon|
\ge |w|-|\zeta_\varepsilon|>\frac14,
\]
so the support of \(f\) specified in
\cref{obj:lem:baseline-translation-affinity} implies
\(F_\varepsilon(w)=0\) for every \(\varepsilon\), and hence \(g_d(w)=0\).
Writing the interval indicator explicitly as
\[
\mathbf1_{I_h}(w)=
\begin{cases}
1,&w\in I_h,\\
0,&w\notin I_h,
\end{cases}
\]
we obtain the integrable envelope
\[
0\le\mathcal D_j(w)
\le\frac{64\pi^2h^2}{d^2}\mathbf1_{I_h}(w).
\]
The interval has length \(1/2+h\le3/4\). Therefore
\[
\int_{\mathbb R}\mathcal D_j(w)\,dw
\le
\left(\frac12+h\right)\frac{64\pi^2h^2}{d^2}
\le\frac{48\pi^2h^2}{d^2}.
\]
Finally, summing this bound over the \(d\) coordinates in the integrated identity gives
\[
\eta_1
=\frac14\sum_{j=1}^d\int_{\mathbb R}\mathcal D_j(w)\,dw
\le\frac d4\,\frac{48\pi^2h^2}{d^2}
=\frac{12\pi^2h^2}{d}.
\]
Together with the preceding steps, this proves all the assertions.
% lean: CausalSmith.Experimentation.ThinnedgraphAdditiveRiskFrontier.integrated_flip_sum_le
\end{enumerate}
\end{proof}

\Cref{obj:lem:walsh-channel-energy} supplies an exact expansion of the normalized response density on the positive-density region of the mixture. Its constant coefficient is one, and each nonconstant coefficient has mean zero under that mixture. The aggregate bound quantifies the squared size of all nonconstant terms, including their order weights, at scale \(h^2/d\). This gives a single measure of response information that can be used across the possible hidden-source allocations.

\subsection{Conditional reference laws and hidden allocation}

The complete-record comparison combines response information with the source associations revealed by the audit. Conditional on the retained graph, the reference law preserves the actual assignment marginal and replaces the distinct block responses by independent uniform sign-mixture responses. The required capacities and conditional discrepancies are defined together.



The reference law in \cref{obj:lem:hidden-allocation-contraction} keeps every assignment coordinate in the comparison. The reconstruction property in \cref{obj:lem:block-law} connects the distinct responses to the full outcome record, including their repeated recipient coordinates. The remaining-capacity event measures how much source allocation is still hidden after observing the graph. Under a smallness condition on the total Walsh energy, that contraction result bounds the conditional discrepancy averaged over the complete graph.

\begin{lemmav}{lem:hidden-allocation-contraction}[Hidden allocation contraction bound]
In the complete block experiment, let \(n\) be the population size, \(B\) the number of blocks, \(d\) the number of sources per block, and \(m=Bd\). Suppose:
\begin{itemize}
\item \textbf{(Block sizes.)} \(n,B,d\) are integers satisfying \(n\ge4\), \(B\ge1\), \(d\ge1\), and \(2Bd\le n\).
\item \textbf{(Amplitude and audit.)} The amplitude satisfies \(0\le h\le1/4\). The retention probability satisfies \(0\le q\le1\); the audit marks on all ordered off-diagonal pairs are independent Bernoulli variables with success probability \(q\), independent of the assignment.
\item \textbf{(Assignment.)} The assignment law satisfies \cref{obj:ass:assignment}.
\item \textbf{(Channel smallness.)} Let \(\gamma_s\) be the integrated squared Walsh coefficient of order \(s\) for the block-response channel at amplitude \(h\), and define its total and order-weighted nonconstant energies by
\[
\eta=\sum_{s=1}^{d}\binom{d}{s}\gamma_s,
\qquad
\eta_1=\sum_{s=1}^{d}s\binom{d}{s}\gamma_s.
\]
These satisfy \(4e\eta<1\).
\end{itemize}
Let \(H\) be the complete labeled retained graph. For each block, let \(u_\ell\) be its remaining hidden-source capacity, and define the total remaining capacity and source-association revelation probability by
\[
M=\sum_{\ell=1}^{B}u_\ell,
\qquad
p=1-(1-q)^d.
\]
Conditional on \(H\), let \(Q_H\) have the actual assignment marginal and independent distinct block responses with the uniform sign-mixture response density \(g_d\). Let \(L_+\) be the density ratio of the positive-prior conditional law \(P_{+1,h}(Z,y\mid H)\) relative to \(Q_H\), and define
\[
\chi_H^2=\mathbb E_{Q_H}\!\left[(L_+-1)^2\right].
\]
Then
\[
\mathbb E_H\!\left[\mathbf1\{M\ge m/4\}\chi_H^2\right]
\le
\frac{\exp(eBp\eta_1)}{1-4e\eta}-1.
\]
The expectation is taken under the actual marginal law of the complete retained graph, including its labels and detailed retained-edge subsets.
\end{lemmav}

\begin{proof}[Proof of \cref{obj:lem:hidden-allocation-contraction}]
We use the convention \(0^0=1\) and interpret empty products as one.

\begin{enumerate}
\item First establish the scalar series estimate that will close the argument. Define
\[
e:=\exp(1),
\qquad
\rho_{\nu_{\mathrm{rev}},\nu_{\mathrm{hid}}}
:=
\min\left\{1,\frac{4(\nu_{\mathrm{rev}}+\nu_{\mathrm{hid}})}{B}\right\}
\quad
(\nu_{\mathrm{rev}},\nu_{\mathrm{hid}}\in\mathbb N_0),
\]
and
\[
\mathfrak w_{\nu_{\mathrm{rev}},\nu_{\mathrm{hid}}}
:=
\binom B{\nu_{\mathrm{rev}}}
\binom{B-\nu_{\mathrm{rev}}}{\nu_{\mathrm{hid}}}
\rho_{\nu_{\mathrm{rev}},\nu_{\mathrm{hid}}}^{\nu_{\mathrm{hid}}}
(p\eta_1)^{\nu_{\mathrm{rev}}}\eta^{\nu_{\mathrm{hid}}}
\quad
(0\le\nu_{\mathrm{rev}},\nu_{\mathrm{hid}}\le B).
\]
By \cref{obj:lem:walsh-channel-energy},
\(0\le\eta\le\eta_1\). Also \(0\le(1-q)^d\le1\), so \(0\le p\le1\).
Consequently all these terms are nonnegative.

Define the nonconstant sum by
\[
\mathfrak S_{\mathrm{nc}}
:=
\sum_{\substack{0\le\nu_{\mathrm{rev}},\nu_{\mathrm{hid}}\le B\\
\nu_{\mathrm{rev}}+\nu_{\mathrm{hid}}>0}}
\mathfrak w_{\nu_{\mathrm{rev}},\nu_{\mathrm{hid}}}
=
\sum_{\nu_{\mathrm{rev}}=0}^B
\sum_{\nu_{\mathrm{hid}}=0}^B
\mathfrak w_{\nu_{\mathrm{rev}},\nu_{\mathrm{hid}}}-1.
\]
The equality follows because \(\mathfrak w_{0,0}=1\).

The falling-factorial formula for binomial coefficients gives
\[
\binom B{\nu_{\mathrm{rev}}}
\le\frac{B^{\nu_{\mathrm{rev}}}}{\nu_{\mathrm{rev}}!},
\qquad
\binom{B-\nu_{\mathrm{rev}}}{\nu_{\mathrm{hid}}}
\le\frac{B^{\nu_{\mathrm{hid}}}}{\nu_{\mathrm{hid}}!}.
\]
Using \(B>0\) and the definition of \(\rho\), we obtain
\[
\begin{aligned}
&\binom B{\nu_{\mathrm{rev}}}
\binom{B-\nu_{\mathrm{rev}}}{\nu_{\mathrm{hid}}}
\rho_{\nu_{\mathrm{rev}},\nu_{\mathrm{hid}}}^{\nu_{\mathrm{hid}}}
\\
&\qquad\le
\frac{B^{\nu_{\mathrm{rev}}}}{\nu_{\mathrm{rev}}!}
4^{\nu_{\mathrm{hid}}}
\frac{(\nu_{\mathrm{rev}}+\nu_{\mathrm{hid}})^{\nu_{\mathrm{hid}}}}
{\nu_{\mathrm{hid}}!}
\\
&\qquad\le
\frac{B^{\nu_{\mathrm{rev}}}}{\nu_{\mathrm{rev}}!}
4^{\nu_{\mathrm{hid}}}
\exp(\nu_{\mathrm{rev}}+\nu_{\mathrm{hid}})
=
\frac{(eB)^{\nu_{\mathrm{rev}}}}{\nu_{\mathrm{rev}}!}
(4e)^{\nu_{\mathrm{hid}}}.
\end{aligned}
\]
The second inequality bounds one nonnegative term of the exponential series by its sum. It includes \(\nu_{\mathrm{hid}}=0\), with the stated power convention. Thus
\[
\begin{aligned}
\mathfrak S_{\mathrm{nc}}+1
&\le
\left(
\sum_{\nu_{\mathrm{rev}}=0}^B
\frac{(eBp\eta_1)^{\nu_{\mathrm{rev}}}}{\nu_{\mathrm{rev}}!}
\right)
\left(
\sum_{\nu_{\mathrm{hid}}=0}^B
(4e\eta)^{\nu_{\mathrm{hid}}}
\right)
\\
&\le
\exp(eBp\eta_1)\sum_{\nu_{\mathrm{hid}}=0}^{\infty}
(4e\eta)^{\nu_{\mathrm{hid}}}
=
\frac{\exp(eBp\eta_1)}{1-4e\eta}.
\end{aligned}
\]
Here \(0\le4e\eta<1\) justifies the geometric series. Therefore
\[
\mathfrak S_{\mathrm{nc}}
\le
\frac{\exp(eBp\eta_1)}{1-4e\eta}-1.
\]
% lean: CausalSmith.Experimentation.ThinnedgraphAdditiveRiskFrontier.contraction_nonconstant_sum_bound

\item We next identify the design and its conditional reference. Define
\[
\mu_Z:=\operatorname{Law}(Z),
\qquad
\mu_{\mathrm{ass}}
:=
\bigotimes_{j\in V}\operatorname{Bernoulli}(1/2),
\qquad
\mu_{\mathrm{aud}}
:=
\bigotimes_{\substack{(j,i)\in V^2\\j\ne i}}
\operatorname{Bernoulli}(q).
\]
The audit measure is a probability measure because \(q\in[0,1]\).
The first marginal of \(\mu_Z\otimes\mu_{\mathrm{aud}}\) is therefore
\(\mu_Z\). The assignment condition in
\cref{obj:ass:assignment}, together with assignment and audit independence, gives
\[
\mu_Z=\mu_{\mathrm{ass}},
\qquad
\operatorname{Law}(Z,W)
=
\mu_{\mathrm{ass}}\otimes\mu_{\mathrm{aud}}.
\]
The source partition is drawn independently of this design, and the retained graph is a function of the partition and audit marks. Hence
\[
\operatorname{Law}(H,Z)
=
\operatorname{Law}(H)\otimes\mu_{\mathrm{ass}}.
\]
In particular, the assignment signs remain independent uniform signs conditional on every graph atom of positive mass.

Define the response reference measure on \(\mathbb R^B\) by
\[
\nu_{\mathrm{resp}}(dy)
:=
\prod_{\ell=1}^B g_d(y_\ell)\,dy_\ell.
\]
Each \(g_d\) is a probability density, being a finite mixture of translates of the baseline density in \cref{obj:lem:block-law}. Thus
\[
Q_H=\mu_{\mathrm{ass}}\otimes\nu_{\mathrm{resp}}.
\]
Every realized retained graph is compatible with its sampled source partition. Positive graph atoms therefore admit compatible partitions, and they exhaust the graph marginal almost surely. All graphwise constructions below concern these atoms; auxiliary graph functions are extended by zero at null atoms.
% lean: CausalSmith.Experimentation.ThinnedgraphAdditiveRiskFrontier.hidden_allocation_contraction

\item Fix such a graph \(H\). Let the fixed source population, the revealed source sets, and the hidden source set be
\[
\mathcal S_{\mathrm{src}}:=\bigcup_{\ell=1}^B S_\ell,
\qquad
R_\ell:=\{j\in\mathcal S_{\mathrm{src}}:
\exists i\in C_\ell,\ (j,i)\in H\},
\qquad
\mathcal D:=\mathcal S_{\mathrm{src}}\setminus\bigcup_{\ell=1}^B R_\ell.
\]
The union defining \(\mathcal S_{\mathrm{src}}\) is the same fixed set for every source partition. Compatibility makes the sets \(R_\ell\) pairwise disjoint. Write
\[
r_\ell:=|R_\ell|,
\qquad
u_\ell=d-r_\ell,
\qquad
M=\sum_{\ell=1}^B u_\ell=|\mathcal D|.
\]
For \(z\in\{0,1\}^V\), define
\[
X_j(z):=2z_j-1,
\qquad
X_j:=X_j(Z),
\qquad
A_\ell(z):=\sum_{j\in R_\ell}X_j(z),
\qquad
K(z):=\sum_{j\in\mathcal D}z_j.
\]

The family of hidden-label completions is
\[
\mathfrak C_H
:=
\left\{
(\mathcal D_1,\ldots,\mathcal D_B):
\mathcal D=\bigsqcup_{\ell=1}^B\mathcal D_\ell,\
|\mathcal D_\ell|=u_\ell
\right\}.
\]
It is nonempty and has cardinality
\[
|\mathfrak C_H|=\frac{M!}{\prod_{\ell=1}^B u_\ell!}.
\]
For integers \(0\le k_\ell\le u_\ell\) with
\(\sum_\ell k_\ell=K(z)\), counting the allocations of treated and untreated hidden labels separately gives
\[
\begin{aligned}
&\frac{
\left|
\left\{
(\mathcal D_\ell)\in\mathfrak C_H:
|\mathcal D_\ell\cap\{j:z_j=1\}|=k_\ell
\text{ for every }\ell
\right\}
\right|
}{|\mathfrak C_H|}
\\
&\qquad=
\frac{
\displaystyle
\frac{K(z)!}{\prod_\ell k_\ell!}
\frac{(M-K(z))!}{\prod_\ell(u_\ell-k_\ell)!}
}{
\displaystyle
\frac{M!}{\prod_\ell u_\ell!}
}
=
\frac{\prod_\ell\binom{u_\ell}{k_\ell}}{\binom M{K(z)}}.
\end{aligned}
\]
The count is zero when the total-degree condition fails. Since \(0\le K(z)\le M\), its denominator is positive, including \(M=K(z)=0\).

Expanding the coefficient-extraction formula in
\cref{obj:lem:block-law} and applying this count gives the exact identity
\[
\lambda_{+1,h}(y\mid H,z)
=
\frac1{|\mathfrak C_H|}
\sum_{(\mathcal D_\ell)\in\mathfrak C_H}
\prod_{\ell=1}^B
f\left(
y_\ell-\frac h{2d}
\left(A_\ell(z)+\sum_{j\in\mathcal D_\ell}X_j(z)\right)
\right).
\]
Thus the conditional density is the uniform-completion average of the original labeled row densities.

If \(g_d(y_\ell)=0\), every translated density contributing to its nonnegative finite mixture is zero. The last display then gives
\(\lambda_{+1,h}(y\mid H,z)=0\). Consequently a version of the likelihood is
\[
L_+(z,y)
=
\begin{cases}
\displaystyle
\frac{\lambda_{+1,h}(y\mid H,z)}
{\prod_{\ell=1}^B g_d(y_\ell)},
&\prod_{\ell=1}^B g_d(y_\ell)>0,\\[6pt]
0,&\prod_{\ell=1}^B g_d(y_\ell)=0.
\end{cases}
\]
The response reference assigns full measure to the first case. The conditional density identification and the product form of \(Q_H\) therefore yield
\[
\chi_H^2
=
\int_{\mathbb R^B}
\int_{\{0,1\}^V}
\bigl(L_+(z,y)-1\bigr)^2
\,\mu_{\mathrm{ass}}(dz)\,\nu_{\mathrm{resp}}(dy).
\]
Finite channel expansions are bounded, so the integral rearrangements below are justified. Multiplying this identity by the remaining-capacity indicator preserves it almost surely under the actual graph marginal.
% lean: CausalSmith.Experimentation.ThinnedgraphAdditiveRiskFrontier.hiddenChiSq_good_integral_eq_actualDensityEnergy

\item We now express this density energy through labeled coefficients. For
\(0\le k\le M\), define
\[
\Psi_k(X_{\mathcal D})
:=
\binom Mk^{-1}
\sum_{\substack{J_{\mathrm{hid}}\subseteq\mathcal D\\
|J_{\mathrm{hid}}|=k}}
\prod_{j\in J_{\mathrm{hid}}}X_j,
\qquad
X_{\mathcal D}:=(X_j)_{j\in\mathcal D}.
\]
Every denominator is positive; in particular \(\Psi_0=1\).

For a hidden-degree vector, define
\[
\boldsymbol k:=(k_1,\ldots,k_B)
\in\prod_{\ell=1}^B\{0,\ldots,u_\ell\},
\qquad
|\boldsymbol k|:=\sum_{\ell=1}^B k_\ell,
\qquad
\omega(\boldsymbol k):=\prod_{\ell=1}^B\binom{u_\ell}{k_\ell}.
\]
For any fixed completion \((\mathcal D_\ell)\), averaging over all permutations of the \(M\) hidden labels gives
\[
\begin{aligned}
&\frac1{M!}\sum_{\pi\in\operatorname{Sym}(\mathcal D)}
\sum_{\substack{
J_{\mathrm{hid},\ell}\subseteq\mathcal D_\ell\\
|J_{\mathrm{hid},\ell}|=k_\ell\ \text{for every }\ell}}
\prod_{\ell=1}^B
\prod_{j\in J_{\mathrm{hid},\ell}}X_{\pi(j)}
\\
&\qquad=
\omega(\boldsymbol k)\,
\Psi_{|\boldsymbol k|}(X_{\mathcal D}).
\end{aligned}
\]
Indeed, there are \(\omega(\boldsymbol k)\) selections. For each selection, its union has size \(|\boldsymbol k|\), and a uniform permutation maps that union uniformly onto subsets of this size. The identity also holds for \(M=0\), when the permutation and selection families are singletons. Uniform completion averaging is invariant under these permutations.

For each revealed subset \(J\subseteq\bigcup_\ell R_\ell\), define
\[
r_{J,\ell}:=|J\cap R_\ell|,
\]
and, for \(0\le k\le M\) and \(y\in\mathbb R^B\), define
\[
\mathcal C_{J,k}(y)
:=
\sum_{\substack{
\boldsymbol k\in\prod_\ell\{0,\ldots,u_\ell\}\\
|\boldsymbol k|=k}}
\omega(\boldsymbol k)
\prod_{\ell=1}^B a_{r_{J,\ell}+k_\ell}(y_\ell).
\]
The orders in this formula lie between zero and \(d\).
Applying the row expansion in \cref{obj:lem:walsh-channel-energy} to the completion average, and then grouping by total hidden degree, gives
\[
L_+(Z,y)
=
\sum_{J\subseteq\bigcup_\ell R_\ell}
\left(\prod_{j\in J}X_j\right)
\sum_{k=0}^M
\Psi_k(X_{\mathcal D})\,\mathcal C_{J,k}(y)
\]
for \(\nu_{\mathrm{resp}}\)-almost every \(y\). The disjointness of the revealed row sets identifies a tuple of rowwise revealed subsets with its union \(J\).

Independent uniform signs satisfy
\[
\mathbb E\!\left[
\Psi_k(X_{\mathcal D})\Psi_{k'}(X_{\mathcal D})
\mid H
\right]
=
\frac{\mathbf1\{k=k'\}}{\binom Mk}
\qquad(0\le k,k'\le M).
\]
To see this, expand the two subset sums. The expected product is zero unless the two selected subsets coincide, since a sign appearing exactly once has mean zero. There are \(\binom Mk\) coincident pairs when \(k=k'\). The same argument for revealed subsets, and independence of the revealed and hidden signs, gives
\[
\begin{aligned}
&\mathbb E\!\left[
\left(\prod_{j\in J}X_j\right)\Psi_k(X_{\mathcal D})
\left(\prod_{j\in J'}X_j\right)\Psi_{k'}(X_{\mathcal D})
\mid H
\right]
\\
&\qquad=
\frac{\mathbf1\{J=J',\ k=k'\}}{\binom Mk}.
\end{aligned}
\]
Treatment coordinates outside the source population contribute a factor one.

For a square-integrable function
\(\varpi:\mathbb R^B\to\mathbb R\), write
\[
\|\varpi\|_{\mathrm{resp}}^2
:=
\int_{\mathbb R^B}\varpi(y)^2\,\nu_{\mathrm{resp}}(dy).
\]
The zero-degree coefficient is
\[
\mathcal C_{\varnothing,0}(y)=\prod_{\ell=1}^B a_0(y_\ell)=1
\quad\text{for }\nu_{\mathrm{resp}}\text{-almost every }y.
\]
Subtracting this coefficient, applying sign orthogonality, and integrating over \(y\) proves the exact labeled-energy identity
\[
\chi_H^2
=
\sum_{\substack{
J\subseteq\bigcup_\ell R_\ell,\ 0\le k\le M\\
(J,k)\ne(\varnothing,0)}}
\frac{\|\mathcal C_{J,k}\|_{\mathrm{resp}}^2}{\binom Mk}.
\]
% lean: CausalSmith.Experimentation.ThinnedgraphAdditiveRiskFrontier.actualDensityEnergy_eq_actualLabelEnergy

\item Separate these coefficients by their response supports. Define
\[
\mathcal B_J:=\{\ell:r_{J,\ell}>0\}.
\]
For \(\mathcal E\subseteq\{1,\ldots,B\}\setminus\mathcal B_J\) and
\(0\le k\le M\), let
\[
\mathscr K_{J,\mathcal E,k}
:=
\left\{
\boldsymbol k\in\prod_\ell\{0,\ldots,u_\ell\}:
|\boldsymbol k|=k,\
k_\ell>0\ (\ell\in\mathcal E),\
k_\ell=0\ (\ell\notin\mathcal B_J\cup\mathcal E)
\right\},
\]
and define
\[
\mathcal C_{J,k,\mathcal E}(y)
:=
\sum_{\boldsymbol k\in\mathscr K_{J,\mathcal E,k}}
\omega(\boldsymbol k)
\prod_{\ell=1}^B a_{r_{J,\ell}+k_\ell}(y_\ell).
\]
Every hidden-degree vector belongs to exactly one such family, so
\[
\mathcal C_{J,k}
=
\sum_{\mathcal E\subseteq\{1,\ldots,B\}\setminus\mathcal B_J}
\mathcal C_{J,k,\mathcal E}.
\]

Each summand has positive coefficient order precisely on
\(\mathcal B_J\cup\mathcal E\). If \(\mathcal E\ne\mathcal E'\), one row belongs to exactly one of these supports. At that row the cross product has one positive-order coefficient and one factor \(a_0=1\) almost everywhere. The zero-mean assertion of
\cref{obj:lem:walsh-channel-energy}, together with the product response measure, gives
\[
\int
\mathcal C_{J,k,\mathcal E}(y)
\mathcal C_{J,k,\mathcal E'}(y)\,
\nu_{\mathrm{resp}}(dy)=0.
\]
It follows that
\[
\|\mathcal C_{J,k}\|_{\mathrm{resp}}^2
=
\sum_{\mathcal E\subseteq\{1,\ldots,B\}\setminus\mathcal B_J}
\|\mathcal C_{J,k,\mathcal E}\|_{\mathrm{resp}}^2.
\]
% lean: CausalSmith.Experimentation.ThinnedgraphAdditiveRiskFrontier.actualLabelWalshCoeff_support_energy

\item Suppose now that \(M\ge Bd/4\). Since \(B,d\ge1\), this implies \(M>0\).
For fixed \(J,\mathcal E\), define
\[
\nu_{\mathrm{rev}}:=|\mathcal B_J|,
\qquad
\nu_{\mathrm{hid}}:=|\mathcal E|,
\qquad
U_{\mathrm{act}}:=\sum_{\ell\in\mathcal B_J\cup\mathcal E}u_\ell.
\]
These quantities satisfy
\[
U_{\mathrm{act}}\le M,
\qquad
U_{\mathrm{act}}\le d(\nu_{\mathrm{rev}}+\nu_{\mathrm{hid}}),
\qquad
\frac{U_{\mathrm{act}}}{M}
\le\rho_{\nu_{\mathrm{rev}},\nu_{\mathrm{hid}}}\le1.
\]
The first two inequalities use the nonnegative capacities and \(u_\ell\le d\); the third uses both \(U_{\mathrm{act}}\le M\) and \(M\ge Bd/4\).

Every vector in \(\mathscr K_{J,\mathcal E,k}\) has
\[
k=\sum_{\ell=1}^B k_\ell
\ge\sum_{\ell\in\mathcal E}k_\ell
\ge\nu_{\mathrm{hid}}.
\]
For \(0\le k\le U_{\mathrm{act}}\), the falling-factorial identity gives
\[
\frac{\binom{U_{\mathrm{act}}}{k}}{\binom Mk}
=
\prod_{v=0}^{k-1}
\frac{U_{\mathrm{act}}-v}{M-v}
\le
\left(\frac{U_{\mathrm{act}}}{M}\right)^k.
\]
Every denominator in this product is positive. At \(k=0\), both sides are one. If \(U_{\mathrm{act}}<k\le M\), the numerator binomial coefficient is zero. Consequently, whenever \(0\le k\le M\) and \(k\ge\nu_{\mathrm{hid}}\),
\[
\frac{\binom{U_{\mathrm{act}}}{k}}{\binom Mk}
\le
\rho_{\nu_{\mathrm{rev}},\nu_{\mathrm{hid}}}^{\nu_{\mathrm{hid}}}.
\]
Here the exponent decreases from \(k\) to \(\nu_{\mathrm{hid}}\) because the capacity fraction lies in \([0,1]\).
% lean: CausalSmith.Experimentation.ThinnedgraphAdditiveRiskFrontier.hidden_choose_ratio_good_event

The allocation weights satisfy
\[
\sum_{\boldsymbol k\in\mathscr K_{J,\mathcal E,k}}
\omega(\boldsymbol k)
\le
[x^k]\prod_{\ell\in\mathcal B_J\cup\mathcal E}(1+x)^{u_\ell}
=
\binom{U_{\mathrm{act}}}{k},
\]
because the restricted family is contained in the full family selecting \(k\) slots from these active rows.

Weighted Cauchy--Schwarz gives, pointwise in \(y\),
\[
\begin{aligned}
\mathcal C_{J,k,\mathcal E}(y)^2
&\le
\left(
\sum_{\boldsymbol k\in\mathscr K_{J,\mathcal E,k}}
\omega(\boldsymbol k)
\right)
\\
&\quad{}\times
\sum_{\boldsymbol k\in\mathscr K_{J,\mathcal E,k}}
\omega(\boldsymbol k)
\prod_{\ell=1}^B
a_{r_{J,\ell}+k_\ell}(y_\ell)^2.
\end{aligned}
\]
Extend the energy notation to degree zero by
\[
\gamma_0:=
\int_{\mathbb R}a_0(w)^2g_d(w)\,dw=1.
\]
This equality follows from \(a_0=1\) on positive reference support and \(\int g_d=1\).
Independence of the reference responses yields
\[
\int
\prod_{\ell=1}^B
a_{r_{J,\ell}+k_\ell}(y_\ell)^2\,
\nu_{\mathrm{resp}}(dy)
=
\prod_{\ell=1}^B\gamma_{r_{J,\ell}+k_\ell}.
\]
Thus
\[
\frac{\|\mathcal C_{J,k,\mathcal E}\|_{\mathrm{resp}}^2}{\binom Mk}
\le
\rho_{\nu_{\mathrm{rev}},\nu_{\mathrm{hid}}}^{\nu_{\mathrm{hid}}}
\sum_{\boldsymbol k\in\mathscr K_{J,\mathcal E,k}}
\omega(\boldsymbol k)
\prod_{\ell=1}^B\gamma_{r_{J,\ell}+k_\ell}.
\]
For \(k<\nu_{\mathrm{hid}}\), the restricted family is empty and both sides vanish.

Summing over \(k=0,\ldots,M\) covers each admissible degree vector exactly once. The remaining restrictions are rowwise, so the sum factors:
\[
\begin{aligned}
&\sum_{k=0}^M
\frac{\|\mathcal C_{J,k,\mathcal E}\|_{\mathrm{resp}}^2}{\binom Mk}
\\
&\quad\le
\rho_{\nu_{\mathrm{rev}},\nu_{\mathrm{hid}}}^{\nu_{\mathrm{hid}}}
\prod_{\ell\in\mathcal B_J}
\left(
\sum_{k_\ell=0}^{u_\ell}
\binom{u_\ell}{k_\ell}\gamma_{r_{J,\ell}+k_\ell}
\right)
\prod_{\ell\in\mathcal E}
\left(
\sum_{k_\ell=1}^{u_\ell}
\binom{u_\ell}{k_\ell}\gamma_{k_\ell}
\right).
\end{aligned}
\]
Define the active-support envelope by
\[
\begin{aligned}
\mathcal A_H
:=
\sum_{J\subseteq\bigcup_\ell R_\ell}
\sum_{\mathcal E\subseteq\{1,\ldots,B\}\setminus\mathcal B_J}
&\rho_{|\mathcal B_J|,|\mathcal E|}^{|\mathcal E|}
\prod_{\ell\in\mathcal B_J}
\left(
\sum_{k_\ell=0}^{u_\ell}
\binom{u_\ell}{k_\ell}\gamma_{r_{J,\ell}+k_\ell}
\right)
\\[-2pt]
&{}\times
\prod_{\ell\in\mathcal E}
\left(
\sum_{k_\ell=1}^{u_\ell}
\binom{u_\ell}{k_\ell}\gamma_{k_\ell}
\right).
\end{aligned}
\]
Restoring the coefficient \((J,k)=(\varnothing,0)\) adds exactly one to the labeled energy. The support decomposition and the last bound therefore give
\[
\chi_H^2+1\le\mathcal A_H,
\qquad
\chi_H^2\le\mathcal A_H-1
\quad\text{when }M\ge Bd/4.
\]
If \(M<Bd/4\), multiplying either quantity by the remaining-capacity indicator gives zero. This includes \(M=0\).
% lean: CausalSmith.Experimentation.ThinnedgraphAdditiveRiskFrontier.actualLabelEnergy_le_activeEnvelope

\item Reindex the envelope by revealed-active and hidden-active rows. For integers \(0\le r\le d\), define
\[
E_{\mathrm{rev}}(r)
:=
\sum_{s_{\mathrm{rev}}=1}^{r}
\binom r{s_{\mathrm{rev}}}
\sum_{k_{\mathrm{row}}=0}^{d-r}
\binom{d-r}{k_{\mathrm{row}}}
\gamma_{s_{\mathrm{rev}}+k_{\mathrm{row}}},
\]
and
\[
E_{\mathrm{hid}}(r)
:=
\sum_{k_{\mathrm{row}}=1}^{d-r}
\binom{d-r}{k_{\mathrm{row}}}\gamma_{k_{\mathrm{row}}}.
\]
The sums are empty when their upper limits are below their lower limits.

Since the \(R_\ell\) are disjoint, each global revealed subset \(J\) corresponds bijectively to its rowwise subsets \(J\cap R_\ell\).
For a row in \(\mathcal B_J\), summing over its nonempty revealed subsets gives
\[
\sum_{\substack{J_\ell\subseteq R_\ell\\J_\ell\ne\varnothing}}
\sum_{k_\ell=0}^{u_\ell}
\binom{u_\ell}{k_\ell}\gamma_{|J_\ell|+k_\ell}
=
E_{\mathrm{rev}}(r_\ell).
\]
For a hidden-active row, its factor is \(E_{\mathrm{hid}}(r_\ell)\).
It follows that
\[
\mathcal A_H
=
\sum_{\substack{
\mathcal B_{\mathrm{rev}},\mathcal B_{\mathrm{hid}}
\subseteq\{1,\ldots,B\}\\
\mathcal B_{\mathrm{rev}}\cap\mathcal B_{\mathrm{hid}}=\varnothing}}
\rho_{|\mathcal B_{\mathrm{rev}}|,|\mathcal B_{\mathrm{hid}}|}^{|\mathcal B_{\mathrm{hid}}|}
\prod_{\ell\in\mathcal B_{\mathrm{rev}}}E_{\mathrm{rev}}(r_\ell)
\prod_{\ell\in\mathcal B_{\mathrm{hid}}}E_{\mathrm{hid}}(r_\ell).
\]
The pair of empty row sets contributes exactly one. Removing it gives
\[
\begin{aligned}
\mathcal A_H-1
=
\sum_{\substack{
\mathcal B_{\mathrm{rev}}\cap\mathcal B_{\mathrm{hid}}=\varnothing\\
|\mathcal B_{\mathrm{rev}}|+|\mathcal B_{\mathrm{hid}}|>0}}
&\rho_{|\mathcal B_{\mathrm{rev}}|,|\mathcal B_{\mathrm{hid}}|}^{|\mathcal B_{\mathrm{hid}}|}
\\
&{}\times
\prod_{\ell\in\mathcal B_{\mathrm{rev}}}E_{\mathrm{rev}}(r_\ell)
\prod_{\ell\in\mathcal B_{\mathrm{hid}}}E_{\mathrm{hid}}(r_\ell).
\end{aligned}
\]
The row sets in this display are subsets of \(\{1,\ldots,B\}\).
% lean: CausalSmith.Experimentation.ThinnedgraphAdditiveRiskFrontier.actualLabelActiveEnvelope_sub_one_eq_row_energy

\item We bound the average of this nonconstant envelope under the detailed graph law. Counting a nonempty subset of the \(d\) row slots according to whether it contains a revealed slot gives
\[
\begin{aligned}
E_{\mathrm{rev}}(r)+E_{\mathrm{hid}}(r)
&=
\sum_{s=1}^d
\left(
\sum_{s_{\mathrm{rev}}=0}^s
\binom r{s_{\mathrm{rev}}}
\binom{d-r}{s-s_{\mathrm{rev}}}
\right)\gamma_s
\\
&=
\sum_{s=1}^d\binom ds\gamma_s=\eta.
\end{aligned}
\]
The second equality is the binomial convolution identity, with out-of-range binomial coefficients zero. All energies are nonnegative, so
\[
E_{\mathrm{rev}}(r)\ge0,
\qquad
0\le E_{\mathrm{hid}}(r)\le\eta.
\]

Fix a source partition \(S\). A source has \(d\) possible outgoing arrows, each audited independently. Its association-revelation indicator therefore has probability \(p=1-(1-q)^d\). Different sources use disjoint audit coordinates, so these indicators are independent conditional on \(S\), including across rows. For any specified \(s\)-subset of a row, the probability of at least one revealed member is \(1-(1-p)^s\). Thus
\[
\mathbb E\!\left[E_{\mathrm{rev}}(r_\ell)\mid S\right]
=
\sum_{s=1}^d
\binom ds\bigl(1-(1-p)^s\bigr)\gamma_s.
\]
The Bernoulli inequality gives
\[
(1-p)^s\ge1-sp,
\qquad
1-(1-p)^s\le sp
\quad(s\in\mathbb N_0).
\]
For completeness, the first inequality starts with equality at \(s=0\); multiplying it by \(1-p\ge0\) gives
\[
(1-p)^{s+1}
\ge(1-sp)(1-p)
=1-(s+1)p+sp^2
\ge1-(s+1)p.
\]
Consequently,
\[
\mathbb E\!\left[E_{\mathrm{rev}}(r_\ell)\mid S\right]
\le p\eta_1,
\qquad
\mathbb E\!\left[E_{\mathrm{hid}}(r_\ell)\mid S\right]
\le\eta.
\]

Every term in \(\mathcal A_H-1\) is nonnegative. On the remaining-capacity event its indicator leaves the term unchanged; outside that event the indicated term is zero. Hence
\[
0\le
\mathbf1\{M\ge Bd/4\}(\mathcal A_H-1)
\le\mathcal A_H-1.
\]
For disjoint row sets, conditional independence gives
\[
\begin{aligned}
&\mathbb E\!\left[
\prod_{\ell\in\mathcal B_{\mathrm{rev}}}E_{\mathrm{rev}}(r_\ell)
\prod_{\ell\in\mathcal B_{\mathrm{hid}}}E_{\mathrm{hid}}(r_\ell)
\;\middle|\;S
\right]
\\
&\quad=
\prod_{\ell\in\mathcal B_{\mathrm{rev}}}
\mathbb E[E_{\mathrm{rev}}(r_\ell)\mid S]
\prod_{\ell\in\mathcal B_{\mathrm{hid}}}
\mathbb E[E_{\mathrm{hid}}(r_\ell)\mid S]
\\
&\quad\le
(p\eta_1)^{|\mathcal B_{\mathrm{rev}}|}
\eta^{|\mathcal B_{\mathrm{hid}}|}.
\end{aligned}
\]
Empty row sets give empty products equal to one.

There are
\(\binom B{\nu_{\mathrm{rev}}}\) choices of
\(\mathcal B_{\mathrm{rev}}\) of size \(\nu_{\mathrm{rev}}\), followed by
\(\binom{B-\nu_{\mathrm{rev}}}{\nu_{\mathrm{hid}}}\) choices of a disjoint
\(\mathcal B_{\mathrm{hid}}\) of size \(\nu_{\mathrm{hid}}\).
When their sizes exceed \(B\) in total, the latter coefficient is zero.
Summing the preceding product bounds therefore gives, for every fixed partition,
\[
\mathbb E\!\left[
\mathbf1\{M\ge Bd/4\}(\mathcal A_H-1)
\mid S
\right]
\le\mathfrak S_{\mathrm{nc}}.
\]

The detailed retained graph generated by a partition and audit array is
\[
H(S,W)
:=
\{(j,i):
\exists\ell,\ j\in S_\ell,\ i\in C_\ell,\ W_{ji}=1\}.
\]
Averaging the fixed-partition inequality over the uniform partition law thus averages over the actual marginal of this full graph:
\[
\begin{aligned}
\mathbb E_H\!\left[
\mathbf1\{M\ge Bd/4\}(\mathcal A_H-1)
\right]
&=
\mathbb E_S\!\left[
\mathbb E\!\left[
\mathbf1\{M\ge Bd/4\}(\mathcal A_{H(S,W)}-1)
\mid S
\right]\right]
\\
&\le\mathfrak S_{\mathrm{nc}}.
\end{aligned}
\]
This average includes every source label and every retained-edge subset.
% lean: CausalSmith.Experimentation.ThinnedgraphAdditiveRiskFrontier.retained_row_energy_average_le

\item Combining the exact conditional density and labeled-energy identities with the good-event envelope bound yields
\[
\begin{aligned}
\mathbb E_H\!\left[
\mathbf1\{M\ge m/4\}\chi_H^2
\right]
&\le
\mathbb E_H\!\left[
\mathbf1\{M\ge Bd/4\}(\mathcal A_H-1)
\right]
\\
&\le\mathfrak S_{\mathrm{nc}}
\\
&\le
\frac{\exp(eBp\eta_1)}{1-4e\eta}-1,
\end{aligned}
\]
where \(m=Bd\). This is the asserted bound under the actual complete retained-graph marginal.
% lean: CausalSmith.Experimentation.ThinnedgraphAdditiveRiskFrontier.hidden_allocation_contraction
\end{enumerate}
\end{proof}

The block-size conditions in \cref{obj:lem:hidden-allocation-contraction} place the source and recipient groups within the finite population. Independent auditing and the assignment condition preserve the experimental design used by the conditional response law. The channel-smallness condition keeps the denominator positive and specifies the range in which the contraction bound applies.

The bound combines two aspects of the experiment. The exponential term depends on the revelation probability and the order-weighted response energy, while the denominator depends on the total nonconstant energy. The indicator restricts the averaged discrepancy to graph realizations leaving at least one quarter of the original source capacity hidden. All graph labels and detailed retained subsets remain part of the averaging law, so the comparison applies to the information represented by the complete retained graph.

Together, \cref{obj:lem:baseline-translation-affinity,obj:lem:walsh-channel-energy,obj:lem:hidden-allocation-contraction} provide the distance comparisons, response-energy control, and allocation bound used in \cref{obj:thm:block-testing-scale,obj:thm:nonlocal-testing-answer}. Their roles are complementary: translations control changes in block responses, Walsh energies quantify their dependence on source signs, and the conditional-reference comparison accounts for the hidden allocation under the actual observation design.


\section{Testing bounds and minimax proofs}

The distance comparisons in \cref{obj:lem:baseline-translation-affinity,obj:lem:walsh-channel-energy,obj:lem:hidden-allocation-contraction} support a uniform testing comparison for the complete-record block experiment. The constant-target property of \cref{obj:lem:block-law} then connects testing difficulty to squared estimation risk. Together with the supplied-graph comparison and the observable upper guarantee, these results characterize attainable precision throughout the admissible degree and retention ranges.

\subsection{Uniform block testing and its radius}

For probability laws \(P\) and \(Q\) on a common measurable space, total-variation distance is
\[
\operatorname{TV}(P,Q)
=\sup_{A\text{ measurable}}|P(A)-Q(A)|.
\]
This convention agrees with \cref{obj:lem:baseline-translation-affinity}. The comparison concerns the opposite-sign complete-record mixture laws \(P_{+1,h}\) and \(P_{-1,h}\) of \cref{obj:def:block-prior}, where \(h\) is the block amplitude.

With \(m=Bd\), the source count, and \(p=1-(1-q)^d\), the association-revelation probability specified in \cref{obj:def:testing-handle}, define the testing scale by
\[
s(B,d,q)=
\begin{cases}
\min\{1,d/\sqrt{mp}\},&p>0,\\
1,&p=0.
\end{cases}
\]
The testing radius is the supremum of amplitudes in \([0,1/4]\) for which the two mixture laws have total variation at most one half:
\[
\bar h(B,d,q)
=\sup\left\{h\in[0,1/4]:
\operatorname{TV}(P_{+1,h},P_{-1,h})\le\tfrac12\right\}.
\]
The universal amplitude constant is \(\kappa=(100\pi)^{-1}\).

The preceding distance bounds have complementary roles in this comparison. Baseline translations control changes in response locations, while the Walsh energy bound measures response dependence on source signs. The hidden-allocation comparison accounts for the source associations remaining uncertain after the audit. Through the reconstruction property of \cref{obj:lem:block-law}, these comparisons concern the complete graph, assignment, and outcome record, including the observed global hidden treated-source count. The resulting statement provides both a uniformly small-distance range and a complementary lower bound on distinguishability.

\begin{theoremv}{thm:block-testing-scale}[Block testing scale bounds]
Suppose the block experiment satisfies the following conditions:
\begin{itemize}
\item \textbf{(Population and blocks.)} The population size \(n\), block count \(B\), and degree \(d\) are integers satisfying
\[
n\ge4,\qquad B\ge1,\qquad d\ge1,\qquad 2Bd\le n.
\]
\item \textbf{(Retention.)} The edge-retention probability satisfies \(q\in[0,1]\).
\item \textbf{(Assignment.)} The experimental design satisfies \cref{obj:ass:assignment}.
\item \textbf{(Audit.)} The experimental design satisfies \cref{obj:ass:audit} at retention probability \(q\).
\item \textbf{(Independence.)} The experimental design satisfies \cref{obj:ass:design-independence}.
\end{itemize}
Put \(m=Bd\), the number of sources, and \(p=1-(1-q)^d\), the probability that a source association is revealed. Define the testing scale and universal constant by
\[
s(B,d,q)=
\begin{cases}
\min\{1,d/\sqrt{mp}\},&p>0,\\
1,&p=0,
\end{cases}
\qquad
\kappa=\frac{1}{100\pi}.
\]
Let \(P_{+1,h}\) and \(P_{-1,h}\) be the complete original-record mixture laws under the opposite-sign block priors, using this experimental design, and let \(\operatorname{TV}\) denote total-variation distance. For every real amplitude \(h\),
\[
0\le h\le\kappa s(B,d,q)
\quad\Longrightarrow\quad
\operatorname{TV}(P_{+1,h},P_{-1,h})\le\frac12.
\]
For every real amplitude \(h\) satisfying \(0<h\le1/4\) and \(p>0\),
\[
\operatorname{TV}(P_{+1,h},P_{-1,h})
\ge 1-\frac{7d^2}{8h^2mp}.
\]
Finally, define the testing radius by
\[
\bar h(B,d,q)
=
\sup\left\{
h\in[0,1/4]:
\operatorname{TV}(P_{+1,h},P_{-1,h})\le\frac12
\right\}.
\]
Then
\[
\kappa s(B,d,q)\le\bar h(B,d,q)\le2s(B,d,q).
\]
\end{theoremv}

The scale in \cref{obj:thm:block-testing-scale} balances the amplitude against the number of revealed source associations. At positive revelation probability, the distance lower bound depends on the squared amplitude through \(h^2mp/d^2\), while the uniform upper bound supplies a range of amplitudes at which the signs remain difficult to distinguish. The truncation at one keeps the scale bounded when revelation is sparse. At zero revelation, the constant-amplitude comparison continues to apply. The radius bounds characterize the supremum directly, using comparisons that hold throughout their stated amplitude ranges.

Under the canonical assignment--audit product design, the same comparison can be recorded entirely in terms of the block parameters. The following statement consolidates the testing radius and the uniform amplitude guarantee, and gives the explicit degree-two expression.

\begin{theoremv}{thm:nonlocal-testing-answer}[Nonlocal block testing bounds]
Let the population size \(n\), block count \(B\), degree \(d\), and retention probability \(q\) satisfy:
\begin{itemize}
\item \textbf{(Population.)} \(n\) is an integer with \(n\ge4\).
\item \textbf{(Blocks and degree.)} \(B,d\) are integers with \(B\ge1\) and \(d\ge1\).
\item \textbf{(Capacity.)} \(2Bd\le n\).
\item \textbf{(Retention.)} \(q\in[0,1]\).
\end{itemize}
Define the source count \(m\), association-revelation probability \(p\), and testing scale \(s\) by
\[
m=Bd,\qquad p=1-(1-q)^d,\qquad
s=
\begin{cases}
\min\{1,d/\sqrt{mp}\},&p>0,\\
1,&p=0.
\end{cases}
\]
Let \(P_{+1,h}\) and \(P_{-1,h}\) be the complete original-record block mixture laws under the canonical assignment-audit product design with retention probability \(q\), at amplitude \(h\) and signs \(+1\) and \(-1\), respectively. With \(\operatorname{TV}\) denoting total-variation distance, define the testing radius
\[
\bar h
=\sup\left\{h\in[0,1/4]:
\operatorname{TV}(P_{+1,h},P_{-1,h})\le\tfrac12\right\}.
\]
Then
\[
(100\pi)^{-1}s\le\bar h\le2s.
\]
Moreover, for every real \(h\) satisfying \(0\le h\le(100\pi)^{-1}s\),
\[
\operatorname{TV}(P_{+1,h},P_{-1,h})\le\tfrac12.
\]
If \(d=2\), the scale satisfies
\[
s=
\begin{cases}
\min\left\{1,\sqrt{\dfrac{2}{B(2q-q^2)}}\right\},&q>0,\\
1,&q=0.
\end{cases}
\]
\end{theoremv}

For the canonical design, \cref{obj:thm:nonlocal-testing-answer} expresses testing difficulty through block count, degree, and retention. In the degree-two case, \(2q-q^2\) is the probability that at least one of a source's two outgoing arrows is retained. Its appearance makes explicit how collection changes the testing scale. The uniform amplitude clause supplies the distance guarantee at the amplitude named in \cref{obj:def:testing-handle}.

\subsection{From testing to squared estimation risk}

The estimation reduction uses generic priors over fixed schedules. Write \(\mathsf Q_+\) and \(\mathsf Q_-\), the complete-record mixture laws of \cref{obj:lem:full-record-two-prior-risk}, for the laws obtained by drawing a schedule from the respective prior and then independently running the canonical assignment--audit experiment. Thus, for a measurable record event, its mixture probability is the prior average of its conditional probability given the fixed schedule.

Suppose the positive and negative priors have constant causal targets \(+\Delta\) and \(-\Delta\), respectively, where \(\Delta>0\) is their common target magnitude. The reduction compares the overlap of the two record laws with the squared error required to estimate these separated targets. Its support condition connects prior-average loss to the minimax schedule class, and its measurability condition makes the induced record laws well defined. The following statement includes arbitrary measurable randomized estimators of the complete record.

\begin{lemmav}{lem:full-record-two-prior-risk}[Full-record two-prior risk bound]
Let \(n\) be the population size, \(d\) the degree bound, and \(q\) the audit-retention probability. Suppose:
\begin{itemize}
\item \textbf{(Parameters.)} \(n,d\) are integers satisfying \(n\ge4\) and \(1\le d\le n-1\), and \(q\in[0,1]\).
\item \textbf{(Priors.)} Two probability priors, each represented by a map from a measurable parameter space to fixed schedules, assign probability one to \(\mathcal M_{n,d}\) as defined in \cref{obj:def:model}.
\item \textbf{(Targets.)} For some \(\Delta>0\), the positive and negative priors respectively satisfy \(\tau(\theta)=\Delta\) and \(\tau(\theta)=-\Delta\) almost surely, where \(\tau(\theta)\), the population all-treated-versus-all-control contrast, is
\[
\tau(\theta)=\frac1n\sum_{i=1}^{n}
\bigl(Y_i(1,\ldots,1)-Y_i(0,\ldots,0)\bigr),
\]
and \(Y_i(z)\) denotes unit \(i\)'s additive potential outcome under schedule \(\theta\) at binary assignment \(z\).
\item \textbf{(Measurability.)} For each prior representation, the map from the parameter, assignment, and audit marks to the complete observed record is jointly measurable.
\end{itemize}
Under each prior, draw the schedule first, then independently draw \(Z\), the treatment-assignment vector, with independent Bernoulli-\(1/2\) coordinates, and \(W\), the array of audit marks on all ordered off-diagonal pairs, with independent Bernoulli-\(q\) coordinates, independently of \(Z\). The retained graph \(H\) and complete record \(O\) are
\[
H_{ij}=\mathbf1\{j\to i\text{ in }G\}\,W_{ij}\quad(i\ne j),
\qquad
O=(H,Z,Y(Z)),
\]
where \(G\) is the schedule's graph and \(Y(Z)\) is the vector of realized outcomes. Let \(\mathsf Q_+\) and \(\mathsf Q_-\) be the resulting complete-record mixture laws under the positive and negative priors. Then the minimax squared risk over measurable randomized estimators defined in \cref{obj:def:risk} satisfies
\[
R_{n,d}(q)\ge
\frac{\Delta^2}{2}
\left(1-\operatorname{TV}(\mathsf Q_+,\mathsf Q_-)\right),
\]
where \(\operatorname{TV}\) is the total-variation distance, defined as the supremum over measurable events of the absolute difference between their probabilities.
\end{lemmav}

\begin{proof}[Proof of \cref{obj:lem:full-record-two-prior-risk}]
\begin{enumerate}
\item Define the complete-record space and the uniform seed law by
\[
\mathcal O
=
\{0,1\}^{\{(i,j)\in V^2:i\ne j\}}
\times\{0,1\}^{V}\times\mathbb R^{V},
\qquad
\nu_{\mathrm{seed}}(d\xi)
=
\mathbf1_{[0,1]}(\xi)\,d\xi.
\]
Since \(q\in[0,1]\), the assignment and audit laws are probability laws. Joint measurability of the record maps ensures that their pushforwards are probability laws, and mixing them under probability priors preserves unit mass. Thus
\[
\mathsf Q_+(\mathcal O)=\mathsf Q_-(\mathcal O)=1,
\qquad
\nu_{\mathrm{seed}}(\mathbb R)=1.
\]
The laws augmented by the independent estimator seed are
\[
\mathsf Q_\pm^{\mathrm{aug}}
:=
\mathsf Q_\pm\otimes\nu_{\mathrm{seed}}.
\]
% lean: CausalSmith.Experimentation.ThinnedgraphAdditiveRiskFrontier.mixtureLaw_probability

\item We first establish
\begin{equation}\label{eq:full-record-two-prior-tv-seed}
\operatorname{TV}
\bigl(\mathsf Q_+^{\mathrm{aug}},\mathsf Q_-^{\mathrm{aug}}\bigr)
=
\operatorname{TV}(\mathsf Q_+,\mathsf Q_-).
\end{equation}
For any measurable event \(\mathcal A\subseteq\mathcal O\times\mathbb R\), define
\[
\psi_{\mathcal A}(o)
:=
\int_{\mathbb R}\mathbf1_{\mathcal A}(o,\xi)\,
\nu_{\mathrm{seed}}(d\xi),
\qquad o\in\mathcal O.
\]
This is measurable and takes values in \([0,1]\), and product integration gives
\[
\mathsf Q_\pm^{\mathrm{aug}}(\mathcal A)
=
\int_{\mathcal O}\psi_{\mathcal A}(o)\,\mathsf Q_\pm(do).
\]
The layer-cake identity gives, for each sign,
\[
\int_{\mathcal O}\psi_{\mathcal A}(o)\,\mathsf Q_\pm(do)
=
\int_{(0,1]}
\mathsf Q_\pm\bigl(\{o\in\mathcal O:v\le\psi_{\mathcal A}(o)\}\bigr)\,dv.
\]
Consequently, the defining event bound for total variation yields
\[
\begin{aligned}
\left|
\mathsf Q_+^{\mathrm{aug}}(\mathcal A)
-\mathsf Q_-^{\mathrm{aug}}(\mathcal A)
\right|
&\le
\int_{(0,1]}
\left|
\mathsf Q_+\bigl(\{o:v\le\psi_{\mathcal A}(o)\}\bigr)
-\mathsf Q_-\bigl(\{o:v\le\psi_{\mathcal A}(o)\}\bigr)
\right|\,dv\\
&\le \operatorname{TV}(\mathsf Q_+,\mathsf Q_-).
\end{aligned}
\]
Taking the supremum over \(\mathcal A\) proves one inequality. For the reverse inequality, every measurable \(\mathcal C\subseteq\mathcal O\) satisfies
\[
\mathsf Q_\pm^{\mathrm{aug}}(\mathcal C\times\mathbb R)
=
\mathsf Q_\pm(\mathcal C).
\]
Taking the supremum over these cylinder events proves the asserted equality.
% lean: Causalean.Stat.tvDist_compProd_eq

\item Fix any measurable randomized estimator \(T\) allowed by
\cref{obj:def:risk}, and define its worst-case risk and its two mixture risks by
\[
\mathcal R_T
:=
\sup_{\theta\in\mathcal M_{n,d}}\mathcal L(T;\theta,q),
\qquad
\mathcal L_{\mathrm{mix},\pm}(T)
:=
\int_{\mathcal O\times\mathbb R}
\bigl(T(o,\xi)\mp\Delta\bigr)^2\,
\mathsf Q_\pm^{\mathrm{aug}}(d(o,\xi)).
\]
These quantities take values in \([0,\infty]\).

Write the two prior representations as
\[
\vartheta_\pm:\Xi_\pm\longrightarrow
\{\text{fixed schedules on }V\},
\qquad
\pi_\pm(\Xi_\pm)=1,
\]
where \(\Xi_\pm\) are the measurable parameter spaces and \(\pi_\pm\) their probability measures. Their almost-sure support and target properties are
\[
\vartheta_\pm(\zeta)\in\mathcal M_{n,d},
\qquad
\tau(\vartheta_\pm(\zeta))=\pm\Delta
\quad\text{for }\pi_\pm\text{-almost every }\zeta\in\Xi_\pm.
\]
Nonnegative product integration, followed by integration over the prior, therefore gives
\[
\begin{aligned}
\mathcal L_{\mathrm{mix},\pm}(T)
&=
\int_{\Xi_\pm}
\mathcal L(T;\vartheta_\pm(\zeta),q)\,\pi_\pm(d\zeta)\\
&\le
\int_{\Xi_\pm}\mathcal R_T\,\pi_\pm(d\zeta)
=
\mathcal R_T.
\end{aligned}
\]
Indeed, conditional on the prior parameter, the record has the fixed-schedule experiment law and the seed remains independent and uniform. The constant-target property identifies the conditional squared loss, and membership in the class of \cref{obj:def:model} bounds that loss by its supremum. These identities and inequalities remain valid for infinite risks.
% lean: CausalSmith.Experimentation.ThinnedgraphAdditiveRiskFrontier.mixture_seeded_risk_le_worst

\item Define the measurable negative-decision event
\[
\mathcal E_T
:=
\{(o,\xi)\in\mathcal O\times\mathbb R:T(o,\xi)<0\}.
\]
By the definition of total variation,
\[
\mathsf Q_-^{\mathrm{aug}}(\mathcal E_T)
-\mathsf Q_+^{\mathrm{aug}}(\mathcal E_T)
\le
\operatorname{TV}
\bigl(\mathsf Q_+^{\mathrm{aug}},\mathsf Q_-^{\mathrm{aug}}\bigr).
\]
Since the augmented laws are probability laws, it follows that
\[
\begin{aligned}
\mathsf Q_+^{\mathrm{aug}}(\mathcal E_T)
+\mathsf Q_-^{\mathrm{aug}}(\mathcal E_T^c)
&=
1+\mathsf Q_+^{\mathrm{aug}}(\mathcal E_T)
-\mathsf Q_-^{\mathrm{aug}}(\mathcal E_T)\\
&\ge
1-\operatorname{TV}
\bigl(\mathsf Q_+^{\mathrm{aug}},\mathsf Q_-^{\mathrm{aug}}\bigr).
\end{aligned}
\]
% lean: Causalean.Stat.one_sub_tvDist_le_test

On \(\mathcal E_T\), \(T(o,\xi)<0\) and \(\Delta>0\), so
\((T(o,\xi)-\Delta)^2\ge\Delta^2\). On its complement,
\(T(o,\xi)\ge0\), so \((T(o,\xi)+\Delta)^2\ge\Delta^2\).
Together with nonnegativity off the respective events, these give the pointwise bounds
\[
\Delta^2\mathbf1_{\mathcal E_T}(o,\xi)
\le (T(o,\xi)-\Delta)^2,
\qquad
\Delta^2\mathbf1_{\mathcal E_T^c}(o,\xi)
\le (T(o,\xi)+\Delta)^2.
\]
Integrating under the corresponding augmented laws yields
\[
\Delta^2\mathsf Q_+^{\mathrm{aug}}(\mathcal E_T)
\le\mathcal L_{\mathrm{mix},+}(T),
\qquad
\Delta^2\mathsf Q_-^{\mathrm{aug}}(\mathcal E_T^c)
\le\mathcal L_{\mathrm{mix},-}(T).
\]
% lean: CausalSmith.Experimentation.ThinnedgraphAdditiveRiskFrontier.threshold_squared_error_le_risk

Combining these bounds in order gives
\[
\begin{aligned}
\Delta^2\left(
1-\operatorname{TV}
\bigl(\mathsf Q_+^{\mathrm{aug}},\mathsf Q_-^{\mathrm{aug}}\bigr)
\right)
&\le
\Delta^2\left(
\mathsf Q_+^{\mathrm{aug}}(\mathcal E_T)
+\mathsf Q_-^{\mathrm{aug}}(\mathcal E_T^c)
\right)\\
&\le
\mathcal L_{\mathrm{mix},+}(T)
+\mathcal L_{\mathrm{mix},-}(T)\\
&\le 2\mathcal R_T.
\end{aligned}
\]
Cancelling the positive finite factor \(2\), which is valid also for extended nonnegative risks, proves
\[
\frac{\Delta^2}{2}\left(
1-\operatorname{TV}
\bigl(\mathsf Q_+^{\mathrm{aug}},\mathsf Q_-^{\mathrm{aug}}\bigr)
\right)
\le\mathcal R_T.
\]
% lean: CausalSmith.Experimentation.ThinnedgraphAdditiveRiskFrontier.two_prior_squared_testing

\item Substitute the total-variation identity in \cref{eq:full-record-two-prior-tv-seed}. The resulting bound holds for every admissible \(T\), so \cref{obj:def:risk} gives
\[
\frac{\Delta^2}{2}
\left(1-\operatorname{TV}(\mathsf Q_+,\mathsf Q_-)\right)
\le
\inf_T\mathcal R_T
=
R_{n,d}(q).
\]
% lean: CausalSmith.Experimentation.ThinnedgraphAdditiveRiskFrontier.full_record_two_prior_risk
\end{enumerate}
\end{proof}

The lower bound in \cref{obj:lem:full-record-two-prior-risk} combines target separation with record-law overlap. A larger constant target magnitude increases the squared-loss consequence of ambiguity, while a smaller total-variation distance preserves more overlap between the experiments. For the block priors, \cref{obj:lem:block-law} supplies support in the schedule class and the target magnitude \(mh/n\). The testing guarantees of \cref{obj:thm:block-testing-scale,obj:thm:nonlocal-testing-answer} therefore provide the record-law comparison needed by this reduction at the prescribed testing amplitude.

\subsection{Combining the comparisons across degrees and retention}

The block and supplied-graph comparisons address complementary parts of the admissible parameter range. The block experiment in \cref{obj:lem:block-law} requires room for equal numbers of sources and recipients. Its testing scale captures the dependence on association revelation, and its target magnitude records the fraction of the population occupied by recipients. The lower comparison in \cref{obj:lem:supplied-block-record-lower} applies throughout \(1\le d\le n-1\) and \(q\in[0,1]\), including degrees for which the population capacity restricts the block layout. It also identifies the degree-dependent estimation difficulty present when the true graph is supplied.

On the attainable side, \cref{obj:thm:observable-upper} controls the observable rule of \cref{obj:def:upper-rule}. The moment identity in \cref{obj:lem:score-audit} separates assignment variation from audit variation, and the upper envelope of \cref{obj:def:upper-envelope} governs the choice between the clipped score and zero. These ingredients supply the upper comparison for the named rule in \cref{obj:def:frontier-rule}; the block testing reduction and the supplied-graph lower comparison supply its minimax counterpart.

The risk scale \(F(n,d,q)\) of \cref{obj:def:frontier-scale} packages these comparisons through the source-association revelation probability. The positive-retention expression separates the degree and collection contributions and truncates their sum at one, representing saturation of the minimax risk order. For population sequences, these two contributions become the respective consistency criteria. The precision pair below records the already-defined scale and rule so that the explicit-constant statement can refer to them jointly.

\begin{definitionv}{def:frontier-handle}[Precision pair $(F,T^\star_{n,d,q})$]
The precision handle is the scale–rule pair \((F,T^\star_{n,d,q})\), defined by
\[
F(n,d,q)=
\begin{cases}
\min\left\{1,\dfrac{d^2}{n[1-(1-q)^d]}\right\},&q>0,\\
1,&q=0,
\end{cases}
\qquad
T^\star_{n,d,q}(O)=
\begin{cases}
\max\{-1,\min\{1,T_q(O)\}\},&q>0,\ u(n,d,q)<1,\\
0,&\text{otherwise}.
\end{cases}
\]
Here \(O\) is the complete observed record, \(T_q\) is the observable inverse-inclusion score, and \(u\) is the upper-risk envelope. The associated positive-retention elbow expression is
\[
\min\left\{1,\frac{d^2}{n}+\frac{d}{nq}\right\}.
\]
For admissible population sequences \(d_n,q_n\), the associated consistency criteria are
\[
\frac{d_n^2}{n}\longrightarrow0,
\qquad
\frac{d_n}{nq_n}\longrightarrow0,
\]
with the second ratio assigned \(+\infty\) when \(q_n=0\). The uniform risk comparison, equivalence with the elbow expression, and necessity and sufficiency of these criteria are proved properties of this pair. The original-record converse uses the complete conditional likelihood and retains the retained graph \(H\), assignment vector \(Z\), realized outcomes \(Y(Z)\), and observed global hidden treated-source count \(K\).
\end{definitionv}

The following statement supplies explicit universal constants for the primary frontier in \cref{obj:thm:observed-record-precision-frontier} and consolidates its endpoint and sequence deductions. The scale, rule, and estimator class remain those of \cref{obj:def:frontier-scale,obj:def:frontier-rule,obj:def:risk}.

\begin{theoremv}{thm:frontier-answer}[Explicit risk and consistency frontier]
Consider the canonical independent assignment--audit product design, with known edge-retention probability \(q\). Let \(\mathcal M_{n,d}\) be the fixed graph-and-coefficient schedules satisfying the common off-diagonal in-degree and out-degree bound \(d\) and the unit coefficient-mass bound. Write \(\mathcal L(T;\theta,q)\) for the expected squared error of \(T(O,\xi)\) in estimating \(\tau(\theta)\), the population all-treated-versus-all-control contrast, where \(O\) is the complete observed record and \(\xi\) is the independent uniform randomization seed. The minimax risk is
\[
R_{n,d}(q)
=
\inf_T\sup_{\theta\in\mathcal M_{n,d}}\mathcal L(T;\theta,q),
\]
where the infimum ranges over all measurable real-valued randomized estimators of the complete record.

Suppose that:
\begin{itemize}
\item \textbf{(Population size.)} \(n\) is a natural number with \(n\ge4\).
\item \textbf{(Degree bound.)} \(d\) is a natural number with \(1\le d\le n-1\).
\item \textbf{(Retention.)} \(q\in[0,1]\).
\end{itemize}
Then the risk scale \(F(n,d,q)\) and attaining rule \(T^\star_{n,d,q}\) of \cref{obj:def:frontier-scale,obj:def:frontier-rule} satisfy
\[
\frac{F(n,d,q)}{2560000\pi^2}
\le R_{n,d}(q)
\le
\sup_{\theta\in\mathcal M_{n,d}}
\mathcal L(T^\star_{n,d,q};\theta,q)
\le 24F(n,d,q).
\]

For every pair of sequences \(d_n\in\mathbb N\) and \(q_n\in\mathbb R\) satisfying
\[
1\le d_n\le n-1,\qquad q_n\in[0,1]
\qquad\text{for every }n\ge4,
\]
there exists a sequence of measurable randomized estimators \(T_n\) of the complete record such that
\[
\sup_{\theta\in\mathcal M_{n,d_n}}
\mathcal L(T_n;\theta,q_n)\longrightarrow0
\]
if and only if
\[
\frac{d_n^2}{n}\longrightarrow0
\qquad\text{and}\qquad
\begin{cases}
+\infty,&q_n=0,\\[2pt]
\dfrac{d_n}{nq_n},&q_n\ne0
\end{cases}
\longrightarrow0.
\]
All limits are as \(n\to\infty\).

For every finite population \(V\), every \(n,d\in\mathbb N\), every \(q\in\mathbb R\), and every record \(O\) on \(V\), the total observable rules of \cref{obj:def:frontier-rule,obj:def:upper-rule} agree:
\[
T^\star_{n,d,q}(O)=T^{\mathrm{up}}(O).
\]
The estimator attaining the displayed risk bound is the seeded measurable representative of this upper rule.

Finally, for every \(n\ge4\) and \(1\le d\le n-1\),
\[
F(n,d,0)=1,
\qquad
F(n,d,1)=\min\left\{1,\frac{d^2}{n}\right\}.
\]
For every \(q\in(0,1]\),
\[
\frac12\min\left\{1,\frac{d^2}{n}+\frac{d}{nq}\right\}
\le F(n,d,q)
\le
(1-e^{-1})^{-1}
\min\left\{1,\frac{d^2}{n}+\frac{d}{nq}\right\}.
\]
\end{theoremv}

The finite-population comparison in \cref{obj:thm:frontier-answer} gives matching risk scales with universal constants. The positive-retention expression makes the two contributions explicit: degree growth affects precision even with a supplied graph, and partial collection adds a retention-dependent contribution. At full retention the scale is \(\min\{1,d^2/n\}\); at zero retention it equals one. The attaining rule translates the upper comparison into an observable procedure using the known design parameters.

For admissible population sequences, sufficiency is constructive: the seeded measurable representatives of \(T^\star_{n,d_n,q_n}\) in \cref{obj:def:frontier-rule} attain uniformly vanishing squared risk whenever both displayed ratios tend to zero. Necessity applies to every sequence of measurable randomized estimators of the complete record: uniformly vanishing risk requires both the degree contribution and the collection contribution to vanish. Assigning the second ratio \(+\infty\) at zero retention incorporates that boundary into the same characterization.


\section{Verification note}

Proofs of the results appear in the appendices.

The population restrictions in \cref{obj:ass:in-degree,obj:ass:out-degree,obj:ass:coefficient-mass} define the schedule class of \cref{obj:def:model}. The assignment, audit, and independence conditions in \cref{obj:ass:assignment,obj:ass:audit,obj:ass:design-independence} specify the experimental laws. These restrictions and laws enter as hypotheses; the checked conclusions are conditional on them and on each result's stated parameter ranges.

The published-model correspondence in \cref{obj:lem:published-model} carries a generated theorem-local disclosure identifying its cited dependency and preserving the source locators. For that dependency, machine checking covers the cited conclusion's use and the remaining deductions, subject to the disclosed trust boundary. That local disclosure supplies the precise certification scope for the comparison.


\section{Proofs of the main results}\label{sec:deferred-proofs}

\begin{proof}[Proof of \cref{obj:thm:observable-upper}]
\begin{enumerate}
\item
The assignment--audit product design is a probability measure because its factors are products of Bernoulli probability laws with parameters \(1/2\) and \(q\in[0,1]\). For fixed \(n,d,q\), the branch condition in \cref{obj:def:upper-rule} is constant on the record space. On its positive-retention branch, the score is a finite sum of outcome coordinates multiplied by functions of the finite graph and assignment coordinates; clipping is continuous. The other branch is constant. Thus \(T^{\mathrm{up}}\) is real-valued and measurable on the complete record space.

Its deterministic extension to the randomized-estimator domain is
\[
(O,\xi)\longmapsto T^{\mathrm{up}}(O).
\]
This is an admissible estimator in \cref{obj:def:risk}, so
\[
R_{n,d}(q)
\le
\sup_{\theta\in\mathcal M_{n,d}}
\mathcal L(T^{\mathrm{up}};\theta,q).
\]
% lean: CausalSmith.Experimentation.ThinnedgraphAdditiveRiskFrontier.thinnedDesign_probabilityDesign
% lean: CausalSmith.Experimentation.ThinnedgraphAdditiveRiskFrontier.upperRule_measurable
% lean: Causalean.Stat.minimaxValueENNReal_le_worstCaseRisk

\item
Fix a schedule
\[
\theta=(G,a,t,b)\in\mathcal M_{n,d}.
\]
Because the estimator ignores the independent uniform seed, integration over that seed gives
\[
\begin{aligned}
\mathcal L(T^{\mathrm{up}};\theta,q)
&=
\mathbb E_{P_{\theta,q}}
\left[
\int_0^1
\bigl(T^{\mathrm{up}}(O)-\tau(\theta)\bigr)^2\,d\xi
\right]\\
&=
\mathbb E_{P_{\theta,q}}
\left[
\bigl(T^{\mathrm{up}}(O)-\tau(\theta)\bigr)^2
\right].
\end{aligned}
\]
We bound this loss separately on the two branches of \cref{obj:def:upper-rule}.
% lean: CausalSmith.Experimentation.ThinnedgraphAdditiveRiskFrontier.sqLoss_seedless

\item
First suppose
\[
q>0
\qquad\text{and}\qquad
u(n,d,q)<1.
\]
The product design satisfies \cref{obj:ass:assignment,obj:ass:audit,obj:ass:design-independence}. Together with the population, degree, and schedule hypotheses, these conditions permit application of \cref{obj:lem:score-audit}. It supplies
\[
\mathbb E_{P_{\theta,q}}[T_q]=\tau(\theta)
\]
and hence
\[
\begin{aligned}
\mathbb E_{P_{\theta,q}}
\left[(T_q-\tau(\theta))^2\right]
&=\operatorname{Var}_{P_{\theta,q}}(T_q)\\
&=\operatorname{Var}_{Z}(T_G)
+\frac{4(1-q)}{n^2q}
\sum_{i\in V}
\bigl|\{j:(j,i)\in G\}\bigr|
\,\mathbb E_Z[Y_i(Z)^2],
\end{aligned}
\]
as well as
\[
\operatorname{Var}_{Z}(T_G)\le\frac{5(d+1)^2}{n}.
\]

For each \(i\in V\) and binary assignment \(z\in\{0,1\}^V\), the additive outcome formula and coefficient-mass bound give
\[
\begin{aligned}
|Y_i(z)|
&=
\left|a_i+t_i z_i+\sum_{j:(j,i)\in G}b_{ij}z_j\right|\\
&\le
|a_i|+|t_i z_i|
+\sum_{j:(j,i)\in G}|b_{ij}z_j|\\
&\le
|a_i|+|t_i|
+\sum_{j:(j,i)\in G}|b_{ij}|\\
&\le1,
\end{aligned}
\]
where the penultimate inequality uses \(z_j\in\{0,1\}\), and the last uses \cref{obj:ass:coefficient-mass}. Consequently,
\[
\mathbb E_Z[Y_i(Z)^2]\le\mathbb E_Z[1]=1.
\]
The in-degree bound therefore yields
\[
\begin{aligned}
\sum_{i\in V}
\bigl|\{j:(j,i)\in G\}\bigr|
\,\mathbb E_Z[Y_i(Z)^2]
&\le
\sum_{i\in V}\bigl|\{j:(j,i)\in G\}\bigr|\\
&\le\sum_{i\in V}d
=nd.
\end{aligned}
\]
Since \(n>0\), \(q>0\), and \(q\le1\),
\[
\frac{4(1-q)}{n^2q}\ge0.
\]
Combining the preceding bounds gives
\[
\begin{aligned}
\mathbb E_{P_{\theta,q}}
\left[(T_q-\tau(\theta))^2\right]
&\le
\frac{5(d+1)^2}{n}
+\frac{4(1-q)}{n^2q}\,nd\\
&=
\frac{5(d+1)^2}{n}
+\frac{4d(1-q)}{nq}
=u(n,d,q).
\end{aligned}
\]
All expectations here are integrable because the assignment--audit space is finite.
% lean: CausalSmith.Experimentation.ThinnedgraphAdditiveRiskFrontier.abs_potentialOutcome_le_one
% lean: CausalSmith.Experimentation.ThinnedgraphAdditiveRiskFrontier.auditScore_squared_error_le

\item
To control clipping, first bound the target. Define the constant assignments by
\[
\mathbf1=(1)_{i\in V},
\qquad
\mathbf0=(0)_{i\in V}.
\]
For every \(i\in V\),
\[
Y_i(\mathbf1)-Y_i(\mathbf0)
=t_i+\sum_{j:(j,i)\in G}b_{ij},
\]
so
\[
\begin{aligned}
|Y_i(\mathbf1)-Y_i(\mathbf0)|
&\le |t_i|+
\left|\sum_{j:(j,i)\in G}b_{ij}\right|\\
&\le |t_i|+\sum_{j:(j,i)\in G}|b_{ij}|\\
&\le1.
\end{aligned}
\]
The last inequality follows from \cref{obj:ass:coefficient-mass} and \(|a_i|\ge0\). Averaging these contrasts gives
\[
\begin{aligned}
|\tau(\theta)|
&=
\frac1n
\left|\sum_{i\in V}
\bigl(Y_i(\mathbf1)-Y_i(\mathbf0)\bigr)\right|\\
&\le
\frac1n\sum_{i\in V}
|Y_i(\mathbf1)-Y_i(\mathbf0)|\\
&\le\frac1n\sum_{i\in V}1
=1.
\end{aligned}
\]
Here \(n\ge4\) ensures \(n>0\).
% lean: CausalSmith.Experimentation.ThinnedgraphAdditiveRiskFrontier.abs_tte_le_one

\item
For arbitrary real values with domains
\[
v_{\mathrm{clip}}\in\mathbb R,
\qquad
s_{\mathrm{target}}\in[-1,1],
\]
clipping satisfies
\[
\left|\max\{-1,\min\{1,v_{\mathrm{clip}}\}\}
-s_{\mathrm{target}}\right|
=
\begin{cases}
s_{\mathrm{target}}+1,
&v_{\mathrm{clip}}<-1,\\
1-s_{\mathrm{target}},
&v_{\mathrm{clip}}>1,\\
|v_{\mathrm{clip}}-s_{\mathrm{target}}|,
&-1\le v_{\mathrm{clip}}\le1.
\end{cases}
\]
In the first case,
\[
s_{\mathrm{target}}+1
\le s_{\mathrm{target}}-v_{\mathrm{clip}}
=|v_{\mathrm{clip}}-s_{\mathrm{target}}|,
\]
and in the second,
\[
1-s_{\mathrm{target}}
\le v_{\mathrm{clip}}-s_{\mathrm{target}}
=|v_{\mathrm{clip}}-s_{\mathrm{target}}|.
\]
The third case gives equality. Squaring therefore proves
\[
\bigl(\max\{-1,\min\{1,v_{\mathrm{clip}}\}\}
-s_{\mathrm{target}}\bigr)^2
\le
(v_{\mathrm{clip}}-s_{\mathrm{target}})^2.
\]
Apply this inequality pointwise with \(v_{\mathrm{clip}}=T_q(O)\) and \(s_{\mathrm{target}}=\tau(\theta)\). By the selected branch of \cref{obj:def:upper-rule},
\[
\begin{aligned}
\mathcal L(T^{\mathrm{up}};\theta,q)
&\le
\mathbb E_{P_{\theta,q}}
\left[(T_q-\tau(\theta))^2\right]\\
&\le u(n,d,q)
=\min\{1,u(n,d,q)\}.
\end{aligned}
\]
The nonnegative real expectation agrees with the extended nonnegative loss used in the risk definition.
% lean: CausalSmith.Experimentation.ThinnedgraphAdditiveRiskFrontier.clipped_sq_error_le
% lean: CausalSmith.Experimentation.ThinnedgraphAdditiveRiskFrontier.observable_upper

\item
Now suppose the branch condition fails:
\[
\neg\bigl(q>0\ \text{and}\ u(n,d,q)<1\bigr).
\]
If \(q>0\), this implies \(u(n,d,q)\ge1\). If \(q\) is not positive, the hypothesis \(q\in[0,1]\) gives \(q=0\), and \cref{obj:def:upper-envelope} gives
\[
u(n,d,0)=+\infty\ge1.
\]
Thus in either case
\[
\min\{1,u(n,d,q)\}=1,
\qquad
T^{\mathrm{up}}(O)=0.
\]
The target bound derived above then yields
\[
\begin{aligned}
\mathcal L(T^{\mathrm{up}};\theta,q)
&=\mathbb E_{P_{\theta,q}}[\tau(\theta)^2]\\
&\le\mathbb E_{P_{\theta,q}}[1]
=1
=\min\{1,u(n,d,q)\}.
\end{aligned}
\]
The equality \(\mathbb E_{P_{\theta,q}}[1]=1\) uses the probability normalization of the product design.
% lean: CausalSmith.Experimentation.ThinnedgraphAdditiveRiskFrontier.observable_upper

\item
Both branches establish the same bound for every \(\theta\in\mathcal M_{n,d}\). Taking the supremum gives
\[
\sup_{\theta\in\mathcal M_{n,d}}
\mathcal L(T^{\mathrm{up}};\theta,q)
\le\min\{1,u(n,d,q)\}.
\]
Together with the admissibility inequality established at the outset, this proves
\[
R_{n,d}(q)
\le
\sup_{\theta\in\mathcal M_{n,d}}
\mathcal L(T^{\mathrm{up}};\theta,q)
\le
\min\{1,u(n,d,q)\}.
\]
% lean: Causalean.Stat.worstCaseRiskENNReal_le
\end{enumerate}
\end{proof}

\begin{proof}[Proof of \cref{obj:thm:observed-record-precision-frontier}]
\begin{enumerate}
\item \textit{Upper risk bound.}
Fix admissible \(n,d,q\). Define the association-revelation probability and a numerical constant by
\[
p=1-(1-q)^d,
\qquad
\alpha_{\mathrm{ret}}=1-e^{-1}\in(0,1).
\]
We first establish
\[
\alpha_{\mathrm{ret}}\min\{1,dq\}\le p\le\min\{1,dq\}.
\]
For the upper bound, induction on the nonnegative integer degree starts with
\(1-(1-q)^0=0\). The induction step is
\[
\begin{aligned}
1-(1-q)^{d+1}
&=q+(1-q)\bigl[1-(1-q)^d\bigr]\\
&\le q+(1-q)dq
=(d+1)q-dq^2
\le(d+1)q.
\end{aligned}
\]
Also \(p\le1\), since \(1-q\ge0\).

For the lower bound, the exponential inequality \(1-q\le e^{-q}\) gives
\[
(1-q)^d\le e^{-dq},
\qquad
p\ge1-e^{-dq}.
\]
If \(dq\le1\), convexity of the exponential yields
\[
e^{-dq}\le(1-dq)e^0+dq\,e^{-1},
\qquad
1-e^{-dq}\ge\alpha_{\mathrm{ret}}dq.
\]
If \(dq>1\), monotonicity gives
\[
e^{-dq}\le e^{-1},
\qquad
1-e^{-dq}\ge\alpha_{\mathrm{ret}}.
\]
In particular, \(p>0\) whenever \(q>0\).

By \cref{obj:thm:observable-upper,obj:def:frontier-rule},
\[
R_{n,d}(q)
\le
\sup_{\theta\in\mathcal M_{n,d}}
\mathcal L(T^\star_{n,d,q};\theta,q)
\le\min\{1,u(n,d,q)\},
\]
and the attaining rule is measurable and real-valued. For \(q>0\), all denominators below are positive, and the probability bounds imply
\[
\frac{d^2}{n}\le\frac{d^2}{np},
\qquad
\frac{d}{nq}\le\frac{d^2}{np}.
\]
Using \(d+1\le2d\) and \(1-q\le1\) in
\cref{obj:def:upper-envelope}, we obtain
\[
\begin{aligned}
u(n,d,q)
&=\frac{5(d+1)^2}{n}+\frac{4d(1-q)}{nq}\\
&\le20\frac{d^2}{n}+4\frac{d}{nq}
\le24\frac{d^2}{np}.
\end{aligned}
\]
If \(d^2/(np)\le1\), then
\[
\min\{1,u(n,d,q)\}
\le u(n,d,q)
\le24\frac{d^2}{np}
=24F(n,d,q).
\]
If \(d^2/(np)>1\), then
\[
\min\{1,u(n,d,q)\}\le1\le24=24F(n,d,q).
\]
At \(q=0\), the envelope is \(+\infty\) and the scale is \(1\), so
\[
\min\{1,u(n,d,0)\}=1\le24F(n,d,0).
\]
Consequently, throughout the admissible range,
\[
R_{n,d}(q)
\le
\sup_{\theta\in\mathcal M_{n,d}}
\mathcal L(T^\star_{n,d,q};\theta,q)
\le24F(n,d,q).
\]
% lean: CausalSmith.Experimentation.ThinnedgraphAdditiveRiskFrontier.frontierEstimator_risk_upper

\item \textit{Lower risk bound.}
The definition gives \(F(n,d,q)\le1\). First suppose that
\(d^2/n\ge1/16\). The canonical product design satisfies
\cref{obj:ass:assignment,obj:ass:audit,obj:ass:design-independence}, so
\cref{obj:lem:supplied-block-record-lower} supplies
\[
\begin{aligned}
R_{n,d}(q)
&\ge\frac{1}{160000\pi^2}
       \min\left\{1,\frac{(d+1)^2}{n}\right\}\\
&\ge\frac{1}{2560000\pi^2}
\ge\frac{F(n,d,q)}{2560000\pi^2}.
\end{aligned}
\]
Here the second inequality follows from
\[
\frac{(d+1)^2}{n}\ge\frac{d^2}{n}\ge\frac1{16}.
\]

Now suppose that \(d^2/n<1/16\). Since \(d\ge1\),
\[
4d\le16d^2<n.
\]
Define the block count and source count by
\[
B=\left\lfloor\frac{n}{2d}\right\rfloor,
\qquad
m=Bd.
\]
The floor inequalities give
\[
2m\le n<2m+2d.
\]
Together with \(4d<n\), these imply
\[
B\ge1,
\qquad
\frac n4\le m\le\frac n2,
\qquad
2Bd\le n.
\]
Use the fixed source and recipient construction of \cref{obj:synth_4}.

Define the testing scale, numerical constant, and chosen amplitude by
\[
s=
\begin{cases}
\min\{1,d/\sqrt{mp}\},&p>0,\\
1,&p=0,
\end{cases}
\qquad
\kappa=\frac1{100\pi},
\qquad
h=\kappa s.
\]
Since \(m,d>0\), the definition gives \(0<s\le1\). Moreover,
\[
0<h\le\kappa<\frac14,
\]
where the last inequality follows from \(\pi>3\).

We next compare the risk and testing scales:
\[
F(n,d,q)\le s^2.
\]
For \(q>0\), we have \(p>0\) and
\[
\left(\frac d{\sqrt{mp}}\right)^2=\frac{d^2}{mp},
\qquad
\frac{d^2}{np}\le\frac{d^2}{mp}.
\]
If \(d/\sqrt{mp}\le1\), then
\[
F(n,d,q)\le\frac{d^2}{np}\le\frac{d^2}{mp}=s^2.
\]
If \(d/\sqrt{mp}>1\), then \(s^2=1\ge F(n,d,q)\).
For \(q=0\), \(p=0\) and \(F(n,d,0)=s^2=1\).
Thus the comparison covers the entire retention range.

Take the opposite-sign priors of \cref{obj:def:block-prior} at amplitude \(h\).
By \cref{obj:lem:block-law}, they are supported almost surely in
\(\mathcal M_{n,d}\), with constant targets \(\pm\Delta\), where
\[
\Delta=\frac{mh}{n}>0.
\]
The record maps are jointly measurable: the partition and binary coordinates
are finite, and the outcomes depend affinely on the real baseline coordinates.
The block and amplitude conditions established above allow
\cref{obj:thm:block-testing-scale} to give
\[
\operatorname{TV}(P_{+1,h},P_{-1,h})\le\frac12.
\]
Applying \cref{obj:lem:full-record-two-prior-risk} therefore yields
\[
R_{n,d}(q)
\ge\frac{\Delta^2}{2}
       \bigl(1-\operatorname{TV}(P_{+1,h},P_{-1,h})\bigr)
\ge\frac{\Delta^2}{4}.
\]
Finally, \(m/n\ge1/4\) implies
\[
\frac{\kappa s}{4}
\le\frac{m\kappa s}{n}
=\Delta.
\]
Hence
\[
\begin{aligned}
\frac{F(n,d,q)}{2560000\pi^2}
&\le\frac{s^2}{2560000\pi^2}
\le\frac{s^2}{640000\pi^2}\\
&=\frac14\left(\frac{\kappa s}{4}\right)^2
\le\frac{\Delta^2}{4}
\le R_{n,d}(q).
\end{aligned}
\]
This proves the lower bound in both degree regimes, including \(q=0\).
% lean: CausalSmith.Experimentation.ThinnedgraphAdditiveRiskFrontier.frontier_minimax_lower

\item \textit{Necessity of the sequence conditions.}
Fix admissible sequences \(d_n,q_n\), and suppose that measurable randomized
estimators \(T_n\) satisfy
\[
\rho_{\mathrm{risk},n}
:=
\sup_{\theta\in\mathcal M_{n,d_n}}
\mathcal L(T_n;\theta,q_n)
\longrightarrow0,
\qquad n\ge4.
\]
The preceding lower bound and the definition of minimax risk give
\[
0\le
\frac{F(n,d_n,q_n)}{2560000\pi^2}
\le R_{n,d_n}(q_n)
\le\rho_{\mathrm{risk},n}.
\]
Squeezing to zero and multiplying by the positive constant
\(2560000\pi^2\), we obtain
\[
F(n,d_n,q_n)\longrightarrow0.
\]

To extract the sequence conditions, we establish the endpoint and
positive-retention comparisons used here. Directly from
\cref{obj:def:frontier-scale}, since \(d\ge1\),
\[
F(n,d,0)=1,
\qquad
F(n,d,1)=\min\left\{1,\frac{d^2}{n}\right\}.
\]
For \(0<q\le1\), the established bounds on \(p\) give
\[
\frac{d^2}{n}\le\frac{d^2}{np},
\qquad
\frac{d}{nq}\le\frac{d^2}{np}.
\]
The lower probability estimate also gives, by splitting at \(dq=1\),
\[
\alpha_{\mathrm{ret}}\frac{d^2}{np}
\le
\begin{cases}
\dfrac{d}{nq},&dq\le1,\\[4pt]
\dfrac{d^2}{n},&dq>1.
\end{cases}
\]
Consequently,
\[
\frac{d^2}{n}+\frac{d}{nq}\le2\frac{d^2}{np},
\qquad
\alpha_{\mathrm{ret}}\frac{d^2}{np}
\le\frac{d^2}{n}+\frac{d}{nq}.
\]

For the lower truncated comparison, if \(d^2/(np)\ge1\), then
\[
\frac12\min\left\{1,\frac{d^2}{n}+\frac{d}{nq}\right\}
\le1=F(n,d,q).
\]
If \(d^2/(np)<1\), then
\[
\frac12\min\left\{1,\frac{d^2}{n}+\frac{d}{nq}\right\}
\le\frac12\left(\frac{d^2}{n}+\frac{d}{nq}\right)
\le\frac{d^2}{np}
=F(n,d,q).
\]
For the upper comparison, \(0<\alpha_{\mathrm{ret}}<1\) gives both
\[
\alpha_{\mathrm{ret}}F(n,d,q)\le1,
\qquad
\alpha_{\mathrm{ret}}F(n,d,q)
\le\alpha_{\mathrm{ret}}\frac{d^2}{np}
\le\frac{d^2}{n}+\frac{d}{nq}.
\]
Combining these inequalities proves
\[
\frac12\min\left\{1,\frac{d^2}{n}+\frac{d}{nq}\right\}
\le F(n,d,q)
\le
\alpha_{\mathrm{ret}}^{-1}
\min\left\{1,\frac{d^2}{n}+\frac{d}{nq}\right\}.
\]
% lean: CausalSmith.Experimentation.ThinnedgraphAdditiveRiskFrontier.frontier_elbow

Since \(F(n,d_n,q_n)\to0\), eventually
\[
F(n,d_n,q_n)<\frac12.
\]
The zero-retention identity then forces \(q_n>0\) eventually. On that tail,
the lower comparison implies
\[
\frac{d_n^2}{n}+\frac{d_n}{nq_n}<1:
\]
otherwise its truncated value would be \(1\), forcing
\(F(n,d_n,q_n)\ge1/2\). Removing the truncation now yields
\[
0\le\frac{d_n^2}{n}
\le\frac{d_n^2}{n}+\frac{d_n}{nq_n}
\le2F(n,d_n,q_n),
\]
and
\[
0\le\frac{d_n}{nq_n}
\le\frac{d_n^2}{n}+\frac{d_n}{nq_n}
\le2F(n,d_n,q_n).
\]
Both real ratios therefore tend to zero.

Define the extended audit ratio for every \(n\ge4\) by
\[
\rho_{\mathrm{audit},n}
=
\begin{cases}
+\infty,&q_n=0,\\[2pt]
\dfrac{d_n}{nq_n},&q_n>0.
\end{cases}
\]
It agrees eventually with the second real ratio, so
\[
\frac{d_n^2}{n}\longrightarrow0,
\qquad
\rho_{\mathrm{audit},n}\longrightarrow0.
\]
% lean: CausalSmith.Experimentation.ThinnedgraphAdditiveRiskFrontier.frontierScale_tendsto_iff

\item \textit{Sufficiency of the sequence conditions.}
Conversely, suppose
\[
\frac{d_n^2}{n}\longrightarrow0,
\qquad
\rho_{\mathrm{audit},n}\longrightarrow0.
\]
Convergence of the extended ratio gives
\(\rho_{\mathrm{audit},n}<1\) eventually. Thus \(q_n>0\) eventually,
and on that tail its finite real value satisfies
\[
\frac{d_n}{nq_n}\longrightarrow0.
\]
It follows that
\[
\frac{d_n^2}{n}+\frac{d_n}{nq_n}\longrightarrow0.
\]
The upper scale comparison proved above gives, eventually,
\[
0\le F(n,d_n,q_n)
\le\alpha_{\mathrm{ret}}^{-1}
\left(\frac{d_n^2}{n}+\frac{d_n}{nq_n}\right),
\]
and hence \(F(n,d_n,q_n)\to0\).

Choose the measurable rule sequence
\[
T_n(O,\xi)=T^\star_{n,d_n,q_n}(O),
\qquad n\ge4.
\]
The upper risk bound then supplies
\[
0\le
\sup_{\theta\in\mathcal M_{n,d_n}}
\mathcal L(T_n;\theta,q_n)
\le24F(n,d_n,q_n)
\longrightarrow0.
\]
Together with necessity, this proves the asserted equivalence.
% lean: CausalSmith.Experimentation.ThinnedgraphAdditiveRiskFrontier.precision_frontier_explicit

\item \textit{Universal constants and endpoint scales.}
Set
\[
c=\frac1{2560000\pi^2},
\qquad
C_0=24,
\qquad
c_1=\frac12,
\qquad
C_1=(1-e^{-1})^{-1}.
\]
These constants are positive because \(\pi>0\) and \(e^{-1}<1\).
The upper and lower risk bounds establish the risk chain with \(c,C_0\);
the sequence argument establishes consistency in both directions.
The endpoint identities and positive-retention comparison established above
give the remaining assertions with \(c_1,C_1\).
% lean: CausalSmith.Experimentation.ThinnedgraphAdditiveRiskFrontier.observed_record_precision_frontier
\end{enumerate}
\end{proof}

\begin{proof}[Proof of \cref{obj:lem:published-model}]
\begin{enumerate}
\item For a directed graph \(G'\) on \(V\), write
\[
\mathcal N^{\mathrm{in}}_i(G')
:=\{j\in V:(j,i)\in G'\},
\qquad
\mathcal N^{\mathrm{out}}_j(G')
:=\{i\in V:(j,i)\in G'\}.
\]
For a published schedule \(\vartheta=(G^\vartheta,c^\vartheta)\) and a binary assignment \(z\in\{0,1\}^V\), its polynomial outcome is
\[
Y_i^{\mathrm{pub}}(\vartheta,z)
:=
\sum_{J\subseteq\mathcal N^{\mathrm{in}}_i(G^\vartheta)}
c^\vartheta_{i,J}\prod_{j\in J}z_j.
\]
The published conventions of
\citep[Sections 3.1--3.2, Assumptions 1--2, equations (3.2)--(3.3)]{CortezRodriguezEichhornYu2023SNIPE}
give, for an interaction-order-one schedule,
\[
Y_i^{\mathrm{pub}}(\vartheta,z)
=
c^\vartheta_{i,\varnothing}
+\sum_{j\in\mathcal N^{\mathrm{in}}_i(G^\vartheta)}
c^\vartheta_{i,\{j\}}z_j,
\qquad
Y_{\max}(\vartheta)
=
\max_{i\in V}
\left(
|c^\vartheta_{i,\varnothing}|
+\sum_{j\in\mathcal N^{\mathrm{in}}_i(G^\vartheta)}
|c^\vartheta_{i,\{j\}}|
\right).
\]
These conventions permit a zero singleton coefficient on a graph arrow.

Let \(\theta^+\) denote the augmentation of an additive schedule \(\theta=(G,a,t,b)\):
\[
\theta^+:=(\widetilde G,c^{\mathrm{pub}}),
\qquad
\widetilde G:=G\cup\{(i,i):i\in V\}.
\]
Its singleton coefficients satisfy
\[
c^{\mathrm{pub}}_{i,\{j\}}
=
\begin{cases}
t_i,&j=i,\\
b_{ij},&j\ne i\text{ and }(j,i)\in G,\\
0,&j\ne i\text{ and }(j,i)\notin G.
\end{cases}
\]
The empty-set coefficient is \(a_i\). If
\(J\not\subseteq\mathcal N^{\mathrm{in}}_i(\widetilde G)\), then \(J\ne\varnothing\).
When \(J=\{j\}\), membership outside this neighborhood forces both
\(j\ne i\) and \((j,i)\notin G\), so its coefficient is zero.
When \(J\) is not a singleton, all singleton contributions in the coefficient construction vanish. Likewise, \(|J|>1\) forces \(J\) to be neither empty nor a singleton. Thus
\[
c^{\mathrm{pub}}_{i,J}=0
\quad\text{if}\quad
J\not\subseteq\mathcal N^{\mathrm{in}}_i(\widetilde G)
\ \text{or}\ |J|>1.
\]
This proves the interaction-order-one restriction.

Since \(G\) is off-diagonal,
\[
\mathcal N^{\mathrm{in}}_i(\widetilde G)
=\{i\}\,\dot\cup\,\mathcal N^{\mathrm{in}}_i(G),
\qquad
\mathcal N^{\mathrm{out}}_j(\widetilde G)
=\{j\}\,\dot\cup\,\mathcal N^{\mathrm{out}}_j(G).
\]
Splitting the coefficient sum at the own coordinate therefore gives
\[
|c^{\mathrm{pub}}_{i,\varnothing}|
+\sum_{j\in\mathcal N^{\mathrm{in}}_i(\widetilde G)}
|c^{\mathrm{pub}}_{i,\{j\}}|
=
|a_i|+|t_i|
+\sum_{j\in\mathcal N^{\mathrm{in}}_i(G)}|b_{ij}|.
\]
For \(\theta\in\mathcal M_{n,d}\), the coefficient-mass requirement in
\cref{obj:def:model} bounds each right-hand side by one. Taking the maximum proves
\[
Y_{\max}(\theta^+)\le1.
\]
The same disjoint neighborhood decompositions and the two degree bounds in
\cref{obj:def:model} give
\[
|\mathcal N^{\mathrm{in}}_i(\widetilde G)|
=1+|\mathcal N^{\mathrm{in}}_i(G)|\le d+1,
\qquad
|\mathcal N^{\mathrm{out}}_j(\widetilde G)|
=1+|\mathcal N^{\mathrm{out}}_j(G)|\le d+1.
\]
Every diagonal arrow belongs to \(\widetilde G\) by construction. These identities also cover \(d=0\), with empty off-diagonal neighborhoods and zero neighborhood sums.
% lean: CausalSmith.Experimentation.ThinnedgraphAdditiveRiskFrontier.augment_publishedClass

\item Conversely, let \(\vartheta=(G^\vartheta,c^\vartheta)\) satisfy the stated published restrictions and contain every diagonal arrow. Define its stripped schedule by
\[
\vartheta^-:=(G^-,a^-,t^-,b^-),
\qquad
G^-:=\{(j,i)\in G^\vartheta:j\ne i\},
\]
\[
a_i^-:=c^\vartheta_{i,\varnothing},
\qquad
t_i^-:=c^\vartheta_{i,\{i\}},
\qquad
b_{ij}^-:=c^\vartheta_{i,\{j\}}
\quad(i,j\in V).
\]
The graph \(G^-\) is off-diagonal. A source in a published in-neighborhood is either the recipient itself or a distinct source retained in \(G^-\). The analogous partition holds for recipients in an out-neighborhood. Consequently,
\[
\mathcal N^{\mathrm{in}}_i(G^\vartheta)
=\{i\}\,\dot\cup\,\mathcal N^{\mathrm{in}}_i(G^-),
\qquad
\mathcal N^{\mathrm{out}}_j(G^\vartheta)
=\{j\}\,\dot\cup\,\mathcal N^{\mathrm{out}}_j(G^-).
\]
Hence
\[
1+|\mathcal N^{\mathrm{in}}_i(G^-)|\le d+1,
\qquad
1+|\mathcal N^{\mathrm{out}}_j(G^-)|\le d+1.
\]
Cancelling one gives both degree bounds by \(d\). The published envelope bounds each recipient's coefficient mass, and splitting off the diagonal term yields
\[
\begin{aligned}
|a_i^-|+|t_i^-|
+\sum_{j\in\mathcal N^{\mathrm{in}}_i(G^-)}|b_{ij}^-|
&=
|c^\vartheta_{i,\varnothing}|
+\sum_{j\in\mathcal N^{\mathrm{in}}_i(G^\vartheta)}
|c^\vartheta_{i,\{j\}}|\\
&\le Y_{\max}(\vartheta)\le1.
\end{aligned}
\]
These are precisely the membership requirements in \cref{obj:def:model}, so
\[
\vartheta^-\in\mathcal M_{n,d}.
\]
% lean: CausalSmith.Experimentation.ThinnedgraphAdditiveRiskFrontier.strip_scheduleClass

\item Re-augmenting \(\vartheta^-\) recovers the published graph:
\[
G^-\cup\{(i,i):i\in V\}=G^\vartheta.
\]
Indeed, for \(j=i\), both graphs contain the arrow by the diagonal hypothesis; for \(j\ne i\), both contain it exactly when \((j,i)\in G^\vartheta\).

Write \(c^{(\vartheta^-)^+}\) for the coefficients of the re-augmented schedule. For every recipient \(i\) and every
\(J\subseteq\mathcal N^{\mathrm{in}}_i(G^\vartheta)\), we claim
\[
c^{(\vartheta^-)^+}_{i,J}=c^\vartheta_{i,J}.
\]
If \(J=\varnothing\), both coefficients equal \(a_i^-\).
If \(|J|>1\), both vanish by their order-one restrictions.
In the remaining case, \(J\ne\varnothing\) and \(|J|\le1\), so \(J=\{j\}\) for some \(j\). If \(j=i\), the re-augmented coefficient is
\(t_i^-=c^\vartheta_{i,\{i\}}\). If \(j\ne i\), the neighborhood condition gives
\((j,i)\in G^\vartheta\), hence \((j,i)\in G^-\), and the re-augmented coefficient is
\(b_{ij}^-=c^\vartheta_{i,\{j\}}\).
This proves recovery of every coefficient on a subset of the recipient's in-neighborhood.
% lean: CausalSmith.Experimentation.ThinnedgraphAdditiveRiskFrontier.augment_strip_equivalent

\item For any additive schedule, stripping its augmentation recovers its graph because
\[
(j,i)\in G\cup\{(k,k):k\in V\},\quad j\ne i
\quad\Longleftrightarrow\quad
(j,i)\in G,
\]
using the off-diagonal property of \(G\). The empty and own-singleton coefficient identities, followed by the off-diagonal singleton identity on graph arrows, give
\[
a^{(\theta^+)^-}=a,\qquad
t^{(\theta^+)^-}=t,\qquad
b^{(\theta^+)^-}_{ij}=b_{ij}
\quad\text{for every }(j,i)\in G.
\]
Together with graph recovery, these are the asserted additive-schedule correspondence.
% lean: CausalSmith.Experimentation.ThinnedgraphAdditiveRiskFrontier.strip_augment_equivalent

\item For any additive schedule, the order-one outcome formula and the disjoint augmented in-neighborhood give, for every \(i\in V\) and \(z\in\{0,1\}^V\),
\[
\begin{aligned}
Y_i^{\mathrm{pub}}(\theta^+,z)
&=
c^{\mathrm{pub}}_{i,\varnothing}
+\sum_{j\in\mathcal N^{\mathrm{in}}_i(\widetilde G)}
c^{\mathrm{pub}}_{i,\{j\}}z_j\\
&=
a_i+t_i z_i
+\sum_{j\in\mathcal N^{\mathrm{in}}_i(G)}b_{ij}z_j\\
&=Y_i(z).
\end{aligned}
\]
The second equality separates the own coordinate; every remaining source is distinct from \(i\) and lies on an arrow of \(G\), so its singleton coefficient is \(b_{ij}\).
% lean: CausalSmith.Experimentation.ThinnedgraphAdditiveRiskFrontier.pubOutcome_augment

\item Define the published population contrast by
\[
\tau^{\mathrm{pub}}(\vartheta)
:=
\frac1n\sum_{i\in V}
\left(
Y_i^{\mathrm{pub}}(\vartheta,\mathbf1)
-Y_i^{\mathrm{pub}}(\vartheta,\mathbf0)
\right),
\qquad
\mathbf1_j:=1,\quad \mathbf0_j:=0
\quad(j\in V).
\]
Applying the outcome identity at these two assignments gives
\[
\tau^{\mathrm{pub}}(\theta^+)
=
\frac1n\sum_{i\in V}\bigl(Y_i(\mathbf1)-Y_i(\mathbf0)\bigr)
=\tau(\theta).
\]
% lean: CausalSmith.Experimentation.ThinnedgraphAdditiveRiskFrontier.pubTte_augment

\item Let
\[
\operatorname{diag}(V):=\{(i,i):i\in V\},
\qquad
W_{ij}:=W(j,i)\quad(j\ne i),
\]
where \(W_{ij}\) is the audit mark for the arrow from \(j\) to \(i\).
For each assignment and audit realization, the retained graph and complete record are
\[
H:=\{(j,i)\in G:W_{ij}=1\},
\qquad
O:=\bigl(H,Z,(Y_i(Z))_{i\in V}\bigr).
\]
Define the record transformation
\[
\mathcal A(O)
:=
\bigl(H\cup\operatorname{diag}(V),Z,(Y_i(Z))_{i\in V}\bigr).
\]
On every off-diagonal pair,
\[
\mathbf1\{(j,i)\in\widetilde G\}\,W_{ij}
=
\mathbf1\{(j,i)\in G\}\,W_{ij}.
\]
Thus the retained off-diagonal graphs agree. The assignment coordinate agrees identically, and the outcome coordinates agree by the outcome identity already proved. Therefore, writing \(\widetilde O_{\theta^+}\) for the published record,
\[
\widetilde O_{\theta^+}=\mathcal A(O).
\]
On the space of records whose graph contains every diagonal arrow, the inverse transformation is
\[
\mathcal A^{-1}
\bigl(\widetilde H,Z,(Y_i(Z))_{i\in V}\bigr)
=
\bigl(\widetilde H\setminus\operatorname{diag}(V),
Z,(Y_i(Z))_{i\in V}\bigr).
\]
Adding and removing these known coordinates are measurable operations and are inverse to one another.
% lean: CausalSmith.Experimentation.ThinnedgraphAdditiveRiskFrontier.pubRecordOf_augment

\item Fix the common assignment-and-audit law satisfying
\cref{obj:ass:assignment,obj:ass:audit,obj:ass:design-independence},
and let \(\mathfrak C\) be a published schedule class such that
\[
\theta\in\mathcal M_{n,d}\quad\Longrightarrow\quad
\theta^+\in\mathfrak C.
\]
For any published schedule \(\vartheta\), its record under this channel is
\[
\widetilde O_\vartheta
:=
\left(
\{(j,i)\in G^\vartheta:j\ne i,\ W_{ij}=1\}
\cup\operatorname{diag}(V),
\ Z,\
(Y_i^{\mathrm{pub}}(\vartheta,Z))_{i\in V}
\right).
\]
As in \cref{obj:def:risk}, use an independent randomization seed with
\[
\xi\sim\operatorname{Unif}[0,1],
\qquad
\xi\mathbin{\perp\!\!\!\perp}(Z,W).
\]
For a measurable published estimator \(T^{\mathrm{pub}}\), define its squared risk and the corresponding additive-record estimator by
\[
\mathcal L^{\mathrm{pub}}(T^{\mathrm{pub}};\vartheta,q)
:=
\mathbb E\!\left[
\bigl(T^{\mathrm{pub}}(\widetilde O_\vartheta,\xi)
-\tau^{\mathrm{pub}}(\vartheta)\bigr)^2
\right],
\qquad
T(O,\xi):=T^{\mathrm{pub}}(\mathcal A(O),\xi).
\]
The latter estimator is measurable because \(\mathcal A\) is measurable. For every additive schedule, the record and target identities give the pointwise equality
\[
\bigl(T(O,\xi)-\tau(\theta)\bigr)^2
=
\bigl(T^{\mathrm{pub}}(\widetilde O_{\theta^+},\xi)
-\tau^{\mathrm{pub}}(\theta^+)\bigr)^2.
\]
Taking expectations under the same assignment, audit, and seed law therefore gives
\[
\mathcal L(T;\theta,q)
=
\mathcal L^{\mathrm{pub}}(T^{\mathrm{pub}};\theta^+,q).
\]
These expectations are understood as nonnegative extended integrals, so the identity also covers infinite risks.

For each \(\theta\in\mathcal M_{n,d}\), class inclusion now yields
\[
\mathcal L(T;\theta,q)
\le
\sup_{\vartheta\in\mathfrak C}
\mathcal L^{\mathrm{pub}}(T^{\mathrm{pub}};\vartheta,q).
\]
Taking the supremum over \(\theta\), and then using the infimum in
\cref{obj:def:risk}, gives
\[
R_{n,d}(q)
\le
\sup_{\theta\in\mathcal M_{n,d}}\mathcal L(T;\theta,q)
\le
\sup_{\vartheta\in\mathfrak C}
\mathcal L^{\mathrm{pub}}(T^{\mathrm{pub}};\vartheta,q).
\]
This holds for every measurable published estimator. Taking their infimum proves
\[
R_{n,d}(q)
\le
\inf_{T^{\mathrm{pub}}}
\sup_{\vartheta\in\mathfrak C}
\mathcal L^{\mathrm{pub}}(T^{\mathrm{pub}};\vartheta,q),
\]
which is the claimed minimax risk transfer under the common observation channel.
% lean: CausalSmith.Experimentation.ThinnedgraphAdditiveRiskFrontier.minimaxRisk_le_pubMinimaxRisk
\end{enumerate}
\end{proof}

\begin{proof}[Proof of \cref{obj:thm:supported-boundaries}]
\begin{enumerate}
\item Choose the universal constants
\[
c:=\frac{1}{2560000\pi^2},
\qquad
C_0:=21.
\]
Both are positive. Fix admissible \(n,d,q\) and a joint assignment--audit law satisfying the stated assumptions. Independence and the prescribed marginals give
\[
\operatorname{Law}(Z,W)
=
\left(\bigotimes_{j\in V}\operatorname{Bernoulli}(1/2)\right)
\otimes
\left(\bigotimes_{(j,i)\in V^2:j\ne i}
\operatorname{Bernoulli}(q)\right).
\]
Thus the experiment has the independent assignment--audit product design. % lean: CausalSmith.Experimentation.ThinnedgraphAdditiveRiskFrontier.design_eq_thinnedDesign
Applying \cref{obj:thm:observable-upper}, with the envelope defined in \cref{obj:def:upper-envelope}, yields
\[
R_{n,d}(q)\le \min\{1,u(n,d,q)\}.
\]
% lean: CausalSmith.Experimentation.ThinnedgraphAdditiveRiskFrontier.observable_upper

\item By \cref{obj:lem:supplied-block-record-lower},
\[
R_{n,d}(q)\ge
\frac{1}{160000\pi^2}
\min\left\{1,\frac{(d+1)^2}{n}\right\}.
\]
Since
\[
0<c=\frac{1}{2560000\pi^2}
\le \frac{1}{160000\pi^2},
\qquad
\min\left\{1,\frac{(d+1)^2}{n}\right\}\ge0,
\]
we obtain
\[
R_{n,d}(q)\ge
c\min\left\{1,\frac{(d+1)^2}{n}\right\}.
\]
% lean: CausalSmith.Experimentation.ThinnedgraphAdditiveRiskFrontier.supported_boundaries

\item We next establish the degree-one frontier bound
\[
R_{n,1}(q)\ge cF(n,1,q),
\]
where the risk scale of \cref{obj:def:frontier-scale} specializes to
\[
F(n,1,q)=
\begin{cases}
\min\{1,1/(nq)\},&q>0,\\
1,&q=0.
\end{cases}
\qquad
0<F(n,1,q)\le1.
\]
The auxiliary degree \(1\) is admissible because \(n\ge4\).

First suppose \(1/n\ge1/16\). Then
\[
\frac4n\ge\frac1{16},
\qquad
\min\left\{1,\frac4n\right\}\ge\frac1{16}.
\]
Using \cref{obj:lem:supplied-block-record-lower} at degree one and
\(2560000=16\cdot160000\), we conclude
\[
cF(n,1,q)
\le c
\le
\frac{1}{160000\pi^2}\min\left\{1,\frac4n\right\}
\le R_{n,1}(q).
\]

Now suppose \(1/n<1/16\). Then \(n>16\), and in particular \(4<n\). Define the integer block count and source count by
\[
B:=\left\lfloor\frac n2\right\rfloor,
\qquad
m:=B.
\]
The division remainder satisfies
\[
0\le n-2B<2,
\qquad
2m\le n<2m+2.
\]
Consequently \(B\ge1\), and, since \(2<n/2\),
\[
n<2m+2<2m+\frac n2
\quad\Longrightarrow\quad
\frac n4\le m.
\]
Thus
\[
B\ge1,\qquad 2m\le n,\qquad \frac n4\le m\le n.
\]
% lean: CausalSmith.Experimentation.ThinnedgraphAdditiveRiskFrontier.frontier_block_rounding

Define the degree-one testing scale and testing constant by
\[
s:=s(B,1,q)=
\begin{cases}
\min\{1,1/\sqrt{mq}\},&q>0,\\
1,&q=0,
\end{cases}
\qquad
\kappa:=\frac1{100\pi}.
\]
These are the quantities in \cref{obj:thm:block-testing-scale}, since the association-revelation probability at degree one is
\[
p:=1-(1-q)^1=q.
\]
For \(q>0\), both entries in the minimum defining \(s\) are positive; for \(q=0\), \(s=1\). Hence
\[
0<s\le1,\qquad \kappa>0.
\]
% lean: CausalSmith.Experimentation.ThinnedgraphAdditiveRiskFrontier.testingScale_pos_le_one

We claim that
\[
F(n,1,q)\le s^2.
\]
For \(q>0\), the positive denominators and \(m\le n\) imply
\[
\left(\frac1{\sqrt{mq}}\right)^2=\frac1{mq},
\qquad
\frac1{nq}\le\frac1{mq}.
\]
If \(1/\sqrt{mq}\le1\), then
\[
F(n,1,q)\le\frac1{nq}
\le\frac1{mq}=s^2.
\]
If \(1/\sqrt{mq}>1\), then
\[
F(n,1,q)\le1=s^2.
\]
At \(q=0\), the definitions give
\[
F(n,1,0)=1=s^2.
\]
This proves the claim over the entire retention range.
% lean: CausalSmith.Experimentation.ThinnedgraphAdditiveRiskFrontier.frontierScale_le_testingScale_sq

Set the testing amplitude and target magnitude to
\[
h:=\kappa s,
\qquad
\Delta:=\frac{mh}{n}.
\]
Since \(\pi>3\),
\[
0<h\le\kappa=\frac1{100\pi}<\frac14,
\qquad
\Delta>0.
\]
Take the opposite-sign block priors of \cref{obj:def:block-prior} with \(B\) blocks, degree one, and amplitude \(h\). The conditions \(B\ge1\), \(2B\le n\), and \(0<h\le1/4\) allow us to apply \cref{obj:lem:block-law}. It gives, under the respective priors,
\[
\theta\in\mathcal M_{n,1},
\qquad
\tau(\theta)=\Delta
\quad\text{and}\quad
\tau(\theta)=-\Delta
\quad\text{almost surely}.
\]
The block record maps are jointly measurable: their graph and assignment coordinates are finite, and their outcome coordinates are affine functions of the baselines for each fixed partition and assignment. Since \(h=\kappa s\), \cref{obj:thm:block-testing-scale} gives
\[
\operatorname{TV}(P_{+1,h},P_{-1,h})\le\frac12.
\]
Therefore \cref{obj:lem:full-record-two-prior-risk} yields
\[
R_{n,1}(q)
\ge
\frac{\Delta^2}{2}
\left(1-\operatorname{TV}(P_{+1,h},P_{-1,h})\right)
\ge
\frac{\Delta^2}{4}
=
\frac14\left(\frac{m\kappa s}{n}\right)^2.
\]
% lean: CausalSmith.Experimentation.ThinnedgraphAdditiveRiskFrontier.block_testing_minimax_lower

Finally, the rounding bound gives
\[
\frac mn\ge\frac14,
\qquad
\frac{\kappa s}{4}
\le\frac{m\kappa s}{n}.
\]
Both sides of the latter inequality are positive, so
\[
\left(\frac{\kappa s}{4}\right)^2
\le
\left(\frac{m\kappa s}{n}\right)^2,
\qquad
\frac14\left(\frac{\kappa s}{4}\right)^2
=
\frac{s^2}{640000\pi^2}.
\]
Combining these estimates with \(F(n,1,q)\le s^2\) gives
\[
\begin{aligned}
cF(n,1,q)
&\le \frac{s^2}{2560000\pi^2}\\
&\le \frac{s^2}{640000\pi^2}\\
&=\frac14\left(\frac{\kappa s}{4}\right)^2\\
&\le\frac14\left(\frac{m\kappa s}{n}\right)^2\\
&\le R_{n,1}(q).
\end{aligned}
\]
Together with the first case, this proves the degree-one frontier bound.
% lean: CausalSmith.Experimentation.ThinnedgraphAdditiveRiskFrontier.frontier_minimax_lower

\item The model class of \cref{obj:def:model} is increasing in its degree bound:
\[
\mathcal M_{n,1}\subseteq\mathcal M_{n,d}
\qquad(d\ge1).
\]
Indeed, the coefficient-mass condition is the same, and each degree bound at most one is also at most \(d\). For every measurable randomized estimator \(T\),
\[
\sup_{\theta\in\mathcal M_{n,1}}\mathcal L(T;\theta,q)
\le
\sup_{\theta\in\mathcal M_{n,d}}\mathcal L(T;\theta,q).
\]
Taking infima over the common estimator class gives
\[
R_{n,1}(q)\le R_{n,d}(q).
\]
% lean: CausalSmith.Experimentation.ThinnedgraphAdditiveRiskFrontier.minimaxRisk_mono_degree

We also have
\[
\frac1{1+nq}\le F(n,1,q).
\]
For \(q>0\), \(nq>0\), and
\[
\frac1{1+nq}\le1,
\qquad
\frac1{1+nq}\le\frac1{nq}.
\]
These two inequalities imply the claim by the degree-one formula for \(F\). At \(q=0\), both sides equal one. Multiplication by \(c>0\), followed by the bounds just proved, therefore gives
\[
\frac{c}{1+nq}
\le cF(n,1,q)
\le R_{n,1}(q)
\le R_{n,d}(q).
\]
% lean: CausalSmith.Experimentation.ThinnedgraphAdditiveRiskFrontier.degree_one_frontier_reciprocal

\item Combining the two lower bounds and using \(c>0\), we obtain
\[
\begin{aligned}
R_{n,d}(q)
&\ge
\max\left\{
c\min\left\{1,\frac{(d+1)^2}{n}\right\},
\frac{c}{1+nq}
\right\}\\
&=
c\max\left\{
\min\left\{1,\frac{(d+1)^2}{n}\right\},
\frac1{1+nq}
\right\}.
\end{aligned}
\]
% lean: CausalSmith.Experimentation.ThinnedgraphAdditiveRiskFrontier.supported_boundaries

\item At \(q=0\), the retention lower bound and the upper bound from \cref{obj:thm:observable-upper} give
\[
c\le R_{n,d}(0)\le1.
\]
% lean: CausalSmith.Experimentation.ThinnedgraphAdditiveRiskFrontier.supported_boundaries

\item At \(q=1\), the audit contribution in \cref{obj:def:upper-envelope} vanishes:
\[
u(n,d,1)=\frac{5(d+1)^2}{n}.
\]
Thus
\[
R_{n,d}(1)\le
\min\left\{1,\frac{5(d+1)^2}{n}\right\}.
\]
If \((d+1)^2/n\le1\), then
\[
\min\left\{1,\frac{5(d+1)^2}{n}\right\}
\le\frac{5(d+1)^2}{n}
\le
21\min\left\{1,\frac{(d+1)^2}{n}\right\}.
\]
If \((d+1)^2/n>1\), then
\[
\min\left\{1,\frac{5(d+1)^2}{n}\right\}
\le1\le21
=
21\min\left\{1,\frac{(d+1)^2}{n}\right\}.
\]
Together with the degree lower bound, these estimates prove
\[
c\min\left\{1,\frac{(d+1)^2}{n}\right\}
\le R_{n,d}(1)
\le
C_0\min\left\{1,\frac{(d+1)^2}{n}\right\}.
\]
% lean: CausalSmith.Experimentation.ThinnedgraphAdditiveRiskFrontier.supported_boundaries

\item It remains to bound the degree-one upper envelope. Suppose first that \(q>0\). Since \(q\le1\),
\[
u(n,1,q)
=
\frac{20}{n}+\frac{4(1-q)}{nq}
=
\frac{4+16q}{nq}
\le\frac{20}{nq}.
\]
Consequently,
\[
\min\{1,u(n,1,q)\}
\le
\min\left\{1,\frac{20}{nq}\right\}.
\]
If \(nq\le20\), then
\[
\min\left\{1,\frac{20}{nq}\right\}
\le1\le\frac{21}{1+nq}.
\]
If \(nq>20\), then the positive denominators and
\[
20(1+nq)\le21nq
\]
give
\[
\min\left\{1,\frac{20}{nq}\right\}
\le\frac{20}{nq}
\le\frac{21}{1+nq}.
\]
At \(q=0\), the envelope is \(+\infty\), so
\[
\min\{1,u(n,1,0)\}=1\le21=\frac{21}{1+nq}.
\]
Hence, throughout \(q\in[0,1]\),
\[
\min\{1,u(n,1,q)\}\le\frac{21}{1+nq}.
\]
Applying the upper bound from \cref{obj:thm:observable-upper} and the retention lower bound already obtained proves
\[
\frac{c}{1+nq}
\le R_{n,1}(q)
\le\frac{C_0}{1+nq}.
\]
% lean: CausalSmith.Experimentation.ThinnedgraphAdditiveRiskFrontier.degree_one_envelope_upper
\end{enumerate}
\end{proof}

\begin{proof}[Proof of \cref{obj:thm:block-testing-scale}]
Write
\[
m:=Bd,\qquad s:=s(B,d,q),\qquad
\kappa:=\frac1{100\pi},\qquad
\mathcal D(h):=\operatorname{TV}(P_{+1,h},P_{-1,h}).
\]
The design assumptions identify the joint assignment--audit law with the product of its specified marginals.

\begin{enumerate}
\item \textit{The small-amplitude domain.}
Since \(m,d>0\), the definition of the scale gives
\[
0<s\le1.
\]
Indeed, when \(p>0\), both entries in its defining minimum are positive; when \(p=0\), \(s=1\). Also, \(\pi>3\) gives
\[
0<\kappa<\frac14.
\]
Consequently, every amplitude satisfying \(0\le h\le\kappa s\) satisfies
\[
0\le h\le\kappa\le\frac14.
\]
Fix such an amplitude while proving the small-amplitude bound.
% lean: CausalSmith.Experimentation.ThinnedgraphAdditiveRiskFrontier.testing_small_amplitude_mem

\item \textit{High retention.}
First suppose \(p\ge1/2\). We obtain a comparison bound by conditioning on the source partition \(S\) from \cref{obj:def:block-prior}. For each source \(j\) and block \(\ell\), define
\[
X_j:=2Z_j-1\in\{-1,1\},
\qquad
X_{\ell,\mathrm{sum}}:=\sum_{j\in S_\ell}X_j,
\qquad
\mathcal J:=\bigcup_{\ell=1}^{B}S_\ell,
\quad |\mathcal J|=m.
\]
Here \(\mathcal J\) is the fixed set of labeled sources. Under prior sign \(\sigma\in\{-1,1\}\), the distinct response of block \(\ell\) is
\[
y_\ell=U_\ell+\frac{\sigma h}{2d}X_{\ell,\mathrm{sum}}.
\]
Thus, conditional on \(S,Z,W\), its response density is
\[
\lambda^S_{\sigma,h}(y\mid Z,W)
:=
\prod_{\ell=1}^{B}
f\!\left(y_\ell-\frac{\sigma h}{2d}X_{\ell,\mathrm{sum}}\right),
\qquad y\in\mathbb R^B.
\]
The law of \((Z,W)\) is common to both signs. The product-translation bound in
\cref{obj:lem:baseline-translation-affinity} gives
\[
\int_{\mathbb R^B}
\left(
\sqrt{\lambda^S_{+1,h}(y\mid Z,W)}
-\sqrt{\lambda^S_{-1,h}(y\mid Z,W)}
\right)^2\,dy
\le
\frac{4\pi^2h^2}{d^2}
\sum_{\ell=1}^{B}X_{\ell,\mathrm{sum}}^2.
\]
For every fixed partition, independent centered treatment signs satisfy
\[
\mathbb E_Z X_{\ell,\mathrm{sum}}^2
=
\sum_{j\in S_\ell}\mathbb E_Z X_j^2
+
\sum_{\substack{j,j'\in S_\ell\\j\ne j'}}
\mathbb E_Z[X_jX_{j'}]
=d.
\]
Averaging the preceding translation bound therefore yields
\[
\mathbb E_{Z,W}\!
\int_{\mathbb R^B}
\left(
\sqrt{\lambda^S_{+1,h}}
-\sqrt{\lambda^S_{-1,h}}
\right)^2\,dy
\le 4\pi^2h^2\frac Bd.
\]
Define the conditional complete-record laws by
\[
P^S_{\sigma,h}:=\operatorname{Law}_{\sigma,h}(H,Z,Y(Z)\mid S).
\]
The common-marginal Hellinger bound in
\cref{obj:lem:baseline-translation-affinity}, followed by the record-copying map supplied by
\cref{obj:lem:block-law}, gives
\[
\operatorname{TV}(P^S_{+1,h},P^S_{-1,h})
\le 2\pi h\sqrt{B/d}.
\]
The copying map preserves the retained graph, the complete assignment, and all response coordinates. For any measurable record event, its probability difference under the two mixtures is bounded by the mean of the conditional total variations. Hence
\[
\mathcal D(h)
\le
\mathbb E_S\operatorname{TV}(P^S_{+1,h},P^S_{-1,h})
\le 2\pi h\sqrt{B/d}.
\]

Since \(p>0\), the definition of \(s\) also gives
\[
h\sqrt{mp}\le\kappa d,
\qquad
h^2mp\le\kappa^2d^2.
\]
Using \(p\ge1/2\) and \(m=Bd\), we obtain
\[
h^2\frac Bd
=\frac{h^2m}{d^2}
\le2\kappa^2.
\]
The required constant calibration is
\[
\left(2\pi h\sqrt{B/d}\right)^2
\le8\pi^2\kappa^2
=\frac8{10000}
\le\frac14.
\]
The quantity being squared is nonnegative, so \(\mathcal D(h)\le1/2\).
% lean: CausalSmith.Experimentation.ThinnedgraphAdditiveRiskFrontier.testing_small_high_retention

\item \textit{Few sources.}
Suppose now \(p<1/2\) and \(m<32\). The conditional-partition comparison just derived applies throughout the small-amplitude domain. Since \(d\ge1\) and \(h\le\kappa\),
\[
\frac Bd=\frac{m}{d^2}\le32,
\qquad
h^2\frac Bd\le32\kappa^2.
\]
Therefore
\[
\left(2\pi h\sqrt{B/d}\right)^2
\le128\pi^2\kappa^2
=\frac{128}{10000}
\le\frac14,
\]
and again \(\mathcal D(h)\le1/2\).
% lean: CausalSmith.Experimentation.ThinnedgraphAdditiveRiskFrontier.testing_small_few_sources

\item \textit{Remaining capacity at low retention.}
It remains to consider \(p<1/2\) and \(m\ge32\). For the retained graph \(H\), define
\[
\rho_j(H):=\mathbf1\{\exists i:(j,i)\in H\},
\qquad j\in\mathcal J,
\]
and, on the support of its actual marginal, define
\[
\begin{aligned}
r_\ell(H)&:=
\bigl|\{j\in\mathcal J:\exists i\in C_\ell,\ (j,i)\in H\}\bigr|,\\
u_\ell(H)&:=d-r_\ell(H),\\
M(H)&:=\sum_{\ell=1}^{B}u_\ell(H),\\
M_{\mathrm{unrev}}(H)&:=\sum_{j\in\mathcal J}(1-\rho_j(H)).
\end{aligned}
\]
Every unrevealed source occupies a remaining source slot in its true block. Thus, for almost every retained graph,
\[
M_{\mathrm{unrev}}(H)\le M(H).
\]
By \cref{obj:lem:block-law}, the variables \(\rho_j\) are independent Bernoulli variables with success probability \(p\). Define the capacity event
\[
\mathcal E_{\mathrm{cap}}:=\{H:M(H)\ge m/4\}.
\]
Its complement satisfies
\[
\Pr(\mathcal E_{\mathrm{cap}}^c)
\le\Pr(M_{\mathrm{unrev}}<m/4).
\]
Independence gives the inverse-third moment identity
\[
\mathbb E_H[3^{-M_{\mathrm{unrev}}}]
=
\prod_{j\in\mathcal J}\mathbb E_H[3^{-(1-\rho_j)}]
=
\left(p+\frac{1-p}{3}\right)^m
\le\left(\frac23\right)^m.
\]
Since \(M_{\mathrm{unrev}}<m/4\) implies
\(3^{-M_{\mathrm{unrev}}}\ge3^{-m/4}\), Markov's inequality yields
\[
\Pr(\mathcal E_{\mathrm{cap}}^c)
\le
3^{m/4}\left(\frac23\right)^m.
\]
The fourth power of this nonnegative upper bound satisfies
\[
\left(3^{m/4}\left(\frac23\right)^m\right)^4
=
\left(\frac{16}{27}\right)^m
\le
\left(\frac{16}{27}\right)^{32}
\le
\left(\frac1{16}\right)^4.
\]
The last comparison is equivalent to
\(16^{36}\le27^{32}\), which follows from \(8^{48}\le9^{48}\). Taking fourth roots proves
\[
\Pr(\mathcal E_{\mathrm{cap}}^c)\le\frac1{16}.
\]
% lean: CausalSmith.Experimentation.ThinnedgraphAdditiveRiskFrontier.hidden_count_good_event_of_retention_half

\item \textit{The good-event chi-squared bound.}
Use the Walsh coefficients and energies of
\cref{obj:lem:walsh-channel-energy}. Explicitly, for \(w\in\mathbb R\) and integers \(1\le k\le d\), these are
\[
g_d(w):=
2^{-d}\sum_{\varepsilon\in\{-1,1\}^d}
f\!\left(w-\frac h{2d}\sum_{j=1}^{d}\varepsilon_j\right),
\]
\[
a_k(w):=
\begin{cases}
\displaystyle
\frac{2^{-d}}{g_d(w)}
\sum_{\varepsilon\in\{-1,1\}^d}
f\!\left(w-\frac h{2d}\sum_{j=1}^{d}\varepsilon_j\right)
\prod_{j=1}^{k}\varepsilon_j,
&g_d(w)>0,\\
0,&g_d(w)=0,
\end{cases}
\]
and
\[
\gamma_k:=\int_{\mathbb R}a_k(w)^2g_d(w)\,dw,
\qquad
\eta:=\sum_{k=1}^{d}\binom dk\gamma_k,
\qquad
\eta_1:=\sum_{k=1}^{d}k\binom dk\gamma_k.
\]
The cited energy bound asserts
\[
0\le\eta\le\eta_1\le\frac{12\pi^2h^2}{d}.
\]
In the present amplitude range,
\[
h^2mp\le\kappa^2d^2.
\]
For \(p=0\), this follows because its left side is zero. For \(p>0\), it follows by squaring
\(h\sqrt{mp}\le\kappa d\), as above. Consequently,
\[
\eta
\le\eta_1
\le\frac{12\pi^2\kappa^2}{d}
\le\frac3{2500},
\]
and
\[
0\le Bp\eta_1
\le
\frac{12\pi^2h^2Bp}{d}
=
\frac{12\pi^2h^2mp}{d^2}
\le\frac3{2500}.
\]

All conditional quantities in this step and the next are used on the finite support of the retained-graph marginal, where \(\Pr\{H\}>0\). Define the reference law and conditional likelihood ratios by
\[
Q_H(dZ,dy):=
\operatorname{Law}(Z)(dZ)\prod_{\ell=1}^{B}g_d(y_\ell)\,dy,
\]
\[
L_{\sigma,H}(Z,y):=
\begin{cases}
\displaystyle
\frac{\lambda_{\sigma,h}(y\mid H,Z)}
{\prod_{\ell=1}^{B}g_d(y_\ell)},
&\prod_{\ell=1}^{B}g_d(y_\ell)>0,\\
0,&\prod_{\ell=1}^{B}g_d(y_\ell)=0,
\end{cases}
\qquad
\chi_H^2:=\mathbb E_{Q_H}[(L_{+1,H}-1)^2].
\]
Here \(\lambda_{\sigma,h}\) is the conditional density supplied by
\cref{obj:lem:block-law}. Independence makes the conditional assignment marginal equal to \(\operatorname{Law}(Z)\).

Put \(e:=\exp(1)\). The standard exponential bounds give
\[
e<3,\qquad
4e\eta<\frac1{50},\qquad
eBp\eta_1\le\frac1{100}.
\]
In particular, \(4e\eta<1\), as required by
\cref{obj:lem:hidden-allocation-contraction}. The exponential series and the geometric series give
\[
\exp(eBp\eta_1)
\le\exp(1/100)
\le\sum_{k=0}^{\infty}(1/100)^k
=\frac{100}{99}
\le\frac{51}{50}.
\]
The contraction bound therefore yields
\[
\begin{aligned}
\mathbb E_H[\mathbf1_{\mathcal E_{\mathrm{cap}}}\chi_H^2]
&\le
\frac{\exp(eBp\eta_1)}{1-4e\eta}-1\\
&\le
\frac{51/50}{49/50}-1
=\frac2{49}
\le\frac1{16}.
\end{aligned}
\]
% lean: CausalSmith.Experimentation.ThinnedgraphAdditiveRiskFrontier.testing_small_good_chiSq_bound

\item \textit{From conditional chi-squared divergence to the complete record.}
The baseline density is even. The conditional-density formula in
\cref{obj:lem:block-law} and the definition of \(g_d\) consequently give
\[
\lambda_{-1,h}(y\mid H,Z)
=
\lambda_{+1,h}(-y\mid H,Z),
\qquad
g_d(-w)=g_d(w).
\]
Thus response reflection \((Z,y)\mapsto(Z,-y)\) preserves \(Q_H\) and exchanges the two conditional laws. In particular,
\[
\mathbb E_{Q_H}[(L_{-1,H}-1)^2]=\chi_H^2.
\]
Each component translation in a block satisfies
\[
f\!\left(w-\frac h{2d}\sum_{j=1}^{d}\varepsilon_j\right)
\le2^d g_d(w).
\]
Averaging over the finite partition mixture and conditioning on \(H\) therefore gives
\[
0\le L_{\sigma,H}\le\frac{2^{Bd}}{\Pr\{H\}}
\qquad Q_H\text{-almost everywhere}.
\]
This also ensures that both squared likelihood deviations are integrable.

For either conditional sign law, the density formula for total variation and Cauchy--Schwarz give
\[
\operatorname{TV}\!\left(
\operatorname{Law}_{\sigma,h}(Z,y\mid H),Q_H
\right)
=
\frac12\mathbb E_{Q_H}|L_{\sigma,H}-1|
\le\frac12\sqrt{\chi_H^2}.
\]
The triangle inequality hence implies
\[
\operatorname{TV}\!\left(
\operatorname{Law}_{+1,h}(Z,y\mid H),
\operatorname{Law}_{-1,h}(Z,y\mid H)
\right)
\le\sqrt{\chi_H^2}.
\]
Both signs have the same retained-graph marginal. Averaging conditional event differences, using the reconstruction in
\cref{obj:lem:block-law}, bounding conditional total variation by one on the complement of the capacity event, and applying Cauchy--Schwarz on that event yield
\[
\begin{aligned}
\mathcal D(h)
&\le
\Pr(\mathcal E_{\mathrm{cap}}^c)
+
\mathbb E_H[
\mathbf1_{\mathcal E_{\mathrm{cap}}}\sqrt{\chi_H^2}]\\
&\le
\Pr(\mathcal E_{\mathrm{cap}}^c)
+
\sqrt{\mathbb E_H[
\mathbf1_{\mathcal E_{\mathrm{cap}}}\chi_H^2]}\\
&\le\frac1{16}+\frac14
\le\frac12.
\end{aligned}
\]
This completes the small-amplitude bound in all three cases, including \(p=0\).
% lean: CausalSmith.Experimentation.ThinnedgraphAdditiveRiskFrontier.blockMixture_tvDist_good_event_chiSq_le

\item \textit{An observable reverse statistic.}
Now fix \(0<h\le1/4\) and \(p>0\). Choose the first labeled recipient in each block,
\[
i_\ell:=\min C_\ell\in C_\ell,
\]
which exists because \(d\ge1\). Define the revealed source-sign sum and the scalar statistic by
\[
A_\ell(H,Z):=
\sum_{\substack{j\in\mathcal J\\
\exists i\in C_\ell,\ (j,i)\in H}}X_j,
\qquad
V_{\mathrm{test}}(O):=
\sum_{\ell=1}^{B}A_\ell(H,Z)Y_{i_\ell}(Z).
\]
The response-copying property gives \(Y_{i_\ell}(Z)=y_\ell\). Thus this statistic is measurable from the complete record, and projection gives
\[
\mathcal D(h)\ge
\operatorname{TV}\!\left(
\operatorname{Law}_{P_{+1,h}}(V_{\mathrm{test}}),
\operatorname{Law}_{P_{-1,h}}(V_{\mathrm{test}})
\right).
\]

For a fixed true partition, the reveal indicator of each source depends on its \(d\) outgoing audit marks. These edge sets are disjoint across sources, so the reveal indicators have independent Bernoulli\((p)\) laws. They are independent of treatment signs and baselines. The statistic consequently has the law obtained by drawing the genuine partition, independent treatment signs, independent reveal indicators, and independent baselines, and then forming
\[
A_\ell=\sum_{j\in S_\ell}\rho_jX_j,
\qquad
V_{\mathrm{test}}
=
\sum_{\ell=1}^{B}
A_\ell\left(
U_\ell+\frac{\sigma h}{2d}X_{\ell,\mathrm{sum}}
\right).
\]
Mixing this identity over the partition and baselines gives the scalar law in the preceding projection inequality.
% lean: CausalSmith.Experimentation.ThinnedgraphAdditiveRiskFrontier.reverseTestStatistic_mixture_law

\item \textit{Integrating the baselines.}
Define the positive test center, variance budget, and finite signal sum by
\[
a_{\mathrm{test}}:=\frac{hmp}{2d}>0,
\qquad
v_{\mathrm{test}}:=\frac{7mp}{64},
\qquad
V_{\mathrm{sig}}:=\sum_{\ell=1}^{B}A_\ell X_{\ell,\mathrm{sum}}.
\]
The baseline density is even and supported on \([-1/4,1/4]\), so
\[
\mathbb E U_\ell=0,
\qquad
\mathbb E U_\ell^2\le\frac1{16}.
\]
For fixed partition, signs, and reveals, independence of the baselines gives
\[
\mathbb E_U\!\left[\left(\sum_{\ell=1}^{B}A_\ell U_\ell\right)^2\right]
=
\sum_{\ell=1}^{B}A_\ell^2\mathbb E U_\ell^2
\le\frac1{16}\sum_{\ell=1}^{B}A_\ell^2.
\]
Moreover,
\[
V_{\mathrm{test}}-\sigma a_{\mathrm{test}}
=
\sum_{\ell=1}^{B}A_\ell U_\ell
+\frac{\sigma h}{2d}(V_{\mathrm{sig}}-mp).
\]
The mixed term has zero baseline expectation. Hence
\[
\mathbb E_U[
(V_{\mathrm{test}}-\sigma a_{\mathrm{test}})^2
\mid S,Z,\rho]
\le
\frac1{16}\sum_{\ell=1}^{B}A_\ell^2
+\frac{h^2}{4d^2}(V_{\mathrm{sig}}-mp)^2.
\]
All required second moments exist because the sums are finite and the baselines have bounded support.
% lean: CausalSmith.Experimentation.ThinnedgraphAdditiveRiskFrontier.reverseLatentLaw_baseline_energy_le

\item \textit{The finite source-row moments.}
Fix a partition and a reveal pattern. For each block define
\[
\mathcal J_{\ell,\mathrm{rev}}
:=\{j\in S_\ell:\rho_j=1\},
\qquad
r_\ell:=|\mathcal J_{\ell,\mathrm{rev}}|\in\{0,\ldots,d\}.
\]
This agrees with the retained-graph count already defined. Independence and centering of the signs give
\[
\mathbb E_Z[A_\ell^2\mid\rho,S]=r_\ell,
\qquad
\mathbb E_Z[A_\ell X_{\ell,\mathrm{sum}}\mid\rho,S]=r_\ell:
\]
in each product expansion, terms with distinct source indices have zero expectation, and each surviving diagonal term equals one.

For deterministic real weights \((\alpha_j)_{j\in S_\ell}\in\mathbb R^d\), expansion of the fourth power leaves precisely the terms in which every sign index occurs an even number of times. Therefore
\[
\begin{aligned}
\mathbb E_Z\left(\sum_{j\in S_\ell}\alpha_jX_j\right)^4
&=
\sum_{j\in S_\ell}\alpha_j^4
+
6\sum_{\substack{j,j'\in S_\ell\\j<j'}}
\alpha_j^2\alpha_{j'}^2\\
&=
3\left(\sum_{j\in S_\ell}\alpha_j^2\right)^2
-2\sum_{j\in S_\ell}\alpha_j^4.
\end{aligned}
\]
Applying this identity to the revealed sum, the full sum, and their sum and difference gives
\[
\begin{aligned}
\mathbb E_Z[A_\ell^4\mid\rho,S]
&=3r_\ell^2-2r_\ell,\\
\mathbb E_Z[X_{\ell,\mathrm{sum}}^4\mid\rho,S]
&=3d^2-2d,\\
\mathbb E_Z[(A_\ell+X_{\ell,\mathrm{sum}})^4\mid\rho,S]
&=3(d+3r_\ell)^2-2(d+15r_\ell),\\
\mathbb E_Z[(A_\ell-X_{\ell,\mathrm{sum}})^4\mid\rho,S]
&=3(d-r_\ell)^2-2(d-r_\ell).
\end{aligned}
\]
For the last two formulas, the weights are respectively \(2,1\) and \(0,-1\) on the revealed and unrevealed coordinates. Using the polynomial identity
\[
(A_\ell+X_{\ell,\mathrm{sum}})^4
+(A_\ell-X_{\ell,\mathrm{sum}})^4
=
2A_\ell^4+12A_\ell^2X_{\ell,\mathrm{sum}}^2
+2X_{\ell,\mathrm{sum}}^4
\]
and substituting the four expectations yields
\[
\mathbb E_Z[
A_\ell^2X_{\ell,\mathrm{sum}}^2\mid\rho,S]
=
r_\ell d+2r_\ell(r_\ell-1)
\le3r_\ell d.
\]
These formulas include \(r_\ell=0\), with empty sums equal to zero.

Since the reveal indicators have mean \(p\),
\[
\mathbb E_\rho[r_\ell\mid S]=dp.
\]
Averaging the row identities and inequality over reveals consequently gives
\[
\mathbb E[A_\ell^2\mid S]=dp,
\qquad
\mathbb E[A_\ell X_{\ell,\mathrm{sum}}\mid S]=dp,
\qquad
\mathbb E[A_\ell^2X_{\ell,\mathrm{sum}}^2\mid S]
\le3d^2p.
\]
% lean: CausalSmith.Experimentation.ThinnedgraphAdditiveRiskFrontier.reverse_revealed_total_mixed_moment

\item \textit{The squared-deviation budget.}
For a fixed partition, different row signals depend on disjoint collections of independent sign--reveal pairs. Their variances therefore add, and the preceding row moments give
\[
\mathbb E[V_{\mathrm{sig}}\mid S]=Bdp=mp,
\]
\[
\begin{aligned}
\mathbb E[(V_{\mathrm{sig}}-mp)^2\mid S]
&=
\sum_{\ell=1}^{B}
\operatorname{Var}(A_\ell X_{\ell,\mathrm{sum}}\mid S)\\
&\le
\sum_{\ell=1}^{B}
\mathbb E[A_\ell^2X_{\ell,\mathrm{sum}}^2\mid S]
\le3Bd^2p.
\end{aligned}
\]
The baseline contribution satisfies
\[
\mathbb E\!\left[\frac1{16}\sum_{\ell=1}^{B}A_\ell^2
\ \middle|\ S\right]
=\frac{Bdp}{16}.
\]
Substitution into the baseline-integrated bound gives, for every partition,
\[
\begin{aligned}
\mathbb E_\sigma[
(V_{\mathrm{test}}-\sigma a_{\mathrm{test}})^2\mid S]
&\le
\frac{Bdp}{16}
+\frac{h^2}{4d^2}(3Bd^2p)\\
&=
\frac{mp}{16}+\frac{3h^2Bp}{4}\\
&\le
\frac{mp}{16}+\frac{3mp}{64}
=\frac{7mp}{64}.
\end{aligned}
\]
The last inequality uses \(h^2\le1/16\) and \(B\le Bd=m\). Averaging over the genuine partition proves, for both signs,
\[
\mathbb E_\sigma[
(V_{\mathrm{test}}-\sigma a_{\mathrm{test}})^2]
\le v_{\mathrm{test}}.
\]
% lean: CausalSmith.Experimentation.ThinnedgraphAdditiveRiskFrontier.reverseSignalEnergy_average_le

\item \textit{The scalar sign test.}
Assign nonnegative values of \(V_{\mathrm{test}}\) to the positive sign. Since \(a_{\mathrm{test}}>0\), the pointwise inequalities
\[
a_{\mathrm{test}}^2\mathbf1\{V_{\mathrm{test}}<0\}
\le(V_{\mathrm{test}}-a_{\mathrm{test}})^2,
\]
\[
a_{\mathrm{test}}^2\mathbf1\{V_{\mathrm{test}}\ge0\}
\le(V_{\mathrm{test}}+a_{\mathrm{test}})^2
\]
bound the two error probabilities by
\[
P_{+1,h}(V_{\mathrm{test}}<0)
\le\frac{v_{\mathrm{test}}}{a_{\mathrm{test}}^2},
\qquad
P_{-1,h}(V_{\mathrm{test}}\ge0)
\le\frac{v_{\mathrm{test}}}{a_{\mathrm{test}}^2}.
\]
For the event \(\{V_{\mathrm{test}}<0\}\), the definition of total variation gives
\[
1-\operatorname{TV}\!\left(
\operatorname{Law}_{P_{+1,h}}(V_{\mathrm{test}}),
\operatorname{Law}_{P_{-1,h}}(V_{\mathrm{test}})
\right)
\le
P_{+1,h}(V_{\mathrm{test}}<0)
+
P_{-1,h}(V_{\mathrm{test}}\ge0).
\]
Combining this with the projection inequality and simplifying the positive denominators yields
\[
\begin{aligned}
\mathcal D(h)
&\ge1-\frac{2v_{\mathrm{test}}}{a_{\mathrm{test}}^2}\\
&=
1-\frac{7d^2}{8h^2mp}.
\end{aligned}
\]
This is the asserted reverse bound.
% lean: CausalSmith.Experimentation.ThinnedgraphAdditiveRiskFrontier.reverse_test_tv_bound

\item \textit{The testing radius.}
Define its admissible set by
\[
\mathcal A_{\mathrm{test}}
:=
\{h\in[0,1/4]:\mathcal D(h)\le1/2\},
\qquad
\bar h:=\sup\mathcal A_{\mathrm{test}}.
\]
The inequalities \(0<s\le1\) and \(0<\kappa<1/4\), together with the small-amplitude bound, show
\[
0\le\kappa s\le\frac14,
\qquad
\mathcal D(\kappa s)\le\frac12.
\]
Thus \(\kappa s\in\mathcal A_{\mathrm{test}}\). This set is nonempty and bounded above by \(1/4\), so
\[
\kappa s\le\bar h.
\]

To prove the upper bound, take any \(h\in\mathcal A_{\mathrm{test}}\). If \(s=1\), then
\[
h\le\frac14\le2s.
\]
If \(s\ne1\), its definition implies
\[
p>0,\qquad s=\frac d{\sqrt{mp}}.
\]
Suppose, for a contradiction, that \(h>2s\). Positivity permits multiplication and squaring, giving
\[
2d<h\sqrt{mp},
\qquad
4d^2<h^2mp.
\]
In particular \(h>0\), and
\[
7d^2<4h^2mp,
\qquad
\frac{7d^2}{8h^2mp}<\frac12.
\]
The reverse bound then gives
\[
\mathcal D(h)
\ge1-\frac{7d^2}{8h^2mp}
>\frac12,
\]
contradicting membership in \(\mathcal A_{\mathrm{test}}\). Every member of that set is therefore at most \(2s\), and taking its supremum proves
\[
\bar h\le2s.
\]
Together with the lower bound, this establishes the final assertion.
% lean: CausalSmith.Experimentation.ThinnedgraphAdditiveRiskFrontier.testing_radius_of_bounds
\end{enumerate}
\end{proof}

\begin{proof}[Proof of \cref{obj:thm:nonlocal-testing-answer}]
\begin{enumerate}
\item The canonical design is the product of the treatment-assignment and audit-mark laws. On the finite population \(V\), with \(|V|=n\), its marginals therefore satisfy
\[
\operatorname{Law}(Z)
=\bigotimes_{j\in V}\operatorname{Bernoulli}(1/2),
\qquad
\operatorname{Law}(W)
=\bigotimes_{\substack{(j,i)\in V^2\\j\ne i}}
\operatorname{Bernoulli}(q).
\]
Indeed, \(q\in[0,1]\) ensures that both factors are probability laws, and projecting their product onto either coordinate gives the corresponding factor. Thus the design satisfies \cref{obj:ass:assignment,obj:ass:audit}.
% lean: CausalSmith.Experimentation.ThinnedgraphAdditiveRiskFrontier.thinnedDesign_assignment
% lean: CausalSmith.Experimentation.ThinnedgraphAdditiveRiskFrontier.thinnedDesign_audit
The two coordinate projections of a product probability law are independent, so
\[
W\mathbin{\perp\!\!\!\perp}Z.
\]
This verifies \cref{obj:ass:design-independence}.
% lean: CausalSmith.Experimentation.ThinnedgraphAdditiveRiskFrontier.thinnedDesign_independent

\item Apply \cref{obj:thm:block-testing-scale} to this design. Its population, block, degree, capacity, and retention hypotheses are precisely the hypotheses here, and the preceding step verifies its three design hypotheses. Its complete-record mixture laws, testing scale, and testing radius coincide with \(P_{+1,h},P_{-1,h},s,\bar h\) in the present statement. Consequently, it supplies
\[
(100\pi)^{-1}s\le\bar h,
\qquad
\bar h\le2s,
\]
as well as
\[
\forall h\in\mathbb R,\qquad
0\le h\le(100\pi)^{-1}s
\ \Longrightarrow\
\operatorname{TV}(P_{+1,h},P_{-1,h})\le\frac12.
\]
% lean: CausalSmith.Experimentation.ThinnedgraphAdditiveRiskFrontier.block_testing_scale

\item Suppose \(d=2\). Expanding the definition of the association-revelation probability gives
\[
m=2B,
\qquad
p=1-(1-q)^2=2q-q^2.
\]
First suppose \(q>0\). Since \(q\le1\) and \(B\ge1\),
\[
p=q(2-q)>0,
\qquad
2Bp>0,
\qquad
\sqrt{2Bp}>0.
\]
The definition of the testing scale therefore gives
\[
s=\min\left\{1,\frac{2}{\sqrt{2Bp}}\right\}.
\]
To simplify the second entry, square it:
\[
\left(\frac{2}{\sqrt{2Bp}}\right)^2
=\frac{4}{2Bp}
=\frac{2}{Bp}.
\]
Both the quotient and the square root below are nonnegative, so
\[
\frac{2}{\sqrt{2Bp}}
=\sqrt{\frac{2}{Bp}},
\qquad
s=\min\left\{1,\sqrt{\frac{2}{B(2q-q^2)}}\right\}.
\]
In the remaining case \(q\le0\), the hypothesis \(q\ge0\) gives
\[
q=0,\qquad p=0,\qquad s=1.
\]
These two cases establish the stated degree-two formula.
% lean: CausalSmith.Experimentation.ThinnedgraphAdditiveRiskFrontier.nonlocal_testing_answer
\end{enumerate}
\end{proof}

\begin{proof}[Proof of \cref{obj:thm:frontier-answer}]
\begin{enumerate}
\item Fix admissible \(n,d,q\). We first establish the upper bound. Define
\[
p_k(q):=1-(1-q)^k
\quad(k\in\mathbb N,\ q\in[0,1]),
\qquad
p:=p_d(q),
\qquad
c_{\mathrm{exp}}:=1-e^{-1}\in(0,1).
\]
The reveal-probability bounds are
\[
c_{\mathrm{exp}}\min\{1,dq\}\le p\le\min\{1,dq\}.
\]
For the upper bound, \(p\le1\), and induction gives \(p_k(q)\le kq\): the initial value is \(p_0(q)=0\), and
\[
p_{k+1}(q)
=q+(1-q)p_k(q)
\le q+(1-q)kq
\le(k+1)q.
\]
For the lower bound, the exponential inequality \(1-q\le e^{-q}\) yields
\[
(1-q)^d\le e^{-dq}.
\]
When \(dq\le1\), convexity of the exponential gives
\[
e^{-dq}\le(1-dq)+dq\,e^{-1},
\qquad
p\ge c_{\mathrm{exp}}dq.
\]
When \(dq>1\), monotonicity gives
\[
e^{-dq}\le e^{-1},
\qquad
p\ge c_{\mathrm{exp}}.
\]
These prove the claimed reveal-probability bounds; in particular, \(p>0\) when \(q>0\).
% lean: CausalSmith.Experimentation.ThinnedgraphAdditiveRiskFrontier.reveal_probability_ge_chord

For admissible parameters with \(q>0\), define
\[
v_{\mathrm{F}}(n,d,q):=\frac{d^2}{np}.
\]
Since \(p\le1\) and \(p\le dq\),
\[
\frac{d^2}{n}\le v_{\mathrm{F}}(n,d,q),
\qquad
\frac{d}{nq}\le v_{\mathrm{F}}(n,d,q).
\]
Using \(d+1\le2d\) and \(1-q\le1\), the envelope in
\cref{obj:def:upper-envelope} therefore satisfies
\[
u(n,d,q)
=\frac{5(d+1)^2}{n}+\frac{4d(1-q)}{nq}
\le20\frac{d^2}{n}+4\frac{d}{nq}
\le24v_{\mathrm{F}}(n,d,q).
\]
If \(v_{\mathrm{F}}(n,d,q)\le1\), this implies
\[
\min\{1,u(n,d,q)\}\le24v_{\mathrm{F}}(n,d,q)=24F(n,d,q).
\]
If \(v_{\mathrm{F}}(n,d,q)>1\), then
\[
\min\{1,u(n,d,q)\}\le1\le24=24F(n,d,q).
\]
At \(q=0\), the same comparison follows from
\[
u(n,d,0)=+\infty,
\qquad
F(n,d,0)=1.
\]
% lean: CausalSmith.Experimentation.ThinnedgraphAdditiveRiskFrontier.upperEnvelope_le_frontierScale

Use the seeded representative
\[
\widehat T^\star_{n,d,q}(O,\xi):=T^\star_{n,d,q}(O)
\qquad
(O\text{ a complete record},\ \xi\in\mathbb R).
\]
It is measurable by \cref{obj:thm:observable-upper} and
\cref{obj:def:frontier-rule,obj:def:upper-rule}. Those results and the envelope comparison give
\[
R_{n,d}(q)
\le
\sup_{\theta\in\mathcal M_{n,d}}
\mathcal L(\widehat T^\star_{n,d,q};\theta,q)
\le
\min\{1,u(n,d,q)\}
\le24F(n,d,q).
\]
% lean: CausalSmith.Experimentation.ThinnedgraphAdditiveRiskFrontier.frontierEstimator_risk_upper

\item For the lower bound, first suppose that
\[
\frac{d^2}{n}\ge\frac1{16}.
\]
The canonical product design satisfies the assignment, audit, and independence hypotheses of \cref{obj:lem:supplied-block-record-lower}, which supplies
\[
R_{n,d}(q)\ge
\frac1{160000\pi^2}
\min\left\{1,\frac{(d+1)^2}{n}\right\}.
\]
Here
\[
\min\left\{1,\frac{(d+1)^2}{n}\right\}\ge\frac1{16},
\qquad
F(n,d,q)\le1.
\]
Consequently,
\[
R_{n,d}(q)
\ge\frac1{2560000\pi^2}
\ge\frac{F(n,d,q)}{2560000\pi^2}.
\]
% lean: CausalSmith.Experimentation.ThinnedgraphAdditiveRiskFrontier.frontier_minimax_lower

\item Now suppose that
\[
\frac{d^2}{n}<\frac1{16}.
\]
Since \(d\ge1\), we have \(4d\le16d^2<n\). Define the block count and source count by
\[
B:=\left\lfloor\frac{n}{2d}\right\rfloor,
\qquad
m:=Bd.
\]
Floor rounding gives
\[
B\ge1,
\qquad
2m\le n<2m+2d,
\qquad
\frac n4\le m\le\frac n2.
\]
Indeed, \(n>4d\) ensures \(B\ge1\), and \(2d<n/2\) together with
\(n<2m+2d\) gives \(m>n/4\).
% lean: CausalSmith.Experimentation.ThinnedgraphAdditiveRiskFrontier.frontier_block_rounding

Define the testing scale on this parameter range by
\[
s:=s(B,d,q)=
\begin{cases}
\min\{1,d/\sqrt{mp}\},&p>0,\\
1,&p=0.
\end{cases}
\]
Thus \(0<s\le1\). When \(q>0\), we have \(p>0\), and
\[
\left(\frac d{\sqrt{mp}}\right)^2=\frac{d^2}{mp},
\qquad
\frac{d^2}{np}\le\frac{d^2}{mp}.
\]
If \(d/\sqrt{mp}\le1\), these identities imply
\[
F(n,d,q)\le\frac{d^2}{np}\le\frac{d^2}{mp}=s^2.
\]
If \(d/\sqrt{mp}>1\), then \(s=1\) and \(F(n,d,q)\le s^2\).
At \(q=0\), \(p=0\) and \(F(n,d,0)=s^2=1\). Hence, throughout the retention range,
\[
F(n,d,q)\le s^2.
\]
% lean: CausalSmith.Experimentation.ThinnedgraphAdditiveRiskFrontier.frontierScale_le_testingScale_sq

Set
\[
\kappa:=\frac1{100\pi},
\qquad
h:=\kappa s=h_\circ(B,d,q),
\qquad
\Delta:=\frac{mh}{n}.
\]
The positivity of \(m,s,\kappa\), and \(\pi>3\), give
\[
0<h\le\kappa<\frac14,
\qquad
\Delta>0.
\]
Consider the opposite-sign priors of \cref{obj:def:block-prior} at this amplitude. To specify their placement, take the sources to be
\[
\{1,\ldots,m\},
\qquad
C_\ell:=\{m+(\ell-1)d+1,\ldots,m+\ell d\}
\quad(1\le\ell\le B).
\]
The capacity bound \(2m\le n\) places these disjoint recipient blocks within the population. Each recipient has \(d\) incoming arrows, each source has \(d\) outgoing arrows, and the remaining degrees are zero.

The baseline density in \cref{obj:def:block-prior} vanishes outside
\([-1/4,1/4]\), so \(|U_\ell|\le1/4\) almost surely for every block. For either
\(\sigma\in\{-1,1\}\), each recipient \(i\in C_\ell\) consequently satisfies
\[
|a_i|+|t_i|+\sum_{j:(j,i)\in G}|b_{ij}|
=
\left|U_\ell-\frac{\sigma h}{2}\right|
+d\frac hd
\le\frac14+\frac{3h}{2}
\le1.
\]
Source and padding rows have coefficient mass zero. Thus both priors are supported on \(\mathcal M_{n,d}\). Summing the spillover contributions by source gives their constant targets:
\[
\tau(\theta)
=\frac1n\sum_{j=1}^{m}\sum_{i:(j,i)\in G}\frac{\sigma h}{d}
=\frac1n\sum_{j=1}^{m}\sigma h
=\frac{\sigma mh}{n}
=\sigma\Delta.
\]
% lean: CausalSmith.Experimentation.ThinnedgraphAdditiveRiskFrontier.block_support

The record maps are jointly measurable: their graph and assignment coordinates are finite, and their outcome coordinates are affine in the baselines. The canonical product design and \(2Bd\le n\) satisfy the hypotheses of \cref{obj:thm:block-testing-scale}, which supplies
\[
\operatorname{TV}(P_{+1,h},P_{-1,h})\le\frac12.
\]
Applying \cref{obj:lem:full-record-two-prior-risk} to these supported priors yields
\[
R_{n,d}(q)
\ge
\frac{\Delta^2}{2}
\left(1-\operatorname{TV}(P_{+1,h},P_{-1,h})\right)
\ge
\frac{\Delta^2}{4}
=
\frac14\left(\frac{m\kappa s}{n}\right)^2.
\]
This bound applies to the full class of measurable randomized estimators in
\cref{obj:def:risk}.
% lean: CausalSmith.Experimentation.ThinnedgraphAdditiveRiskFrontier.block_testing_minimax_lower

Finally, \(m/n\ge1/4\) gives
\[
\frac{\kappa s}{4}\le\frac{m\kappa s}{n},
\qquad
\frac14\left(\frac{\kappa s}{4}\right)^2
=\frac{s^2}{640000\pi^2}.
\]
Combining the preceding inequalities,
\[
\frac{F(n,d,q)}{2560000\pi^2}
\le
\frac{s^2}{2560000\pi^2}
\le
\frac{s^2}{640000\pi^2}
\le
\frac14\left(\frac{m\kappa s}{n}\right)^2
\le R_{n,d}(q).
\]
Together with the large-degree case and the upper bound, this proves the displayed risk chain.
% lean: CausalSmith.Experimentation.ThinnedgraphAdditiveRiskFrontier.frontier_minimax_lower

\item We next prove necessity in the consistency criterion. Let \(d_n,q_n\) be admissible sequences and suppose that measurable randomized estimators \(T_n\) satisfy
\[
\sup_{\theta\in\mathcal M_{n,d_n}}
\mathcal L(T_n;\theta,q_n)\longrightarrow0.
\]
Define the positive constant
\[
C_{\mathrm{low}}:=2560000\pi^2>0.
\]
For every \(n\ge4\), the lower bound and the definition of minimax risk give
\[
0\le
\frac{F(n,d_n,q_n)}{C_{\mathrm{low}}}
\le R_{n,d_n}(q_n)
\le
\sup_{\theta\in\mathcal M_{n,d_n}}
\mathcal L(T_n;\theta,q_n).
\]
Squeezing in the extended nonnegative reals and passing to real values at zero yields
\[
\frac{F(n,d_n,q_n)}{C_{\mathrm{low}}}\longrightarrow0,
\qquad
F(n,d_n,q_n)\longrightarrow0.
\]

To extract the two ratio limits, we establish the explicit positive-retention comparison. Fix admissible \(n,d\) and \(0<q\le1\). The reveal-probability bounds from the first part imply
\[
c_{\mathrm{exp}}v_{\mathrm{F}}(n,d,q)
\le
\begin{cases}
\dfrac{d}{nq},&dq\le1,\\[4pt]
\dfrac{d^2}{n},&dq>1.
\end{cases}
\]
Indeed, in the first case \(p\ge c_{\mathrm{exp}}dq\), and in the second
\(p\ge c_{\mathrm{exp}}\). Along with the two lower comparisons for
\(v_{\mathrm{F}}\), this gives
\[
\frac{d^2}{n}+\frac{d}{nq}\le2v_{\mathrm{F}}(n,d,q),
\qquad
c_{\mathrm{exp}}v_{\mathrm{F}}(n,d,q)
\le\frac{d^2}{n}+\frac{d}{nq}.
\]
If \(v_{\mathrm{F}}(n,d,q)\ge1\), then
\[
\frac12\min\left\{1,\frac{d^2}{n}+\frac{d}{nq}\right\}
\le1=F(n,d,q).
\]
If \(v_{\mathrm{F}}(n,d,q)<1\), then
\[
\frac12\min\left\{1,\frac{d^2}{n}+\frac{d}{nq}\right\}
\le\frac12\left(\frac{d^2}{n}+\frac{d}{nq}\right)
\le v_{\mathrm{F}}(n,d,q)=F(n,d,q).
\]
For the other direction, \(0<c_{\mathrm{exp}}<1\) and
\(F(n,d,q)=\min\{1,v_{\mathrm{F}}(n,d,q)\}\) imply
\[
c_{\mathrm{exp}}F(n,d,q)\le1,
\qquad
c_{\mathrm{exp}}F(n,d,q)
\le\frac{d^2}{n}+\frac{d}{nq}.
\]
Therefore
\[
\frac12\min\left\{1,\frac{d^2}{n}+\frac{d}{nq}\right\}
\le F(n,d,q)
\le
c_{\mathrm{exp}}^{-1}
\min\left\{1,\frac{d^2}{n}+\frac{d}{nq}\right\}.
\]
% lean: CausalSmith.Experimentation.ThinnedgraphAdditiveRiskFrontier.frontier_elbow

For the sequence argument, define on the tail \(n\ge4\)
\[
\rho_{\mathrm{deg}}(n):=\frac{d_n^2}{n}\in[0,\infty),
\qquad
\rho_{\mathrm{audit}}(n):=
\begin{cases}
+\infty,&q_n=0,\\[2pt]
\dfrac{d_n}{nq_n},&q_n\ne0,
\end{cases}
\quad\in[0,+\infty].
\]
All subsequent sequence bounds concern this tail. Since \(F(n,d_n,q_n)\to0\), eventually
\[
F(n,d_n,q_n)<\frac12.
\]
The identity \(F(n,d_n,0)=1\) forces \(q_n>0\) eventually. The lower comparison then gives
\[
\frac12\min\{1,\rho_{\mathrm{deg}}(n)+\rho_{\mathrm{audit}}(n)\}
\le F(n,d_n,q_n)<\frac12.
\]
Thus the sum is eventually less than one, and
\[
0\le\rho_{\mathrm{deg}}(n)
\le\rho_{\mathrm{deg}}(n)+\rho_{\mathrm{audit}}(n)
\le2F(n,d_n,q_n),
\]
\[
0\le\rho_{\mathrm{audit}}(n)
\le\rho_{\mathrm{deg}}(n)+\rho_{\mathrm{audit}}(n)
\le2F(n,d_n,q_n).
\]
Both ratios converge to zero by squeezing. The audit ratio is eventually finite, so its real convergence also gives convergence in the extended nonnegative reals.
% lean: CausalSmith.Experimentation.ThinnedgraphAdditiveRiskFrontier.precision_frontier_explicit

\item Conversely, suppose that
\[
\rho_{\mathrm{deg}}(n)\longrightarrow0,
\qquad
\rho_{\mathrm{audit}}(n)\longrightarrow0,
\]
with the second limit taken in the extended nonnegative reals. Eventually
\(\rho_{\mathrm{audit}}(n)<1\), which excludes \(q_n=0\); admissibility then gives
\(q_n>0\). Passing to real values at zero yields
\[
\frac{d_n}{nq_n}\longrightarrow0,
\qquad
\rho_{\mathrm{deg}}(n)+\rho_{\mathrm{audit}}(n)\longrightarrow0.
\]
The positive-retention comparison gives, eventually,
\[
0\le F(n,d_n,q_n)
\le c_{\mathrm{exp}}^{-1}
\bigl(\rho_{\mathrm{deg}}(n)+\rho_{\mathrm{audit}}(n)\bigr),
\]
so \(F(n,d_n,q_n)\to0\).
% lean: CausalSmith.Experimentation.ThinnedgraphAdditiveRiskFrontier.frontierScale_tendsto_iff

Choose the measurable seeded attaining estimator for every \(n\):
\[
T_n(O,\xi):=\widehat T^\star_{n,d_n,q_n}(O,\xi).
\]
Its risk satisfies, for every \(n\ge4\),
\[
0\le
\sup_{\theta\in\mathcal M_{n,d_n}}\mathcal L(T_n;\theta,q_n)
\le24F(n,d_n,q_n)\longrightarrow0.
\]
The finitely many indices below four have no effect on the limit. This proves sufficiency and completes the equivalence.
% lean: CausalSmith.Experimentation.ThinnedgraphAdditiveRiskFrontier.precision_frontier_explicit

\item The equality of the total observable rules follows directly from their construction:
\[
T^\star_{n,d,q}(O)=T^{\mathrm{up}}(O)
\qquad
\left(
\begin{array}{l}
V\text{ any finite population},\quad n,d\in\mathbb N,\\
q\in\mathbb R,\quad O\text{ any record on }V
\end{array}
\right).
\]
On the admissible parameter range these are the rules of
\cref{obj:def:frontier-rule,obj:def:upper-rule}, and the measurable seeded representative used above is precisely the representative of the upper rule.
% lean: CausalSmith.Experimentation.ThinnedgraphAdditiveRiskFrontier.frontier_answer

\item Finally, \cref{obj:def:frontier-scale} gives the endpoint identities. At zero retention it gives \(F(n,d,0)=1\); at full retention, \(d\ge1\) gives \(1-(1-1)^d=1\). Hence
\[
F(n,d,0)=1,
\qquad
F(n,d,1)=\min\left\{1,\frac{d^2}{n}\right\}.
\]
The positive-retention comparison established above, with
\(c_{\mathrm{exp}}=1-e^{-1}\), is
\[
\frac12\min\left\{1,\frac{d^2}{n}+\frac{d}{nq}\right\}
\le F(n,d,q)
\le
(1-e^{-1})^{-1}
\min\left\{1,\frac{d^2}{n}+\frac{d}{nq}\right\},
\qquad 0<q\le1.
\]
These are the remaining assertions.
% lean: CausalSmith.Experimentation.ThinnedgraphAdditiveRiskFrontier.frontier_elbow
\end{enumerate}
\end{proof}

\bibliographystyle{plainnat}

\bibliography{references}

\end{document}
